\documentclass[paper=a4, fontsize=10pt, jafontsize=10pt]{jlreq}

% LuaLaTeX + jlreq. TeX Live 2024 ships both jlreq and luatexja, and the
% HaranoAji fonts they default to, so this builds with no font installation.
\usepackage{luatexja}
\usepackage[haranoaji]{luatexja-preset}

\usepackage{amsmath}
\usepackage{amssymb}
\usepackage{graphicx}
\usepackage{booktabs}
\usepackage{longtable}
\usepackage{siunitx}
\usepackage{xcolor}
\usepackage[table]{xcolor}
\usepackage{caption}
\usepackage{subcaption}
\usepackage{float}
\usepackage[backend=bibtex,style=authoryear,natbib=true]{biblatex}
\usepackage[hidelinks]{hyperref}

\addbibresource{refs.bib}

% Figures live where the pipeline wrote them, so the report can never show a
% stale copy.
\graphicspath{{../out/figures/}{../out/figures/hokkaido/}{../out/figures/wakayama/}
  {../out/figures/hokkaido_week/}{../out/figures/wakayama_week/}
  {../out/figures/compare/}}

\sisetup{
  detect-all,
  per-mode = symbol,
  range-phrase = {--},
  range-units = single,
}

% Japanese unit names are written as \num{140} 分, not \SI{140}{分}. siunitx
% sets the unit argument in a Latin font, so a CJK unit is silently dropped and
% the reader is left with a bare number: the log says "Missing character: There
% is no 分 (U+5206) in font [lmroman10-regular]" and the PDF says "約 140".
% \DeclareSIUnit does not help -- the glyph is lost at typesetting, not at
% parsing -- and neither does wrapping it in \mbox or \text.

\captionsetup{
  font = small,
  labelfont = bf,
  justification = raggedright,
  singlelinecheck = false,
}

\definecolor{linkblue}{HTML}{2A78D6}
\hypersetup{
  colorlinks = true,
  linkcolor = linkblue,
  citecolor = linkblue,
  urlcolor = linkblue,
}

% Shorthands for the units that appear constantly.
\newcommand{\ms}{\si{\metre\per\second}}
\newcommand{\hpa}{\si{\hecto\pascal}}
\newcommand{\km}{\si{\kilo\metre}}

% A box for the "why this choice" asides, which are the point of this report.
\usepackage{tcolorbox}
\tcbuselibrary{breakable, skins}
\newtcolorbox{whybox}[1]{
  colback = linkblue!4,
  colframe = linkblue!40,
  boxrule = 0.4pt,
  left = 6pt, right = 6pt, top = 5pt, bottom = 5pt,
  title = {#1},
  fonttitle = \bfseries\small,
  breakable,
}

% Code blocks. listings rather than minted: minted shells out to pygmentize,
% which lives only in the conda environment, so a plain `make` outside it
% would fail. listings is pure TeX and always builds.
%
% Only the command names are keywords. Using listings' bash language would
% colour `cd` and `env` while leaving `conda` and `python` plain, which reads
% as random rather than as a rule.
\usepackage{listings}
\usepackage{needspace}
\definecolor{codebg}{HTML}{F7FAFD}
\definecolor{codecomment}{HTML}{6B7280}
\definecolor{codestring}{HTML}{9A4B00}
\lstset{
  basicstyle = \ttfamily\small,
  keywordstyle = \color{linkblue}\bfseries,
  % No \itshape: the Japanese comments would come out as slanted mincho.
  commentstyle = \color{codecomment},
  stringstyle = \color{codestring},
  morekeywords = {conda, python, pytest, make, cd, ls, cat},
  comment = [l]{\#},
  morestring = [b]",
  morestring = [b]',
  showstringspaces = false,
  extendedchars = true,
  columns = fullflexible,
  keepspaces = true,
  breaklines = true,
  breakindent = 2em,
  backgroundcolor = \color{codebg},
  frame = single,
  rulecolor = \color{linkblue!40},
  framerule = 0.4pt,
  framesep = 6pt,
  xleftmargin = 6pt,
  xrightmargin = 6pt,
  aboveskip = \medskipamount,
  belowskip = \medskipamount,
}
% A framed block split two lines before a page break looks like a mistake.
% Start one on a fresh page unless several lines fit.
\newcommand{\codeblock}{\Needspace*{6\baselineskip}}

% jlreq measures the width of the contents-page number box once, when the
% class loads and \thesection is still \arabic{section}. The appendix
% renumbers sections to "付録A", which no longer fits, so the number runs
% into the title ("付録 BCSV 表の一覧"). Widen the box from the appendix on.
\makeatletter
\newcommand{\appendixtocwidth}{%
  \renewcommand{\l@section}{\@dottedtocline{1}{0pt}{4.5em}}%
}
\makeatother


\title{JRA-3Q を用いた高層風の統計的傾向の解析}
\author{西濱大将}
\date{2026年9月}

\begin{document}
\maketitle

\begin{abstract}
スペースバルーンの落下位置は、打ち上げてから着地するまでに通過する全高度の風で決まる。
ところが現行のシミュレータが使える風データはどちらも「予報」であり、
数日先までしか見られないうえ、高度 \SI{24}{\kilo\metre} より上の風を持っていない。
スペースバルーンの飛行はその上まで及ぶ。

本報告は、気象庁第3次長期再解析 JRA-3Q の 2016--2025 年 10 年分から、
北海道・和歌山・愛媛（愛南町）の 3 領域について「いつ・どこで・どの高度の風が、ふつうどちらへどれだけ吹くか」を
統計としてまとめたものである。
高度は地上から約 \SI{40}{\kilo\metre} まで、時期は月を 2 分割した半月単位で扱う。

まとめた統計で個々の日の落下位置をどれだけ正確に予測できるかも測った。
落下位置の誤差は、時期の情報をまったく使わない場合と比べて、
和歌山では \SI{45}{\percent}、愛媛では \SI{47}{\percent} 下がるのに対し、北海道では \SI{17}{\percent} しか下がらない。
同じ手順で作った気候値でも、\textbf{南の 2 領域では時期から落下位置をよく予測できるが、北海道ではあまり予測できない}。
和歌山と愛媛は経度が \ang{2} 離れているが、上空の気候はほとんど同じだった。

さらに、現行シミュレータが風を\textbf{打ち上げ時刻と打ち上げ地点で固定している}近似の
誤差を、再解析の風の中で軌道を計算して測った。
この近似による落下点のずれは \SIrange{2.8}{5.3}{\kilo\metre}、飛行距離の
\SIrange{2}{5}{\percent} にとどまる。
一方、打ち上げが 1 日ずれると落下点は \SIrange{29}{63}{\kilo\metre} 動く。
落下点のずれは、風を固定したことより、日程がずれたことによるもののほうがはるかに大きい。
また、Web アプリのモンテカルロは機体の諸元だけを振っているため、
示している散らばりは実際に起こりうる大きさの 1 割程度である。
残りは風のばらつきによるもので、それは現在まったく考慮されていない。

読者に気象学の予備知識は仮定しない。
各章では「何を計算したか」だけでなく「なぜその計算を選んだか」を述べる。
特に、風向のような円環量を平均してはいけない理由、
平均風速に 2 種類ある理由、
そして 600 個の観測が 600 個分の情報を持たない理由は、
結果の読み方そのものに関わるため丁寧に説明する。
\end{abstract}

\tableofcontents
\clearpage

\section{はじめに}

\subsection{何を知りたいのか}

スペースバルーンは、ヘリウムまたは水素を詰めたゴム気球で観測機器を吊り上げ、
上空 \SI{30}{\kilo\metre} 前後で気球が破裂したあとパラシュートで降りてくる。
打ち上げから回収までおよそ 2--3 時間、水平方向には数十から数百キロメートル流される。

このとき機体を運ぶのは風だけである。
上昇中も下降中も、機体は自分の力で進むわけではなく、
その高度の風にそのまま乗って流されていく。
したがって落下地点は、地上から破裂高度までの\textbf{全高度の風}の積み重ねで決まる。

打ち上げの計画を立てる段階で知りたいのは、次のような問いである。

\begin{itemize}
  \item 3 月の前半に十勝から上げたら、機体はふつうどちらへ何キロ流れるのか。
  \item 落下範囲のばらつきは、どの時期が小さいのか。
  \item 海に落ちる危険が小さいのはいつか。
  \item モンテカルロ計算で風に与えるべき「ばらつき」は、実際どれくらいなのか。
\end{itemize}

いずれも「明日の風」ではなく「その時期の風のふつうの姿」を問うている。
これを気候値（climatology）と呼ぶ。

\subsection{なぜ既存のデータでは答えられないのか}
\label{sec:intro-ceiling}

現在のシミュレータには風の入手経路が 2 つあるが、どちらも気候値を与えない。

\paragraph{Open-Meteo の気圧面 API（\texttt{backend/bfps/core/wind.py}）}
気圧面ごとの風を返してくれるが、扱える一番上の面が \SI{30}{\hecto\pascal}、
高度にして約 \SI{23.8}{\kilo\metre} までしかない。
実装はそれより上を「最上層と同じ風が続く」と仮定して引き伸ばしている。
破裂高度を \SI{30}{\kilo\metre} とすると、この空白は飛行時間の 17 \% にあたる。

この空白の影響の大きさは季節によってまったく違う。
本報告の結果を使って測ったところ、冬は水平移動の 1--5 \% にすぎないのに対し、
夏は北海道で 54 \%、和歌山で \textbf{111 \%}、愛媛では \textbf{140 \%} に達する（第 \ref{sec:ceiling} 節）。
しかも夏は、落下距離が年間で最も短くなるので、打ち上げに向いた季節である。
つまり\textbf{現在の実装は、打ち上げに最も向いた季節に最も大きな誤差を出す}。

\paragraph{Tawhiri / GFS（\texttt{frontend/src/api/tawhiri.ts}）}
こちらは全高度をカバーするが、数値予報の出力であり、
実装上も「12 時間前から 7 日後まで」に限られている。
来月の前半を問い合わせることはできない。

つまり、\textbf{過去の風を時期ごとにまとめたデータが手元に無い}。
本報告はそのデータを作る。

\subsection{本報告がしたこと}

気象庁の長期再解析 JRA-3Q から、

\begin{itemize}
  \item 期間 2016--2025 年（10 年）、
  \item 高度 地上から約 \SI{40}{\kilo\metre}（気圧 \SIrange{1000}{3}{\hecto\pascal} の 38 面）、
  \item 水平 \ang{0.375} 格子（約 \SI{40}{\kilo\metre} 間隔）、
  \item 時間 6 時間ごと（1 日 4 回）
\end{itemize}

のデータを取得し、北海道・和歌山・愛媛（愛南町）それぞれについて、
月を 2 つに割った半月ごと・高度ごと・格子点ごとに
風向・風速の統計をとった。

出力は NetCDF（本体）、Parquet（扱いやすい表形式）、要点だけの CSV、
そして本報告に載せた図である。

\subsection{この報告の読み方}

第 \ref{sec:data} 章でデータの素性を、第 \ref{sec:method} 章で計算方法を述べる。
方法の章がやや長いのは、手順に加えて
\textbf{なぜその方法でなければいけないのか}まで説明しているためである。
風向のような「円になっている量」は、ふつうの平均を取ると誤った答えが出る。
その理屈を飛ばすと、第 \ref{sec:hokkaido} 章以降の図の読み方を間違える。

第 \ref{sec:validation} 章は\textbf{この統計をどこまで信じてよいか}を扱う。
「月を 2 分割した半月」という単位は人間側の都合で決めたもので、
データがその粒度を支えているかは別の話である。
そこを 1 年抜き交差検証で確かめた。結果の図を見る前に読んでほしい。

結果は第 \ref{sec:hokkaido} 章（北海道）、第 \ref{sec:wakayama} 章（和歌山）、
第 \ref{sec:ehime} 章（愛媛）、第 \ref{sec:compare} 章（3 領域の比較）にある。
第 \ref{sec:limits} 章には、この結果で言えないことを書いた。

\section{データ}
\label{sec:data}

\subsection{再解析とは何か}

天気の観測は、あるところでは密で、あるところではまばらである。
陸上には気象台もラジオゾンデもあるが、
太平洋の真ん中の高度 \SI{25}{\kilo\metre} を毎日測っている人はいない。
だから「過去のある日の、ある地点の、ある高度の風」を観測だけから並べることはできない。

\textbf{再解析}（reanalysis）はこの穴を埋める方法である。
数値予報モデル---大気の運動方程式を格子上で解くプログラム---を走らせながら、
その時刻に得られた観測（地上観測、ラジオゾンデ、航空機、船舶、衛星）を
次々に取り込んでモデルの状態を修正していく。
これを過去の全期間にわたって、\textbf{同一のモデルと同一の手法で}やり直したものが再解析である。

\begin{whybox}{なぜ「同一の手法で」が重要なのか}
現業の天気予報モデルは毎年改良される。
そのため、現業の解析結果をそのまま何十年も並べると、
気候が変わったのかモデルが変わったのか区別できない。
再解析は全期間を一つのモデルでやり直すので、
少なくともモデル由来の不連続は持ち込まれない。

ただし観測網そのものは時代とともに変わっている（衛星観測は 1979 年以降に本格化した）。
これは再解析でも消せない。本報告が 2016 年以降に限った理由の一つである。
\end{whybox}

したがって再解析の値は\textbf{観測そのものではない}。
観測に強く拘束された、モデルによる最良推定である。
観測が密なところ（日本の陸上、対流圏）では観測に近く、
まばらなところ（高層、外洋）ではモデルの寄与が大きい。

\subsection{なぜ JRA-3Q か}

再解析にはいくつか選択肢がある。図\ref{fig:f12} に並べた。

\begin{figure}[H]
\centering
\includegraphics[width=1.0\linewidth]{compare/F12_reanalyses.pdf}
\caption[再解析プロダクトの比較]{主な再解析の比較。
左は横軸が水平格子間隔（\textbf{細かいほうが右}）、縦軸が気圧面プロダクトの上端高度。
本用途に適するのは右上で、下の帯は破裂高度 \SI{30}{\kilo\metre} に届かない範囲である。
右は対象期間と、気圧面の数・解析間隔。
\textbf{JRA-3Q の値だけが実測}で、本パイプラインが取得したファイルから読んでいる。
他は各提供元が公表している\textbf{気圧面プロダクト}の仕様を手で入力したもので、
モデル格子やモデル面の仕様ではない。
仕様は更新されうるので、自分の研究に選ぶときは最新の値を確認してほしい。}
\label{fig:f12}
\end{figure}

図から読めることが 3 つある。

\textbf{NCEP/NCAR R1 は使えない。} 気圧面プロダクトが \SI{10}{\hecto\pascal}
（約 \SI{31}{\kilo\metre}）で止まる。ちょうど破裂高度なので、
飛行のいちばん上の部分に風のデータが無いことになる。
格子も \SI{200}{\kilo\metre} 級で、対象領域に数点しか入らない。

\textbf{JRA-55 は解像度が足りない。} \ang{1.25} は和歌山の領域に 4 点しか置けない。
「場所による違い」を論じる本報告の目的に合わない。
2023 年で更新も止まっている。

\textbf{残るのは ERA5、MERRA-2、JRA-3Q の 3 つ}で、
どれも高さと解像度の条件は満たす。ここから先は優劣ではなく選択である。

JRA-3Q を選んだ理由を挙げる。

\begin{enumerate}
  \item \textbf{気象庁自身の製品である。}
        日本周辺の観測、特に国内のラジオゾンデや日本の静止衛星のデータを厚く同化している。
        対象領域が日本である以上、この点で他の再解析より精度が期待できる。
  \item \textbf{高さが足りる。}
        気圧面プロダクトが 45 面あり、上端は \SI{0.01}{\hecto\pascal}（約 \SI{80}{\kilo\metre}）に達する。
        スペースバルーンに必要な \SI{40}{\kilo\metre} までは十分に余裕がある。
        Open-Meteo の \SI{30}{\hecto\pascal} 上限という制約がここで解ける。
  \item \textbf{期間が長い。}
        1947 年 9 月から現在まで。今回は使わないが、将来もっと長い統計が要るときに困らない。
  \item \textbf{入手できる。}
        NSF NCAR の GDEX から無償で、しかも匿名で取得できる。
\end{enumerate}

\begin{whybox}{ERA5 を選ばなかったのは、性能の問題ではない}
図\ref{fig:f12} のとおり ERA5 は JRA-3Q より水平解像度が細かく（\SI{20}{\kilo\metre} 対
\SI{31}{\kilo\metre}）、時間解像度も 1 時間ごとで 6 倍細かい。
気圧面の上端は \SI{1}{\hecto\pascal}（約 \SI{48}{\kilo\metre}）で、
これもスペースバルーンには十分である。
本用途の要件だけを見れば ERA5 のほうが有利な点が多い。

選ばなかった理由は 2 つある。

第一に、\textbf{ERA5 は同じプロジェクトの別のメンバーが並行して解析している}。
2 人で同じ製品を重複して扱うより、別々の再解析を分担したほうが、
あとで突き合わせたときに「製品による違い」が見える。
本報告の数字が JRA-3Q 固有のものなのか再解析一般に言えることなのかは、
本報告だけでは判定できない。それは 2 つを比べて初めて分かる。

第二に、日本周辺という対象に対しては\ref{sec:data} 節の理由 1（気象庁自身の製品で、
国内の観測を厚く同化している）が有利に働く可能性がある。
これは\textbf{検証していない期待であって、測った事実ではない}。
どちらがこの領域で実際に正確かは、2 つの解析が出そろってから比べるべきことである。
\end{whybox}

\subsection{使ったプロダクトと設定}

\begin{table}[H]
\centering
\caption{取得したデータの仕様}
\begin{tabular}{ll}
\toprule
項目 & 値 \\
\midrule
データセット & JRA-3Q \texttt{anl\_p}（等圧面解析値）\\
配布元 & NSF NCAR GDEX, データセット \texttt{d640000} \\
水平解像度 & \ang{0.375}（$960 \times 480$ のガウス格子、約 \SI{40}{\kilo\metre}）\\
鉛直 & 気圧面 38 面（\SIrange{1000}{3}{\hecto\pascal}）\\
時間 & 6 時間ごと（00, 06, 12, 18 UTC）\\
期間 & 2016--2025 年（10 年、14{,}610 時刻）\\
変数 & 東西風 $u$、南北風 $v$、ジオポテンシャル高度 $Z$ \\
領域 & \ang{30}--\ang{48}N / \ang{132}--\ang{150}E から各地域を切り出し \\
\bottomrule
\end{tabular}
\end{table}

\paragraph{なぜ 2016--2025 年なのか}
10 年という長さは 2 つの要求の折り合いである。
短すぎると 1 つのビンに入る標本が減り、統計が暴れる。
長すぎると、前述の観測網の変化や、気候そのものの長期変化が混ざってくる。
また現実的な問題として、10 年分で取得は約 2.5 時間、キャッシュは約 \SI{5}{\giga\byte} である。

\paragraph{なぜ \ang{0.375} なのか}
JRA-3Q には \ang{1.25} の軽いプロダクトもあり、
北海道だけなら \ang{1.25} でも格子点が 54 個取れて図になる。
しかし和歌山の領域（\ang{2.5} $\times$ \ang{3}）には \ang{1.25} 格子だと点が 4 個しか入らず、
「場所による違い」を論じることができない。
\ang{0.375} なら和歌山でも約 60 点、北海道で約 540 点になる。

なお、解像度を上げても\textbf{分かることが増えるとは限らない}点には注意が要る。
JRA-3Q のモデル格子は約 \SI{40}{\kilo\metre} であり、
\ang{0.375} のプロダクトはその内挿である。
山脈 1 本のスケールの現象は、どちらの解像度でも解けていない。

\paragraph{なぜ気圧面なのか}
大気のモデルは地形に沿った座標で計算されるが、配布は気圧面か等高度面が普通である。
気圧面を選んだのは、スペースバルーンの高度計が実質的に気圧計であること、
そして既存のシミュレータが Open-Meteo の気圧面を前提に作られていて
換算の仕方を揃えられることによる。
気圧面から実際の高度への換算は第 \ref{sec:method} 章で述べる。

\paragraph{なぜ \SI{3}{\hecto\pascal} で切ったのか}
\SI{3}{\hecto\pascal} は約 \SI{40}{\kilo\metre} である。
ゴム気球の破裂高度はふつう \SIrange{25}{33}{\kilo\metre} なので、
これで飛行の全区間に余裕を持って収まる。
それより上の 7 面（\SI{2}{\hecto\pascal} 以上）は中間圏に入り、
スペースバルーンには関係がないうえ、図の縦軸を無駄に伸ばして対流圏を見づらくする。

\subsection{ライセンスと引用}

JRA-3Q は \textbf{CC BY-NC-SA 4.0} で提供されている \citep{gdex_d640000}。
これは次の 3 つを意味する。

\begin{itemize}
  \item \textbf{非商用に限る。} 商用目的では使えない。
  \item \textbf{継承。} 本報告や出力データは二次的著作物であり、
        再配布するなら同じ CC BY-NC-SA 4.0 を付けなければならない。
        リポジトリ本体のライセンスとは別である。
  \item \textbf{表示。} 出典を明記しなければならない。
\end{itemize}

学術・技術文書での引用は JRA-3Q の総合報告 \citep{kosaka2024jra3q} に対して行う。
この継承条項があるため、本パイプラインの出力は
Web アプリのビルドには混ぜず、独立したディレクトリに置いている。

\section{方法}
\label{sec:method}

この章は「何をしたか」と同時に「なぜそうしたか」を述べる。
特に \ref{sec:circular} 節から \ref{sec:neff} 節までは、
結果の図の読み方に直結するので飛ばさないでほしい。

\subsection{時間の切り方}
\label{sec:weeks}

1 か月を次の 2 つに分ける。

\begin{center}
\begin{tabular}{lll}
\toprule
ビン & 日 & 日数 \\
\midrule
H1 & 1--15 日 & 15 \\
H2 & 16 日--月末 & 13--16 \\
\bottomrule
\end{tabular}
\end{center}

1 年で 24 ビンになる。

\begin{whybox}{なぜ ISO 週や半旬ではなく「月内固定」なのか}
境界の決め方には候補が 3 つあった。

\textbf{ISO 週}（月曜始まりの通し番号）は年ごとに月との対応がずれる。
「1 月の 2 番目のビン」と言ったとき、年によって指す日付が変わってしまう。

\textbf{半旬}（5 日 $\times$ 73）は気候学では標準的で、標本数が完全に揃う。
しかし月境界をまたぐので、「何月の話か」を言うのに一手間かかる。

\textbf{月内固定}は、「3 月前半」が毎年必ず 3 月 1--15 日を指す。
打ち上げ計画は「来年の 3 月前半に上げたい」という形で立つものなので、
人が計画を立てる単位にそのまま対応する。これを選んだ。
\end{whybox}

\begin{whybox}{なぜ 7 日ではなく 15 日なのか}
本報告は当初、1 か月を 7 日 $\times$ 4 に分けていた（W1 = 1--7 日、
W2 = 8--14 日、W3 = 15--21 日、W4 = 22 日--月末）。
7 日という幅は、データからどこまで細かく区別できるかを調べる前に決めたものである。

第 \ref{sec:validation} 章がそれを実測した。1 年抜き交差検証の誤差は
\numrange{5}{15} 日のあいだで平坦で、1 標準誤差ルールが選ぶのは和歌山と愛媛で
\textbf{15 日}、北海道で 30 日である。
7 日は間違いではない。「区別できない範囲」の中に入っている。
ただし\textbf{データから読み取れる細かさよりも細かかった}。

\textbf{精度は変わらない。} 季節ごとに交差検証をやり直すと、
15 日ビンと 7 日ビンの差は北海道・和歌山の 2 領域 $\times$ 4 季節の 8 通りすべてで
1 標準誤差未満に収まり、うち 6 通りでは 15 日のほうがわずかに良い。

\textbf{変わるのは分位点の精度である。} 6 時間ごとの解析は互いによく似ていて
（ラグ 1 自己相関がおよそ 0.975）、7 日ビンの 280 標本は実質 4 個ぶんの
情報しか持たない（\ref{sec:neff} 節）。
ビンを倍にすると、本報告が主に使う量のノイズがほぼ半分になる。
\SI{200}{\hecto\pascal} の風速 p90 の標準誤差は、
7 日で \SI{1.70}{\metre\per\second}、15 日で \SI{1.04}{\metre\per\second} である。

7 日の製品は \ref{sec:appendix-week} に残してある。
\end{whybox}

代償は H2 だけが 13--16 日と変動することである。
H1 は常に 15 日なので標本数は 10 年で 600 個、
H2 は月によって 520--640 個になる。
したがって\textbf{ビンごとに標本数が違う}。
これを忘れると、たとえば「H2 だけ最大風速が大きい」という見かけの結果が出る
（標本が多いぶん極値を引きやすいだけ）。
このため標本数 $n$ はすべての統計量と並べて保存し、
各領域の章の最初に示す（図\ref{fig:f0-hokkaido}、\ref{fig:f0-wakayama}）。

なお 7 日ビンの日数は 7--10 日で 43 \% 変動していたので、
この点は以前より改善している。

\subsection{気圧面を高度に直す}
\label{sec:altitude}

データは気圧面上にあるが、人が読みたいのは高度である。
換算には 2 段階あり、混同しやすい。

\paragraph{第 1 段階：ジオポテンシャル高度}
JRA-3Q が持っている $Z$ は\textbf{ジオポテンシャル高度}で、単位は gpm（ジオポテンシャルメートル）である。
これは距離ではなく、重力に逆らって単位質量を持ち上げるのに要した仕事を
標準重力 $g_0 = \SI{9.80665}{\metre\per\second\squared}$ で割った量である。

\begin{equation}
  Z = \frac{1}{g_0}\int_0^z g(\phi, z')\, \mathrm{d}z'
\end{equation}

\paragraph{第 2 段階：幾何高度}
実際の重力は高度とともに弱まり、緯度によっても違う。
そのためジオポテンシャル高度と実際の距離（幾何高度 $z$）はずれる。

\begin{equation}
  z = \frac{R_\mathrm{eff}(\phi)\, Z\, g_0}{g(\phi)\, R_\mathrm{eff}(\phi) - Z\, g_0}
\end{equation}

ここで $g(\phi)$ は Somigliana の式による海面重力、
$R_\mathrm{eff}(\phi)$ は List の有効地球半径 \citep{list1968smithsonian} である。

補正量は \ang{45}N・\SI{30}{\kilo\metre} で約 \SI{+140}{\metre}。
一見わずかだが、上層の気圧面どうしの間隔と同じ桁であり、無視すると層が 1 枚ずれかねない。

\begin{whybox}{なぜ変数名を \texttt{height} にしないのか}
コード中では必ず \texttt{\_gpm}（ジオポテンシャル）か \texttt{\_geom\_m}（幾何）を付けている。
単に \texttt{height} と書くと、どちらの意味で使われているか 3 か月後には分からない。
この取り違えは古典的な事故であり、名前で防ぐのが一番安い。
\end{whybox}

なお既存のシミュレータ（\texttt{backend/bfps/core/atmosphere.py}）は、
気圧から高度への換算に国際標準大気 ISA \citep{icao1993isa} を使っている。
ISA は季節も緯度も天気も知らない理想化なので、実際とどれだけ違うかは確かめる価値がある。
そのずれは図\ref{fig:f10-hokkaido} と図\ref{fig:f10-wakayama} に示す。

\subsection{なぜ風向を平均してはいけないのか}
\label{sec:circular}

風向は角度であり、\ang{0} と \ang{360} は同じ向きを指す。
つまり数直線ではなく円の上にある。
円の上の量に、ふつうの算術平均を当てると誤った答えが出る。

\begin{whybox}{北風 2 回の平均が南風になる}
ほぼ真北から吹く風を 2 回観測したとする。1 回目は \ang{359}、2 回目は \ang{1}。
算術平均をとると
\[
  \frac{359 + 1}{2} = \ang{180}
\]
となり、\textbf{真南}という答えが返る。実際には 2 回とも北風だったのに。
\end{whybox}

正しいやり方は、角度を平均しないことである。
各時刻の風をまず東西成分 $u$ と南北成分 $v$ に分解し、
\textbf{成分を平均して}から、最後に一度だけ向きに戻す。

\begin{align}
  \bar{u} &= \frac{1}{N}\sum_i u_i, \qquad \bar{v} = \frac{1}{N}\sum_i v_i \\
  \theta &= \left(\arctan2(\bar{u}, \bar{v}) \cdot \frac{180}{\pi} + 180\right) \bmod 360
\end{align}

本パイプラインでは、この規則をコメントではなく\textbf{構造}で守っている。
統計を集計するモジュールは $u$ と $v$ しか扱わず、
風向という量は出力を書く直前の変換関数にしか現れない。
集計器に角度が入らないので、角度を平均するコードが書けない。

\paragraph{風向の向きの取り方}
本報告の風向はすべて\textbf{気象学の慣習}、すなわち
「風が吹いてくる方位」を北 \ang{0} から時計回りに測る。
西風（西から東へ吹く風、$u > 0$）は \ang{270} である。
既存シミュレータの \texttt{met\_to\_uv} と同じ定義で、両者は数値的に一致する。

\subsection{平均風速が 2 種類ある理由}
\label{sec:twomeans}

平均風速には意味の違う 2 つがあり、両方を出している。

\begin{description}
  \item[ベクトル平均風速 $|\bar{\mathbf{V}}|$]
    まず $u, v$ を平均し、その長さを取る：$\sqrt{\bar{u}^2 + \bar{v}^2}$。
    これは\textbf{正味どれだけ運ばれるか}である。
  \item[スカラー平均風速 $\overline{|\mathbf{V}|}$]
    各時刻の風速を先に計算してから平均する：$\frac{1}{N}\sum_i \sqrt{u_i^2 + v_i^2}$。
    これは\textbf{ふつうどれくらい強く吹いているか}である。
\end{description}

両者は一般に一致しない。数学的に必ず
\begin{equation}
  |\bar{\mathbf{V}}| \le \overline{|\mathbf{V}|}
\end{equation}
が成り立つ（三角不等式）。

\begin{whybox}{ビンの半分が東風、半分が西風だったら}
どちらも \SI{20}{\metre\per\second} だったとする。

スカラー平均は \SI{20}{\metre\per\second}。強い風が吹いていたのは事実である。

ベクトル平均は \SI{0}{\metre\per\second}。ビン全体をならして見れば、機体はどこにも運ばれない。

スペースバルーンの平均的な移流を見たいならベクトル平均、
そのビンで遭遇する風の強さを見たいならスカラー平均。
どちらか一方だけでは足りない。
\end{whybox}

この 2 つの比を\textbf{定常度}（constancy）と呼ぶ。

\begin{equation}
  C = \frac{|\bar{\mathbf{V}}|}{\overline{|\mathbf{V}|}} \in [0, 1]
\end{equation}

$C = 1$ なら常に同じ向きに吹いており、平均風向はそのまま信用してよい。
$C$ が 0 に近いと風向が定まっておらず、\textbf{平均風向には意味がない}。

このため本報告の図では、$C < 0.3$ の領域の風向を灰色でマスクしている。
定まらない風向を自信ありげに描くことは、
このパイプラインが犯しうる最も有害な誤りだからである。

\subsection{分散楕円と、それが実態を表さない場合}
\label{sec:ellipse}

風のばらつきは、$u$ と $v$ の分散共分散行列
\begin{equation}
  \Sigma = \begin{pmatrix} \sigma_u^2 & \sigma_{uv} \\ \sigma_{uv} & \sigma_v^2 \end{pmatrix}
\end{equation}
に全部入っている。この固有値分解から $1\sigma$ の楕円（長軸・短軸・傾き）が出る。
これはモンテカルロ計算で風に与える摂動の大きさを、推測ではなく実測から決められることを意味する。

NetCDF には楕円ではなく $\sigma_u^2, \sigma_v^2, \sigma_{uv}$ の 3 つを保存している。
楕円は情報を落とさずにここから復元できるが、逆はできないからである。

ただし楕円は\textbf{$(u,v)$ の散らばりが 2 変量正規に近い}ことを仮定している%
\footnote{%
「2 変量正規」は、2 つの量が同時にとる値の散らばり方の型のひとつである。
ここでの 2 つの量は東西成分 $u$ と南北成分 $v$ で、
10 年ぶんの解析を $(u,v)$ の点として平面に打つと散布図ができる。
その散布図が\textbf{中心が 1 つだけの、楕円形の雲}になっている状態を指す。
中心から離れるほど点がなだらかに減り、
どの向きに切っても釣鐘型（正規分布）の断面になる。

この型であれば、雲の形は
平均 2 つ（$\bar u$, $\bar v$）と分散共分散行列 $\Sigma$ の 3 つ、
合わせて 5 個の数だけで完全に決まる。
だからこそ楕円 1 つで散らばりを表せる。逆に言えば、
\textbf{楕円で表せるのは、この型に近い散らばりだけ}である。

破れ方は本文のとおりで、
中心が 2 つある（＝雲が 2 つの塊に分かれる）場合が本報告では実際に起きる。
2 つの塊の中間に中心を 1 つ置いた楕円は、
どちらの塊にも入らない場所を「典型」として描いてしまう。
中心が 1 つでも、裾が正規分布より重い（極端な値が想定より多い）場合には
楕円が \SI{90}{\percent} の範囲を過小に見積もる。
本報告が分位点（p90, p95）を別途保存しているのは、後者への備えでもある。
分位点は分布の型を仮定しないので、裾が重くてもそのまま正しい。%
}。
これが破れる場面がはっきりある。
日本上空の成層圏では、風向が 4--5 月と 9--10 月に東西逆転する。
その時期のビンには「強い西風の日」と「強い東風の日」が混ざり、
散らばりは 2 つの塊に分かれる。
このとき平均はほぼゼロ、分散は巨大、楕円は実際には誰も経験しない中間の風を描く。

この二峰性は、破裂高度帯で、しかも毎年決まった時期に起きる。
そのため風配図を用意した（図\ref{fig:f7-hokkaido-h1} 以降）。
風配図は分布そのものを見せるので、
2 つの塊があればそう見える。平均と楕円には見せられない。

\subsection{なぜ分位点も出すのか}

平均だけでは運用判断ができない。
「平均 \SI{15}{\metre\per\second}」でも、
そのビンの 1 割の日が \SI{40}{\metre\per\second} なら、回収計画は変わる。

そこで風速の経験分位点 $p_{10}, p_{25}, p_{50}, p_{75}, p_{90}, p_{95}$ と最大値、
さらに \SI{10}{}、\SI{20}{}、\SI{30}{}、\SI{40}{\metre\per\second} を超える頻度を保存した。

分位点は\textbf{風速（スカラー）に対してのみ}計算している。
風向には大小の順序がないので、分位点という概念が定義できない
（どこを起点に数えるかで答えが変わってしまう）。

なお分位点は平均や分散から計算することができない。
全標本を持っていないと出せない量である。
そのため集計では 1 ビン分の全標本をメモリに載せて厳密に計算している
（北海道の 1 ビンで約 \SI{46}{\mega\byte}、十分に載る）。
近似アルゴリズムは使っていない。

\subsection{600 個の観測は 600 個分の情報ではない}
\label{sec:neff}

H1 のビンには 10 年 $\times$ 15 日 $\times$ 4 回 $=$ 600 個の解析値が入る。
しかしこれは\textbf{独立な観測 600 個}ではない。

6 時間しか離れていない 2 つの解析値は、ほとんど同じ天気を見ている。
低気圧や高気圧が日本上空を通過するには 2--5 日かかるので、
15 日のビンに入る「独立な天気」はせいぜい 3--7 個である。

これを無視して標準誤差を $\sigma/\sqrt{600}$ とすると、
確からしさを 3 倍ほど過大に見積もることになる。

そこで各年の連続した区間内でラグ 1 自己相関 $r_1$ を求め、年をまたいで平均し、
\textbf{有効標本数}
\begin{equation}
  n_\mathrm{eff} = n\,\frac{1 - r_1}{1 + r_1}
\end{equation}
を保存した \citep{wilks2011statistical}。
実際の値は中央値で 76--78、$n$ の 13 \% である（3 領域とも）。
本報告の誤差はすべて $n_\mathrm{eff}$ から計算している。

\begin{whybox}{なぜ年をまたいで自己相関を取らないのか}
10 年分を時間順に並べて一括で自己相関を取ると、
1 月末の標本と翌年 1 月頭の標本が「隣」として扱われてしまう。
両者は 6 時間ではなく 1 年離れており、当然ほぼ無相関である。
その偽の無相関が混ざると $r_1$ が小さく出て、有効標本数を過大評価する。
だから各年の連続区間の中だけで計算し、あとで平均している。
\end{whybox}

\subsection{時刻の扱い}

JRA-3Q の解析時刻は 00, 06, 12, 18 UTC で、日本時間ではそれぞれ
09, 15, 21, 03 JST にあたる（18 UTC は JST では翌日になる）。

日付の区切りは既定で\textbf{UTC} を使った。
月をまたぐ繰り越しが不要で、データ本来の軸だからである。
7 日のビンでは、JST で区切った場合と食い違うのは 28 標本のうち 1 つだけである。

とはいえ主張で済ませるのは望ましくないので、両方を計算して比べた。
北海道で UTC 区切りと JST 区切りを突き合わせると、次のようになる。

\begin{center}
\begin{tabular}{lrr}
\toprule
量 & 最大の差 & RMS \\
\midrule
平均風速 & \SI{0.73}{\metre\per\second} & \SI{0.11}{\metre\per\second} \\
東西風 $\bar u$ & \SI{0.72}{\metre\per\second} & \SI{0.13}{\metre\per\second} \\
定常度 & 0.061 & 0.006 \\
飛行距離 & \SI{1.6}{\kilo\metre} & --- \\
標本数 & 10 個（400 個中） & --- \\
\bottomrule
\end{tabular}
\end{center}

平均風速で最大 \SI{0.73}{\metre\per\second}、
飛行距離では平均 \SI{130}{\kilo\metre} に対して最大 \SI{1.6}{\kilo\metre} である。
本報告のどの結論も、この選択では変わらない。

また、下層には本物の日変化がある（海陸風、夜間の接地逆転）。
そこで \SI{300}{\hecto\pascal} 以下の層に限り、
解析時刻ごとに分けた副次的な統計も出している。
図\ref{fig:f8-hokkaido} が示すとおり、夜（03 JST）と昼（15 JST）の平均風速の差は
3 領域・4 季節のどれでも \SI{1}{\metre\per\second} 未満で、日々の変動（数十 \si{\metre\per\second}）に比べて小さい。
それが\textbf{上層では時刻を区別せず平均してよい根拠}になっている。
これも主張ではなく、図で示している。

\subsection{領域平均の取り方}

緯度経度格子をそのまま平均すると、高緯度側の格子を過大に数えてしまう
（同じ \ang{0.375} でも、北の格子は地上での東西幅が狭い）。
JRA-3Q のファイルにはガウス求積の重みが入っているので、それを使って重み付き平均をとった。

\subsection{風から落下位置へ}
\label{sec:drift}

風の統計そのものより、多くの読者が知りたいのは「結局どこに落ちるのか」である。
そこで、既存シミュレータと同じ飛行モデルを風の統計に通し、
落下位置の気候値も作った。

飛行モデルは \texttt{backend/bfps/core/trajectory.py} と同一である。

\begin{itemize}
  \item 上昇は一定速度（既定 \SI{5}{\metre\per\second}）で破裂高度（既定 \SI{30}{\kilo\metre}）まで
  \item 降下は海面等価の降下率（既定 \SI{6}{\metre\per\second}）を空気密度でスケールした
        $v(z) = v_0\sqrt{\rho_0/\rho(z)}$
  \item 水平方向は慣性なしの純移流。機体はその高度の風にそのまま乗る
\end{itemize}

このモデルのもとでは、水平変位は
\begin{equation}
  \Delta x = \sum_k u_k T_k, \qquad \Delta y = \sum_k v_k T_k
  \label{eq:drift}
\end{equation}
と書ける。$T_k$ は気圧面 $k$ の担当する層に、上りと下りで滞在する合計時間である。
$T_k$ は層の厚さと降下則だけで決まり、\textbf{風には依存しない}。

\begin{whybox}{線形であることから何が言えるか}
式 \eqref{eq:drift} は風のプロファイルについて\textbf{線形}である。
そこから 2 つのことが従う。

第一に、\textbf{各解析時刻ごとに着地点を計算してから統計をとれる}。
落下位置の散らばりは推定ではなく実測になる。
もし平均風のプロファイルだけから散らばりを出そうとすると、
「高度 10 km の風のゆらぎと高度 25 km の風のゆらぎがどれくらい連動しているか」
を仮定しなければならない。その仮定は本文のどこにも現れないまま結果を左右してしまう。
層間の相関は保存していないので、仮定を置かずに各時刻の着地点を数えるほうを選んだ。

第二に、\textbf{平均風から求めた変位と、変位の平均が一致する}。
気候値から計算した量では一般には成り立たない性質で、
ここでは成り立つと言い切れるので、そう書いておく。
\end{whybox}

降下が寄与する時間は上昇より短い。
高度 \SI{30}{\kilo\metre} の空気は薄く、密度スケーリングによって降下速度が海面の 10 倍近くになるためである。
既定のパラメータでの飛行時間は約 \num{140} 分（上昇 100 分、降下 40 分）になる。

\subsubsection{\texorpdfstring{$T_k$}{T\_k} をどう計算するか}
\label{sec:drift-layers}

式\eqref{eq:drift} は $T_k$ が分かれば解ける。その $T_k$ の作り方を書いておく。

\paragraph{1. 気圧面に高度を与える}
JRA-3Q はジオポテンシャル高度を別の変数として持っているので、それを取得してある。
ジオポテンシャル高度は重力加速度の緯度・高度変化を吸収した量なので、
幾何学的な高度に直してから使う（\ref{sec:altitude} 節）。
高度は日々数百メートル動くが、層の厚さを決めるのが目的なので、
ビンごとの平均を 1 本使う。

\paragraph{2. 層の境界を決める}
気圧面 $k$ が「担当する」層は、\textbf{隣の気圧面との中点}で区切る。
高度の昇順に並べた $z_k$ に対し、

\begin{equation}
  z_k^{\text{下}} = \frac{z_{k-1} + z_k}{2}, \qquad
  z_k^{\text{上}} = \frac{z_k + z_{k+1}}{2}
  \label{eq:layer-edges}
\end{equation}

とし、これを打ち上げ高度と破裂高度で切り詰める。
層の厚さは $\Delta z_k = \max(z_k^{\text{上}} - z_k^{\text{下}},\, 0)$ で、
破裂高度より完全に上にある気圧面は $\Delta z_k = 0$ になって寄与しない。

\paragraph{3. 上りと下りの滞在時間を足す}
上昇は一定速度、降下はその層の代表高度での密度スケール速度を使う。

\begin{equation}
  T_k = \underbrace{\frac{\Delta z_k}{w_{\uparrow}}}_{\text{上り}}
      + \underbrace{\frac{\Delta z_k}{v_0\sqrt{\rho_0/\rho(z_k)}}}_{\text{下り}}
  \label{eq:layer-seconds}
\end{equation}

実装は \texttt{stats/drift.py} の \texttt{layer\_seconds} である。

\paragraph{実際の数字}
和歌山の中心格子、1 月前半の平均高度プロファイルで計算すると次のようになる。

\begin{table}[H]
\centering
\caption{$T_k$ の計算例（和歌山中心、1 月前半、既定の飛行諸元）。
$v_{\downarrow}$ はその高度での降下速度。}
\label{tab:layer-seconds}
\sisetup{table-format=5.0}
\begin{tabular}{@{}S[table-format=4.0]SSS[table-format=3.1]S[table-format=3.1]S[table-format=3.1]S[table-format=3.1]@{}}
\toprule
{気圧面} & {高度} & {層の厚さ} & {$v_{\downarrow}$} & {上り} & {下り} & {$T_k$} \\
{[\si{\hecto\pascal}]} & {[\si{\metre}]} & {[\si{\metre}]} & {[\si{\metre\per\second}]}
& {[\si{\second}]} & {[\si{\second}]} & {[\si{\second}]} \\
\midrule
1000 &   184 &  103 &  6.1 &  20.7 & 17.1 &  37.7 \\
 850 &  1491 &  234 &  6.5 &  46.7 & 36.2 &  82.9 \\
 500 &  5570 &  740 &  8.0 & 147.9 & 92.6 & 240.5 \\
 200 & 11894 &  821 & 11.8 & 164.2 & 69.5 & 233.6 \\
 100 & 16271 & 1188 & 16.7 & 237.7 & 71.3 & 309.0 \\
  30 & 23787 & 2234 & 30.4 & 446.8 & 73.6 & 520.4 \\
  10 & 31018 & 1281 & 53.5 & 256.3 & 23.9 & 280.2 \\
   5 & 35779 &    0 & 77.8 &   0.0 &  0.0 &   0.0 \\
   3 & 39417 &    0 &102.5 &   0.0 &  0.0 &   0.0 \\
\bottomrule
\end{tabular}
\end{table}

38 層ぶんを足すと $\sum_k T_k = \SI{8223}{\second}$（\num{137.0} 分）になり、
上昇が \num{99.4} 分、降下が \num{37.7} 分である。

表から読み取れることがいくつかある。

\begin{itemize}
  \item \textbf{上の層ほど厚い。} \SI{30}{\hecto\pascal} の層は \SI{2234}{\metre} あり、
        \SI{1000}{\hecto\pascal} の \SI{103}{\metre} の 20 倍を超える。
        気圧は高度に対して指数的に減るので、等間隔の気圧面は上ほど高度間隔が広くなる。
  \item \textbf{下りは上の層でほとんど時間を使わない。} \SI{30}{\hecto\pascal} では
        降下速度が \SI{30.4}{\metre\per\second} まで上がるので、
        \SI{2234}{\metre} を \SI{74}{\second} で通過する。
        同じ層を上るには \SI{447}{\second} かかる。
        成層圏の強い風が落下点を大きく動かすのは、主に\textbf{上りの側}である。
  \item \textbf{破裂より上の 3 層は寄与しない。} \SI{7}{\hecto\pascal}（約 \SI{33}{\kilo\metre}）、
        \SI{5}{\hecto\pascal}、\SI{3}{\hecto\pascal} は $T_k = 0$ になる。
        取得はしてあるが、既定の破裂高度 \SI{30}{\kilo\metre} では使われない。
        破裂高度を上げれば自動的に効きはじめる。
\end{itemize}

\begin{whybox}{最下層は地面ではなく \SI{1000}{\hecto\pascal} から始まる}
表の最下行を見ると、層の下端が \SI{184}{\metre} になっている。
最も低い気圧面そのものの高度であって、地面ではない。
\SI{1000}{\hecto\pascal} より下の空気は、この計算に入っていない。

飛行時間 \SI{8223}{\second} のうち \SI{40}{\second} 弱にあたる。
しかもそこは飛行中でいちばん風が弱く、いちばんゆっくり通る区間である。
落下点への影響は無視できる大きさなので、このままにしてある。

打ち上げ地点の標高が高い場合は、逆にこの層をまたいで始まることになる。
本報告の対象は海岸に近い低地なので問題にならないが、
山地から上げる計画に使うときは \texttt{launch\_altitude\_m} を設定する必要がある。
\end{whybox}

$T_k$ が風に依存しないことは、この 3 段階のどこにも風が出てこないことから確かめられる。
高度・層の厚さ・降下則だけで決まる。
そのため $T_k$ はビンごとに 1 回計算すればよく、
各解析時刻ではその重みに風を掛けて足すだけで着地点が出る。
季節による違いも小さく、和歌山では 1 月 \SI{8223}{\second} に対し 7 月 \SI{8259}{\second} で、
\SI{0.3}{\percent} しか変わらない。

このモデルには 2 つの限界があり、どちらも\textbf{既存シミュレータが置いている仮定と同じ}である。

\begin{itemize}
  \item 飛行中ずっと風のプロファイルが空間的にも時間的にも凍結していると見なしている。
        \texttt{wind.py} が打ち上げ点で 1 回取得したプロファイルを飛行中固定するのと同じ扱いである。
  \item 破裂高度を固定している。したがってここで出る散らばりは
        \textbf{風だけに由来する}もので、気球が早く割れたり粘ったりする効果は含まない。
\end{itemize}

\paragraph{向きの規約に注意}
落下位置の方位は「機体が\textbf{進んでいく}方位」で表す。
風向（風が\textbf{吹いてくる}方位）とはちょうど \ang{180} 逆である。
西風は \ang{270} だが、機体を運ぶ先は \ang{90} である。
両者を取り違えると着地点が打ち上げ地点の反対側に出てしまうので、
コードでも本報告でも別々の名前を使っている
（\texttt{dir\_vector\_mean} と \texttt{drift\_bearing\_deg}）。

\subsection{取得したデータが本当に正しいかの検証}

このパイプラインで最も重篤かつ最も気づきにくい失敗は、
\textbf{取り出した領域がずれていること}である。
緯度軸が北から南へ並んでいるため、範囲指定を逆向きに書くと、
例外は出ずに空の配列が返る。あるいは南半球の同じ形の領域が返ってくる。
どちらの場合も、その後の計算はすべて正常に完走し、
「もっともらしいが間違った」数字の集合ができあがる。

そのため、ダウンロードした全ファイル（2160 個）に対して、次を毎回検査している。

\begin{itemize}
  \item 座標が要求した範囲に格子 1 個分の余裕内で収まっているか
  \item 時刻数が期待通りで、6 時間刻みに欠落がないか
  \item 緯度が昇順に直っているか
  \item $|u|, |v| < \SI{200}{\metre\per\second}$、ジオポテンシャル高度が妥当な範囲か
  \item 全欠測でないか、欠測率が 1 \% 未満か
\end{itemize}

検査の結果はファイルごとの付随 JSON に記録してあるので、
後から「本当に全ファイルで検査が走ったか」を監査できる。

\begin{whybox}{検査で見つかった実例}
最初の試験取得で、1 月のジオポテンシャル高度に $-73$ gpm という負の値が出て検査に引っかかった。
調べると、これは\textbf{正しいデータ}だった。
発達した低気圧の中心では \SI{1000}{\hecto\pascal} 面が海面より下に潜る。
検査の下限を 0 にしていた側が間違っていたので、そちらを直した。

検査は間違いを見つけるためだけでなく、
自分の思い込みを見つけるためにも働く。
\end{whybox}

\section{この統計はどこまで信じてよいか}
\label{sec:validation}

前章までで「月を 2 分割した半月」という単位で統計をとった。
しかし境界の位置は、打ち上げ計画が「3 月前半に上げたい」という形で立つという
\textbf{人間側の都合}で決めたものである。
データからその細かさの違いを読み取れるかどうかは、別の話である。

本章はその点を確かめる。\textbf{15 日という幅はこの章の結果から決めた}ので、
順序としては本章が先で、前章までの統計が後になる
（\ref{sec:weeks} 節）。
当初は 7 日 $\times$ 4 に分けており、本章の結論を受けて改めた。
7 日版は \ref{sec:appendix-week} に残してある。

結果の図を見る前に読んでほしい。

\subsection{問いの立て方}

「傾向があるか」という問いは、そのままでは答えても意味がない。
200 hPa の風速が冬 \SI{75}{\metre\per\second}、夏 \SI{23}{\metre\per\second} なのだから、
季節変化が実在することは検定するまでもない。

意味のある問いはこうである。

\begin{quote}
\textbf{どの時間単位まで細かく見てよいのか。}
日単位か、週単位か、旬か、月か、季節か。
どこから先が信号ではなくノイズなのか。
\end{quote}

これは検定ではなく\textbf{モデル選択}の問題であり、決着のつく形がある。

\begin{whybox}{なぜ決着がつくのか}
ビンを細かくすると、手元の 10 年にはよく合うようになる。
一方でビンあたりの標本数は減り、別の年には合わなくなる。

そこで\textbf{1 年抜き交差検証}を使う。
10 年のうち 1 年を除き、残り 9 年で気候値を作り、除いた年を予測させる。
これを 10 通り繰り返す。ビン幅を横軸にとると、

\begin{itemize}
  \item \textbf{訓練誤差}（同じ年を予測したときの誤差）はビンを細かくするほど単調に下がり続ける
  \item \textbf{検証誤差}（除いた年を予測したときの誤差）はどこかで底を打ち、それより細かいと上がる
\end{itemize}

\textbf{検証誤差が最小になる幅が、データから違いを読み取れる最も細かい時間単位}である。
それより細かく分けると、季節の違いではなく偶然のばらつきを区別していることになる。
\end{whybox}

\subsubsection{10 回ぶんをどうまとめるか}
\label{sec:val-folds}

検証年は固定していない。2016 年から 2025 年までが\textbf{順番に 1 回ずつ}検証年になり、
そのたびに残り 9 年で気候値を作り直す。同じ年を 2 回使うことはない
（\texttt{timescale.py:331}）。

したがって 1 つのビン幅につき誤差が 10 個出る。報告している RMSE はこれを二乗平均したもので、

\begin{equation}
  \mathrm{RMSE} = \sqrt{\frac{1}{10}\sum_{j=1}^{10} \mathrm{RMSE}_j^{\,2}}
  \label{eq:fold-pool}
\end{equation}

$j$ は「$j$ 番目の年を抜いた回」を指す。
つまり本章の数字はすべて\textbf{10 年ぶんを通した成績}であって、
特定の年を予測した成績ではない。

この 10 個は、あとで誤差棒にも使う。
\ref{sec:val-answer} 節の図に付いている帯は、この 10 個のばらつきから作った標準誤差である。
ただし\textbf{年ごとの成績のばらつき}をそのまま使うのではない。
同じ 10 年がどのビン幅も採点しているので、
「ビン幅 A とビン幅 B の差」を年ごとにとってから、その差の標準誤差を計算する。
こうすると、年による成績の良し悪しが差の両側で打ち消し合い、
ビン幅どうしの比較がずっと精密になる（対応のある比較）。
年ごとのばらつきをそのまま使うと、本当にある差がその中に埋もれてしまう。

予測する対象は\textbf{落下位置}にした。
第 \ref{sec:drift} 節の飛行モデルで、各解析時刻の風のプロファイルから
着地変位を厳密に計算できるためである。
風そのものについても同じことを行い、高度別の結果を \ref{sec:val-level} 節に示す。

\subsection{本章で使う 2 つの目盛り}
\label{sec:val-metrics}

以降の図と表は、誤差を \textbf{RMSE}、改善の度合いを\textbf{スキル}で表す。
どちらも本章だけで使う量なので、先に定義しておく。

\subsubsection{「予測」と「実際」が何を指すか}
\label{sec:val-what}

先に、以降の式に出てくる 2 つの語をはっきりさせておく。
どちらも日常語と意味がずれているので、ここを取り違えると章全体を誤解することになる。

\paragraph{「予測」は気候値の引き当てであって、機械学習ではない。}
ここでいう予測モデルは、\textbf{該当するビンに入る訓練年の落下点を平均しただけ}のものである
（\texttt{timescale.py} の \texttt{\_group\_means} と \texttt{\_predict}）。

たとえばビン幅が 7 日で、\textbf{ある年}の 3 月 10 日の落下点を予測する場合、
モデルが答えるのは「\textbf{その年を除く}残り 9 年ぶんの、
3 月 8--14 日に入る全解析の平均落下点」である。
それだけである。学習も重みも反復も無く、\textbf{引き当て表}と言ってよい。

除く年は固定していない。10 年が順番に 1 回ずつ検証年になる
（\ref{sec:val-folds} 節）。

\begin{whybox}{なぜわざと単純なモデルなのか}
本章が知りたいのは「どのくらい細かい時間単位までデータから違いを読み取れるか」であって、
「落下点をどこまで正確に予測できるか」ではない。

複雑なモデルを使うと、出てきた数字が
\textbf{データに構造があるからなのか、モデルが賢いからなのか}が分からなくなる。
ビンの平均という、それ以上単純にしようのないモデルに固定しておけば、
ビン幅を変えたときの誤差の動きは、データの側の性質だけを反映する。

唯一の例外が \ref{sec:val-harmonic} 節の調和解析で、
そこでは $\sin$ と $\cos$ を説明変数にした最小二乗回帰を使う。
これも係数を一次方程式で一度に解くだけで、学習と呼べるものではない。
\end{whybox}

\paragraph{「実際」は実際に飛ばしたスペースバルーンの着地点ではない。}
スペースバルーンは一度も飛ばしていない。本報告に観測された着地点は 1 つも入っていない。

「実際」とは、\textbf{その解析時刻に JRA-3Q が与える実際の風のプロファイルを、
第 \ref{sec:drift} 節の飛行モデルに通して得た落下点}である。
2020 年 3 月 10 日 00 UTC なら、その時刻・その格子点の 38 層ぶんの風を使って
計算した着地変位が「実際」になる。

\begin{table}[H]
\centering
\caption{式\eqref{eq:rmse} の 2 つの量が何から作られているか}
\label{tab:pred-vs-true}
\begin{tabular}{@{}lll@{}}
\toprule
 & 何か & どこから来るか \\
\midrule
予測 & そのビンの平均落下点 & 検証年を除く 9 年ぶんの平均 \\
実際 & その解析時刻の落下点 & その時刻の実際の風 $+$ 飛行モデル \\
\bottomrule
\end{tabular}
\end{table}

したがって本章が測っているのは、
\textbf{「そのビンのふつう」が「その日の実際」からどれだけ離れているか}である。
飛行モデルの誤差や予報の誤差はここに入っていない。
飛行モデルは両側で同じものを使っているので差し引かれ、
予報はそもそも使っていない。
シミュレータ自身の誤差は第 \ref{sec:fidelity} 章で別に測る。

\subsubsection{RMSE とスキル}

\paragraph{RMSE（二乗平均平方根誤差）}
予測が平均してどれだけ外れるかを 1 つの数字にしたものである。
名前の通り、後ろから順に、誤差をとり（error）、二乗し（square）、
平均し（mean）、平方根に戻す（root）。

\begin{equation}
\mathrm{RMSE} = \sqrt{\frac{1}{N}\sum_{i=1}^{N}\left(\text{予測}_i - \text{実際}_i\right)^2}
\label{eq:rmse}
\end{equation}

$N$ は検証に使った解析の総数（10 年 $\times$ 6 時間ごと $\times$ 格子点数）で、
$i$ はその 1 つ 1 つを指す。
2 つの語の意味は \ref{sec:val-what} 節のとおりである。

平方根で戻すので\textbf{単位は元のデータと同じ}になる。
本章の RMSE が \si{\kilo\metre} で書かれているのはそのためで、
「\SI{64}{\kilo\metre}」は「落下点の予測が典型的に \SI{64}{\kilo\metre} 外れる」と読んでよい。

\begin{whybox}{なぜ二乗するのか、なぜ絶対値ではないのか}
誤差をそのまま平均すると、東に \SI{50}{\kilo\metre} のずれと
西に \SI{50}{\kilo\metre} のずれが打ち消し合ってゼロになってしまう。
二乗すれば符号が消えるので、打ち消し合わない。

符号を消すだけなら絶対値の平均でもよい。
二乗を選ぶと、\textbf{大きな誤差ほど重く数えられる}という違いが出る。
\SI{10}{\kilo\metre} の誤差 4 回より、\SI{40}{\kilo\metre} の誤差 1 回のほうが悪く評価される。

スペースバルーンの回収では、この重みづけのほうが実情に合う。
たいていは近くに落ちるが年に一度 \SI{200}{\kilo\metre} 先に落ちる予測は、
いつも \SI{20}{\kilo\metre} ずれる予測より扱いに困る。
\end{whybox}

落下位置は東西と南北を持つ 2 次元の量なので、ここでは\textbf{ベクトル RMSE}を使う。
東に \SI{10}{\kilo\metre}、北に \SI{10}{\kilo\metre} ずれたとき、
これは \SI{14}{\kilo\metre} の誤差 1 個であって、\SI{10}{\kilo\metre} の誤差 2 個ではない。
そこで 2 成分の二乗を成分間で平均せずに\textbf{足してから}平方根をとる
（\texttt{stats/timescale.py} の \texttt{\_sq\_err}）。
成分ごとに平均してしまうと、答えが $\sqrt{2}$ 倍だけ小さく出る。
つまりここでいうベクトル RMSE は、
\textbf{予測落下点と実際の落下点の直線距離}を二乗平均したものである。

\paragraph{スキル}
RMSE は \si{\kilo\metre} なので、領域どうしを比べにくい。
和歌山の \SI{64}{\kilo\metre} と北海道の \SI{78}{\kilo\metre} を並べても、
そもそも飛ぶ距離が違うので、どちらの気候値が役に立っているかは分からない。

そこで\textbf{何も知らない場合からの改善率}に直す。

\begin{equation}
\text{スキル} = 1 - \frac{\mathrm{RMSE}}{\mathrm{RMSE}_{\text{季節を知らない場合}}}
\label{eq:skill}
\end{equation}

スキル 0 は「時期を知っても何も改善しない」、
スキル 1 は「誤差がゼロになる」を意味する。
本章で「誤差削減 \SI{46}{\percent}」と書いてあるのはスキル \num{0.46} のことである。

基準となる「季節を知らない場合」は、\textbf{常に 1 年を 1 つのビンにした予測}、
つまり 10 年間の平均落下点をいつでも答えるモデルに固定してある
（\texttt{timescale.py:305-309}）。
比べるビン幅の一覧の中でいちばん粗いものを基準にすると、
一覧を変えるたびに基準が動いて、別々の実行を比べられなくなるためである。

\begin{whybox}{スキルは 1 にならない}
同じ 3 月前半でも、年によって、日によって天気は違う。
気候値は「そのビンのふつう」しか答えられないので、
当日の天気で決まる分の誤差は、どれだけビンを細かくしても残る。

\ref{sec:val-variance} 節で見るように、この下限はかなり高いところにある。
最良の粒度でも説明できるのは落下位置の分散の \SI{70}{\percent}（和歌山）、
\SI{72}{\percent}（愛媛）、\SI{30}{\percent}（北海道）で、残りは原理的に説明できない。
\end{whybox}

\subsection{手法が壊れていないことを先に確かめる}
\label{sec:val-placebo}

検定が甘ければ、ノイズにも「信号がある」と答えてしまう。
そこで結果を見る前に、\textbf{曜日}を予測変数にして同じ手続きを回した。

自由大気は今日が月曜か火曜かを知らない。
したがって曜日には、原理的にスキルが出てはならない。

\begin{table}[H]
\centering
\caption{曜日プラセボ検定。ここに 0 でない値が出たら、以下の結果は信用できない。}
\begin{tabular}{lr}
\toprule
領域 & 曜日のスキル \\
\midrule
北海道 & $-0.0002$ \\
和歌山 & $-0.0001$ \\
愛媛 & $-0.0001$ \\
\bottomrule
\end{tabular}
\end{table}

3 領域とも絶対値が $0.0002$ 以下である。検定は甘くない。

同じ仕組みが\textbf{検出力}を持つことも確かめてある。
合成データに曜日依存の信号を人工的に埋めると、スキル 0.9 以上で検出する
（\texttt{tests/test\_timescale.py}）。
つまり上の 0 は「何も検出できない鈍い手法」の結果ではなく、
「検出できる手法が、正しく何も見つけなかった」結果である。

\subsection{答え}
\label{sec:val-answer}

\begin{figure}[H]
\centering
\includegraphics[width=\linewidth]{validation/V1-wakayama_skill_curve.pdf}
\caption[和歌山の解像度曲線]{和歌山。左が全範囲、右が \num{61} 日以下の拡大。
灰色の破線が訓練誤差、青が検証誤差。
訓練誤差は細かくするほど下がり続けるのに、検証誤差は底を打って上がる。
この開きが過適合である。帯は最良との差の $\pm 2$ 標準誤差で、
\num{10} 回の 1 年抜きから作った\textbf{対応のある}比較の不確かさである
（\ref{sec:val-folds} 節）。年ごとの成績のばらつきではない。}
\label{fig:v1-wakayama}
\end{figure}

\begin{figure}[H]
\centering
\includegraphics[width=\linewidth]{validation/V1-hokkaido_skill_curve.pdf}
\caption[北海道の解像度曲線]{北海道。同じ形だが、底が浅い。
季節を知ることで減らせる誤差が和歌山よりずっと小さい。}
\label{fig:v1-hokkaido}
\end{figure}

\begin{figure}[H]
\centering
\includegraphics[width=\linewidth]{validation/V1-ehime_skill_curve.pdf}
\caption[愛媛の解像度曲線]{愛媛。和歌山とほぼ同じ形で、
\num{7} 日と \num{15} 日が並んで底をつくり、\num{30} 日で上がり始める。}
\label{fig:v1-ehime}
\end{figure}

\begin{table}[H]
\centering
\caption{時間解像度の答え。誤差は落下位置のベクトル RMSE。}
\begin{tabular}{lrrr}
\toprule
 & 北海道 & 和歌山 & 愛媛 \\
\midrule
季節情報なしの誤差 [\si{\kilo\metre}] & 93.6 & 118.2 & 118.7 \\
到達できる最小誤差 [\si{\kilo\metre}] & 78.1 & 64.9 & 63.2 \\
誤差の削減 & 17 \% & \textbf{45 \%} & \textbf{47 \%} \\
分散でいえば & 30 \% & \textbf{70 \%} & \textbf{72 \%} \\
\midrule
最良のビン幅 & 15 日 & 15 日 & 15 日 \\
\textbf{区別できないビン幅の範囲} & \textbf{7--30 日} & \textbf{5--15 日} & \textbf{5--15 日} \\
\bottomrule
\end{tabular}
\end{table}

同じ 10 年で全部のビン幅を評価しているので、
ビン幅どうしの比較は\textbf{対応のある}比較になる。
年ごとの誤差のばらつき（\SI{3}{}--\SI{5}{\kilo\metre}）よりはるかに細かく判定できる。

\begin{whybox}{同点のときは粗いほうを採る}
検証誤差の曲線は最小点の近くで平らである。
たとえば和歌山の \num{7} 日と \num{15} 日は \SI{0.00}{\kilo\metre} 差で並ぶ。
このとき単純に最小値を取ると、たまたま小数第 3 位が小さかったほうを
「データから読み取れる最も細かい幅」として報告してしまう。

そこで交差検証の慣例である\textbf{1 標準誤差ルール}を使った。
最小値から 1 標準誤差以内に収まるもののうち、\textbf{最も粗いもの}を採る。
粗いほうがデータに対して弱い主張をしており、
区別できない 2 つのうち強いほうを選ぶ理由がないからである。

本章の「区別できないビン幅の範囲」はこの基準で決めている。
\end{whybox}

\begin{table}[H]
\centering
\caption{最良のビン幅との差（対応のある比較）。
SE の 2 倍を超えていれば区別できる。太字が本報告の採用値。
3 領域とも 15 日が最良なので差はゼロになる。}
\begin{tabular}{lrrrrrr}
\toprule
 & \multicolumn{2}{c}{北海道} & \multicolumn{2}{c}{和歌山} & \multicolumn{2}{c}{愛媛} \\
\cmidrule(lr){2-3}\cmidrule(lr){4-5}\cmidrule(lr){6-7}
ビン幅 & 差 [\si{\kilo\metre}] & SE の倍数 & 差 [\si{\kilo\metre}] & SE の倍数 & 差 [\si{\kilo\metre}] & SE の倍数 \\
\midrule
1 日（日）  & $+2.22$ & $+6.9$ & $+1.50$ & $+6.5$ & $+1.46$ & $+6.4$ \\
7 日（週）  & $+0.29$ & $+1.4$ & $+0.00$ & $+0.0$ & $+0.02$ & $+0.1$ \\
\textbf{15 日（半月）} & $\mathbf{0.00}$ & $\mathbf{0.0}$ & $\mathbf{0.00}$ & $\mathbf{0.0}$ & $\mathbf{0.00}$ & $\mathbf{0.0}$ \\
30 日（月） & $+0.08$ & $+0.5$ & $+1.61$ & $+5.4$ & $+1.74$ & $+5.8$ \\
\bottomrule
\end{tabular}
\end{table}

読み取れることは 3 つある。

\begin{itemize}
  \item \textbf{日単位は、どの領域でも最良の幅より明らかに悪い。}
        最良より $+6.4$ から $+6.9$ 標準誤差ぶん悪い。
        1 日ビンには 10 年ぶんでも 40 標本しか入らず、
        自己相関を考えれば実質 7 標本程度である。
  \item \textbf{和歌山と愛媛では月では粗すぎる。} 月単位は最良より $+5.4$ SE（和歌山）、
        $+5.8$ SE（愛媛）悪い。
        亜熱帯ジェットの南北移動が速いためで、
        月単位で平均すると遷移期の変化が消えてしまう。
  \item \textbf{北海道では月で足りる。} 月単位は最良と $+0.5$ SE、すなわち区別できない。
        \textbf{北海道について週単位の統計を作っても、月単位より予測がよくなるわけではない。}
\end{itemize}

\begin{figure}[H]
\centering
\includegraphics[width=\linewidth]{validation/V6_region_comparison.pdf}
\caption[3 領域の比較]{左は絶対誤差、右は季節情報なしと比べた改善率。
和歌山と愛媛は下げ幅が大きく底も鋭く、2 本の曲線はほとんど重なる。
北海道は全体に平坦で、月から週へ動かしても何も変わらない。
点線は 1 日・週・月の位置である。}
\label{fig:v6}
\end{figure}

\subsection{ビンではなく滑らかな曲線で見る}
\label{sec:val-harmonic}

ここまではビン幅を変えて比べてきた。しかしビンという形そのものにも
疑う余地がある。そこで\textbf{まったく別の形のモデル}を同じ土俵に乗せた。

\subsubsection{2 つのモデルを式で並べる}

ビンによる気候値は、こう書ける。$d$ を年内通日、$\hat{\mathbf{y}}(d)$ を
その日の落下位置の予測とすると、

\begin{equation}
  \hat{\mathbf{y}}_\text{ビン}(d) = \sum_{b=1}^{B} \mathbf{c}_b \, \mathbf{1}[\,d \in \text{ビン } b\,],
  \qquad B = \left\lceil \frac{365}{w} \right\rceil
  \label{eq:bin-model}
\end{equation}

$\mathbf{1}[\cdot]$ は「その日がビン $b$ に入れば 1、そうでなければ 0」。
$\mathbf{c}_b$ はビン $b$ の平均落下位置で、これが推定するパラメータである。
つまり\textbf{階段関数}であり、$w = 7$ 日なら段が $B = 52$ 個できる。

調和回帰は、同じものを三角関数の和で書く。

\begin{equation}
  \hat{\mathbf{y}}_\text{調和}(d) = \mathbf{a}_0
    + \sum_{j=1}^{K}\Bigl[\,\mathbf{a}_j \cos\bigl(2\pi j\, \tfrac{d}{365}\bigr)
    + \mathbf{b}_j \sin\bigl(2\pi j\, \tfrac{d}{365}\bigr)\Bigr]
  \label{eq:harmonic-model}
\end{equation}

$j = 1$ が 1 年周期、$j = 2$ が半年周期、というように $j$ 番目が
年を $j$ 等分する周期を表す。推定するのは $\mathbf{a}_0, \dots, \mathbf{a}_K,
\mathbf{b}_1, \dots, \mathbf{b}_K$ の $2K + 1$ 組である。
実装は \texttt{timescale.py} の \texttt{harmonic\_design} で、
式\eqref{eq:harmonic-model} の右辺の列を並べた行列を作り、最小二乗で解く。

\begin{whybox}{これはフーリエ「変換」ではない}
FFT は手元の信号を周波数に分解する道具で、予測はしない。
ここでやっているのは\textbf{回帰}である。
式\eqref{eq:harmonic-model} を 9 年ぶんのデータに最小二乗で当てはめ、
残した 1 年を予測させ、その誤差を測る。
ビンのときとまったく同じ 1 年抜き交差検証（\ref{sec:val-folds} 節）にかける。

式\eqref{eq:bin-model} も式\eqref{eq:harmonic-model} も、
\textbf{どちらも「$d$ の関数を基底の線形結合で表す線形モデル」}である。
違うのは基底だけで、前者は矩形波、後者は三角関数。
同じ土俵に乗るからこそ、誤差を直接比べられる。
\end{whybox}

\subsubsection{なぜこれを見る必要があったか}

理由が 4 つある。

\paragraph{1. ビンの切り位置が恣意的である。}
式\eqref{eq:bin-model} の $\mathbf{1}[\cdot]$ は、1・8・15・22 日で切ると
決めた瞬間に定まる。大気にはその日付に何も起きていない。
位相を 3 日ずらせば、同じ「7 日幅」でも別の数字が出る。
つまりビンの掃引では、\textbf{幅の効果と切り位置の効果が分離できない}。
式\eqref{eq:harmonic-model} には境界が 1 つもないので、この任意性がない。

\paragraph{2. 幅が飛び飛びにしか試せない。}
掃引したのは \num{1}、\num{2}、\num{3}、\num{5}、\num{7}、\num{10}、\num{15}、\num{30}、\num{61}、\num{91}、\num{365} 日の 11 点だけである。
一方 $K$ は 0 から 26 まで 1 刻みで動かせる。
$K$ 次までの調和が表せる最短のスケールは

\begin{equation}
  \tau(K) = \frac{365}{2K}\ \text{日}
  \label{eq:harmonic-resolution}
\end{equation}

で（$\tau$ より短い変動は $K$ 次までの和では作れない）、
$K = 6$ なら \num{30} 日、$K = 12$ なら \num{15} 日、$K = 26$ なら \num{7} 日にあたる。
分母の 2 はナイキストの関係で、1 周期を表すには山と谷の 2 点が要ることによる。

\paragraph{3. パラメータの数で比べられる。}
この比較で結論が決まった。

\begin{table}[H]
\centering
\caption{同じ時間分解能を得るのに要るパラメータ数（1 成分あたり）。
落下位置は 2 成分なので実際はこの 2 倍。}
\label{tab:param-count}
\begin{tabular}{@{}lrlr@{}}
\toprule
\multicolumn{2}{c}{ビン} & \multicolumn{2}{c}{調和} \\
\cmidrule(lr){1-2}\cmidrule(lr){3-4}
幅 [日] & $B$ 個 & 次数 & $2K+1$ 個 \\
\midrule
30 & 12 & $K = 6$  & 13 \\
15 & 24 & $K = 12$ & 25 \\
7  & 52 & $K = 26$ & 53 \\
\bottomrule
\end{tabular}
\end{table}

同じ分解能なら数はほぼ同じである。
ところが\textbf{必要な分解能が違った}。
調和は $K = 6$（パラメータ 13 個）で頭打ちになるのに、
ビンは \num{15} 日（24 個）まで細かくしないと最小にならない。
少ない自由度で同等以上の精度が出るということは、
\textbf{実在する構造が滑らかだ}という証拠になる。

\paragraph{4. 階段関数は物理的に正しくない。}
季節変化は連続である。式\eqref{eq:bin-model} は境界ごとに不連続を主張しており、
その主張が害になっていないかを確かめる必要があった。

\begin{figure}[H]
\centering
\includegraphics[width=0.85\linewidth]{validation/V2-wakayama_harmonic.pdf}
\caption[調和解析（和歌山）]{和歌山。上軸が分解できる時間スケール。
$K$ = 3--6 で頭打ちになり、それ以上は増やしても検証誤差が下がらない。}
\label{fig:v2-wakayama}
\end{figure}

\begin{figure}[H]
\centering
\includegraphics[width=0.85\linewidth]{validation/V2-hokkaido_harmonic.pdf}
\caption[調和解析（北海道）]{北海道。曲線の形は和歌山と同じで、
頭打ちになる位置もほぼ同じ（$K = 5$ と $K = 6$）である。
違うのは下げ幅で、$K = 0$ の \SI{93.6}{\kilo\metre} から
\SI{77.5}{\kilo\metre}（$-17$ \%）までしか下がらない。
和歌山は \SI{118.2}{\kilo\metre} から \SI{64.2}{\kilo\metre}（$-46$ \%）である。
縦軸の範囲が 2 図で違う点に注意。}
\label{fig:v2-hokkaido}
\end{figure}

\begin{figure}[H]
\centering
\includegraphics[width=0.85\linewidth]{validation/V2-ehime_harmonic.pdf}
\caption[調和解析（愛媛）]{愛媛。$K = 5$（分解能 \num{36} 日）で頭打ちになり、
\SI{118.7}{\kilo\metre} から \SI{62.4}{\kilo\metre}（$-47$ \%）まで下がる。
形も深さも和歌山とほぼ同じである。}
\label{fig:v2-ehime}
\end{figure}

\subsubsection{結果}

和歌山の最良は $K = 6$（式\eqref{eq:harmonic-resolution} で分解能 \num{30} 日）で、
そのときの検証誤差は \SI{64.2}{\kilo\metre}。
これは\textbf{どのビン幅よりも小さい}（最良のビンは \SI{64.9}{\kilo\metre}）。
北海道も同様で、$K = 5$（分解能 \num{36} 日）の \SI{77.5}{\kilo\metre} が
最良のビン（\SI{78.1}{\kilo\metre}）を下回る。
愛媛も $K = 5$ の \SI{62.4}{\kilo\metre} が最良のビン（\SI{63.2}{\kilo\metre}）を下回る。

3 つの図を並べると、領域の違いが曲線の\textbf{深さ}だけにあることが分かる。
頭打ちになる $K$ が 5 と 6 でほとんど変わらないので、
実在する季節変化の\textbf{時間スケール}は 3 領域で共通である。
違うのは、その季節変化がどれだけ落下位置を決めているかである。
和歌山と愛媛は季節を知ることで誤差が 46--47 \% 減るのに対し、北海道は 17 \% しか減らない。
北海道では残りを日々の天気が決めている。

\begin{whybox}{週ビンは「週の構造」を捉えているのではない}
滑らかな 6 次の曲線のほうが、\num{7} 日ビンより少ないパラメータで、
しかもよい精度を出す。

つまり実在するのは\textbf{約 1--2 か月スケールの滑らかな季節変化}であって、
週に固有の構造ではない。
週ビンはその滑らかな曲線を階段で近似しているにすぎず、
階段を細かくすれば曲線に近づくが、曲線そのものを使うほうが効率がよい。

「週単位の傾向がある」という言い方は正確ではない。
正しくは「\textbf{1--2 か月スケールの季節変化があり、
それを週ビンで表現することもできる}」である。
\end{whybox}

\subsection{高度によって答えが違う}
\label{sec:val-level}

\begin{figure}[H]
\centering
\includegraphics[width=\linewidth]{validation/V3-wakayama_by_level.pdf}
\caption[高度別（和歌山）]{左は気候値がどれだけ誤差を減らすか、
右は最適なビン幅（帯は最適と区別できない範囲）。
左の横軸は図\ref{fig:v3-hokkaido} と共通の \SI{60}{\percent} までにしてある。
図ごとに目盛りを自動で決めると、各領域の最良の棒がどちらもパネル一杯になり、
2 つの図を見比べても差が分からなくなるためである。}
\label{fig:v3-wakayama}
\end{figure}

\begin{figure}[H]
\centering
\includegraphics[width=\linewidth]{validation/V3-hokkaido_by_level.pdf}
\caption[高度別（北海道）]{北海道。横軸は図\ref{fig:v3-wakayama} と共通なので、
棒の長さをそのまま比べられる。
どの高度でも和歌山より短く、\SI{100}{\hecto\pascal} と \SI{30}{\hecto\pascal} で
ようやく 2 割を超える程度である。
右の帯は和歌山より幅が広く、最適なビン幅が絞り込めていない。}
\label{fig:v3-hokkaido}
\end{figure}

\begin{figure}[H]
\centering
\includegraphics[width=\linewidth]{validation/V3-ehime_by_level.pdf}
\caption[高度別（愛媛）]{愛媛。横軸は図\ref{fig:v3-wakayama} と共通である。
棒の長さも右の帯も、和歌山とほぼ同じになる。}
\label{fig:v3-ehime}
\end{figure}

\begin{table}[H]
\centering
\caption{高度別の誤差削減（風ベクトル）}
\begin{tabular}{lrrr}
\toprule
気圧面 & 北海道 & 和歌山 & 愛媛 \\
\midrule
850 hPa（約 \SI{1.5}{\kilo\metre}） & 4 \% & 6 \% & 6 \% \\
500 hPa（約 \SI{5.5}{\kilo\metre}） & 5 \% & 23 \% & 24 \% \\
200 hPa（約 \SI{12}{\kilo\metre}）  & 7 \% & 38 \% & 39 \% \\
100 hPa（約 \SI{16}{\kilo\metre}）  & 23 \% & \textbf{49 \%} & \textbf{50 \%} \\
30 hPa（約 \SI{24}{\kilo\metre}）   & 24 \% & 36 \% & 37 \% \\
10 hPa（約 \SI{31}{\kilo\metre}）   & 10 \% & 22 \% & 23 \% \\
\bottomrule
\end{tabular}
\end{table}

\textbf{地上付近の風は、季節からほとんど予測できない。}
\SI{850}{\hecto\pascal} での誤差削減は 4--6 \% しかない。

これは運用に直結する。
\textbf{打ち上げ時の地上風は、本統計ではなく当日の予報で判断しなければならない。}
本統計が役に立つのは上空、特に \SIrange{16}{24}{\kilo\metre} 付近である。

北海道の値が全体に小さいことにも注意が要る。
\SI{200}{\hecto\pascal} で 7 \% しかないので、
\textbf{北海道上空の対流圏の風は、時期を知ってもほとんど予測できない}。
落下位置の削減が 17 \% あるのは、全層を積算する過程でノイズが打ち消し合うためである。

\subsection{季節によって答えが違う}

\begin{figure}[H]
\centering
\includegraphics[width=\linewidth]{validation/V4-wakayama_by_season.pdf}
\caption[季節別（和歌山）]{各季節の\textbf{中で}、時期を細かく知ることの価値。
基準はその季節の平均である。
左の横軸は図\ref{fig:v4-hokkaido}・\ref{fig:v4-ehime} と共通の \SI{40}{\percent} までである。}
\label{fig:v4-wakayama}
\end{figure}

和歌山では、季節の中で時期を細かく知る価値が季節によって大きく違う。

\begin{table}[H]
\centering
\caption{和歌山：季節の中で時期を知る価値}
\begin{tabular}{lrr}
\toprule
季節 & 最適なビン幅 & 季節内での誤差削減 \\
\midrule
冬 (12--2 月) & 365 日（＝区別できない） & \textbf{0 \%} \\
春 (3--5 月)  & 30 日 & 8 \% \\
夏 (6--8 月)  & 5 日 & 26 \% \\
秋 (9--11 月) & 7 日 & \textbf{34 \%} \\
\bottomrule
\end{tabular}
\end{table}

\textbf{冬の中では、時期を細かく知っても何も得られない。}
1 月上旬か 2 月下旬かを区別しても、落下位置の予測は改善しない。
冬はジェットが安定して強い平坦な高原状態で、
その中の違いは日々の天気の変動に埋もれている。

一方、秋は 34 \%、夏は 26 \% の改善がある。
季節が速く移り変わる時期には、週の違いで予測がよくなる。

\textbf{「週単位の統計が要るのは遷移期であって、冬ではない」}というのが結論である。

愛媛もまったく同じ形になる（図\ref{fig:v4-ehime}）。
冬 0 \%、春 9 \%、夏 27 \%、秋 37 \% で、最適なビン幅も夏と秋が 7 日である。

\begin{figure}[H]
\centering
\includegraphics[width=\linewidth]{validation/V4-ehime_by_season.pdf}
\caption[季節別（愛媛）]{愛媛。横軸は図\ref{fig:v4-wakayama} と共通である。}
\label{fig:v4-ehime}
\end{figure}

北海道では、そもそもどの季節でも季節内の改善が小さい。

\begin{figure}[H]
\centering
\includegraphics[width=\linewidth]{validation/V4-hokkaido_by_season.pdf}
\caption[季節別（北海道）]{北海道。横軸は図\ref{fig:v4-wakayama} と共通である。
和歌山では秋が \SI{34}{\percent}、愛媛では \SI{37}{\percent} まで伸びるのに対し、
北海道はどの季節も \SI{5}{\percent} 台にとどまる。
季節の\textbf{中で}時期を細かく知っても、落下位置の予測はどの季節でも改善しない。}
\label{fig:v4-hokkaido}
\end{figure}

\begin{table}[H]
\centering
\caption{北海道：季節の中で時期を知る価値}
\begin{tabular}{lrr}
\toprule
季節 & 最適なビン幅 & 季節内での誤差削減 \\
\midrule
冬 (12--2 月) & 5--365 日（区別できない） & 1 \% \\
春 (3--5 月)  & 5--365 日（区別できない） & 2 \% \\
夏 (6--8 月)  & 7--30 日 & 5 \% \\
秋 (9--11 月) & 3--91 日 & 5 \% \\
\bottomrule
\end{tabular}
\end{table}

\subsection{有意であることと、役に立つことは別}
\label{sec:val-significance}

北海道について、3 つの手法が一見食い違う答えを出した。

\begin{table}[H]
\centering
\caption{「週は月に情報を加えるか」への 3 つの答え（北海道）}
\begin{tabular}{llr}
\toprule
手法 & 何を測っているか & 答え \\
\midrule
ブロック順列検定 & 週ラベルが情報をまったく持たない確率 & $p = 0.010$ \\
対応のある交差検証 & 週にすると予測がどれだけ良くなるか & $+0.5$ SE（改善なし） \\
月内残差の半分割相関 & 別の 10 年でも同じ週の構造が出るか & $r = +0.19$ \\
\bottomrule
\end{tabular}
\end{table}

矛盾ではない。三者は違うことを測っている。

順列検定が言っているのは「週ラベルの効果はゼロではない」である。
標本が 14{,}610 個あれば、どんなに小さな効果でも有意になる。
交差検証が言っているのは「その効果は小さすぎて予測を改善しない」である。
半分割相関が言っているのは「その小さな構造は別の 10 年では再現しない」である。

\begin{whybox}{$p$ 値だけを見てはいけない理由}
もし順列検定だけを行っていたら、
「北海道でも週単位に有意な傾向がある（$p = 0.010$）」と報告していただろう。
それは誤りではないが、\textbf{読み手に誤解を与える}。
実際には、週単位にしても落下位置の予測は月単位より改善しない（最良との差は週 \SI{0.29}{\kilo\metre}、月 \SI{0.08}{\kilo\metre}）。

本章で検定ではなくモデル選択（交差検証）を主軸に据えたのは、
この区別をつけるためである。
和歌山では 3 つの手法が揃って「月では粗い」と答える（$p = 0.002$、
交差検証 $+5.4$ SE で月が劣る、残差相関 $r = +0.54$）。
愛媛も同じく揃う（$p = 0.002$、$+5.8$ SE、$r = +0.55$）。
北海道では揃わない。これは、週単位の効果が統計的には検出できても、予測には役立たない大きさだということを示している。
\end{whybox}

\subsection{気候値で説明できる範囲}
\label{sec:val-variance}

\begin{figure}[H]
\centering
\includegraphics[width=0.9\linewidth]{validation/V5-wakayama_variance.pdf}
\caption[分散分解（和歌山）]{和歌山。落下位置の分散のうち、気候値が説明する割合。
残りは当日の天気である。
\textbf{ビンを細かくするほど割合が増えている}が、これは手元の 10 年への
あてはまりであって、予測が良くなったことを意味しない。
同じ 10 年で測っているので、パラメータを増やせば必ず増える。
交差検証（図\ref{fig:v1-wakayama}）が必要だったのはこのためである。
本文の 70 \% は交差検証で測った値（\SI{69.9}{\percent}）で、
この図の \SIrange{69}{74}{\percent} とは別の量である。
和歌山では両者がほぼ一致する。}
\label{fig:v5-wakayama}
\end{figure}

\begin{figure}[H]
\centering
\includegraphics[width=0.9\linewidth]{validation/V5-hokkaido_variance.pdf}
\caption[分散分解（北海道）]{北海道。横軸は和歌山と同じ \SIrange{0}{100}{\percent} である。
気候値が説明するのは \SIrange{32}{39}{\percent} で、
和歌山の \SIrange{69}{74}{\percent} の半分以下にとどまる。
落下位置の 3 分の 2 を当日の天気が決めているということで、
気候値でできることの上限がそもそも低い。
なお本文の 30 \% は交差検証で測った値（\SI{30.3}{\percent}）で、
この図の値より低い。
和歌山では両者がほぼ一致するのに北海道では開くこと自体が、
細かいビンのあてはまりが別の年には持ち越せていない証拠である。}
\label{fig:v5-hokkaido}
\end{figure}

\begin{figure}[H]
\centering
\includegraphics[width=0.9\linewidth]{validation/V5-ehime_variance.pdf}
\caption[分散分解（愛媛）]{愛媛。気候値が説明するのは \SIrange{71}{75}{\percent} で、
和歌山とほぼ同じである。交差検証で測った値は \SI{72}{\percent}。}
\label{fig:v5-ehime}
\end{figure}

最良の粒度でも、説明できるのは分散の 70 \%（和歌山）、72 \%（愛媛）、30 \%（北海道）である。
残りは当日の総観規模の天気で決まり、気候値では原理的に説明できない。

\textbf{気候値は計画のための道具であって、予報の代わりではない。}
和歌山の落下位置は、時期を知ることで散らばりが
\SI{118}{\kilo\metre} から \SI{65}{\kilo\metre} に縮むが、
\SI{65}{\kilo\metre} は残る。愛媛でも \SI{119}{\kilo\metre} から \SI{63}{\kilo\metre} で、
やはり \SI{63}{\kilo\metre} は残る。

\subsection{この検証で前章までの結果はどう変わるか}

結論として、本報告の集計単位は\textbf{半月}（15 日）とした。
そのうえで次の点を明示する。

\begin{itemize}
  \item \textbf{和歌山と愛媛の半月単位の統計には、データの裏付けがある}（月より有意に良い）。
  \item \textbf{北海道の半月単位の統計は、月単位と区別できない}。
        半月ごとの細かい違いは読み取らないこと。
  \item \textbf{日単位の統計は作らない}。3 領域とも最良の幅より明らかに悪い。
  \item \textbf{冬の半月ごとの違いは読み取らない}。季節の中では区別できない。
  \item \textbf{地上付近の風にこの統計を使わない}。誤差削減が 5 \% 前後しかない。
\end{itemize}

第 \ref{sec:limits} 章に、この検証を踏まえた限界の一覧を置いた。

\section{結果：北海道}
\label{sec:hokkaido}

対象は \ang{40.0}--\ang{46.5}N / \ang{138.0}--\ang{149.0}E、
\ang{0.375} 格子で 570 点。周辺の日本海・オホーツク海・太平洋を含む。
代表地点は \texttt{sites.json} に登録した \textbf{VOISING}
（\ang{42.1823}N, \ang{142.8510}E、日高・浦河の東）である。
最寄り格子までの距離は \SI{4.7}{\kilo\metre}。
どの地点を代表にするかはコード側で名前指定してあり
（\texttt{csv\_out.REFERENCE\_SITE}）、
\texttt{sites.json} の並び順には依存しない。

\subsection{まず標本数を確認する}

\begin{figure}[H]
\centering
\includegraphics[width=0.95\linewidth]{hokkaido/F0_sample_counts.pdf}
\caption[標本数]{各ビンに入った解析値の数。前半は必ず 15 日ぶんで 600 個、後半は月の長さで変わり 1 月は 640 個、2 月は閏年が 3 回あるので 532 個になる。この図が合わなければ取得か集計が壊れているので、他のどの図よりも先に確認する。}
\label{fig:f0-hokkaido}
\end{figure}


図 \ref{fig:f0-hokkaido} の標本数は設計どおりである。
前半が各 600 個、後半は月の長さに応じて 532--640 個。
一方、有効標本数（第 \ref{sec:neff} 節）の中央値は \textbf{77.2} で、
生の標本数の 13 \% にすぎない。
つまり「10 年分の 600 観測」は、独立な天気としては 77 個程度の情報しか持っていない。
以降の平均値は、この前提で読む必要がある。

\subsection{年間を通した姿}

\begin{figure}[H]
\centering
\includegraphics[width=1.0\linewidth]{hokkaido/F1_alt_half_speed.pdf}
\caption[時期×高度の風速と風向]{上：平均風速。下：平均風向（吹いてくる向き）。風向には円環用のカラーマップを使っている。順次カラーマップを使うと \ang{0} と \ang{360} の境目に実在しない断層が現れるためである。灰色は定常度が 0.3 未満で塗りを外した領域で、そのビンではそもそも風向が定まっていない。}
\label{fig:f1-hokkaido}
\end{figure}

\begin{figure}[H]
\centering
\includegraphics[width=1.0\linewidth]{hokkaido/F3_profiles_by_month.pdf}
\caption[月別の高度プロファイル]{月ごとの風速プロファイル。線は各ビンの中央値、灰色の帯はその月の 2 ビンを通した p10--p90 の広がり。帯の幅が平均からのずれの大きさを表しており、線だけを見て計画を立ててはいけない。}
\label{fig:f3-hokkaido}
\end{figure}


対流圏界面付近（\SI{200}{\hecto\pascal}、約 \SI{11.8}{\kilo\metre}）の平均風速は
冬 \SI{40.2}{\metre\per\second}、春 \SI{35.0}{\metre\per\second}、
夏 \SI{29.9}{\metre\per\second}、\textbf{秋 \SI{44.5}{\metre\per\second}} である。

冬ではなく秋が最も強いのは、緯度による。
東アジアのジェットは冬に最も強くなるが、同時に南へ下がる。
北海道の緯度帯には、ジェットが南下しきる前の 10--11 月に核が通過する。
後述する和歌山（冬に \SI{75.5}{\metre\per\second}）と比べると、この違いははっきりする。

定常度は対流圏では総じて高い（\SI{200}{\hecto\pascal} で 0.81--0.96）。
つまり\textbf{北海道上空の対流圏の風向は年間を通して安定している}。
一方で、上層には風向が定まらない時期がある。

\subsection{上層では風向が定まらない時期がある}

\begin{figure}[H]
\centering
\includegraphics[width=1.0\linewidth]{hokkaido/F2_alt_half_constancy.pdf}
\caption[定常度]{定常度（ベクトル平均風速 ÷ スカラー平均風速）。1 に近いほど風向が安定し、0 に近いほどビンのなかで風向が入れ替わっている。淡い帯が成層圏の風の反転している高度で、季節とともに上下する。}
\label{fig:f2-hokkaido}
\end{figure}


\SI{10}{\hecto\pascal}（約 \SI{31}{\kilo\metre}、破裂高度帯）の定常度は
夏に 0.99 と極めて高い一方、春は 0.36、秋は 0.31 しかない。
特に定常度が低いのは 10 月前半（0.15）、5 月前半（0.17）、
11 月後半（0.21）、10 月後半（0.23）、3 月前半（0.26）である。

これらは成層圏の風が東西逆転する時期にあたる。
そのビンの 10 年ぶんの解析には「強い西風の日」と「強い東風の日」が混ざっており、
平均するとほぼ打ち消し合う。
図 \ref{fig:f1-hokkaido} の下パネルでこれらのビンが灰色になっているのはそのためで、
\textbf{データが無いのではなく、平均風向という量に意味がない}。

\begin{figure}[H]
\centering
\includegraphics[width=1.0\linewidth]{hokkaido/F7-1_wind_roses.pdf}
\caption[風配図（前半）]{代表格子点の風配図、\textbf{月の前半}（1--15 日）。
平均と楕円では見せられない「分布そのもの」を示す。
ローブが 2 つに割れている場面が、平均風がほぼゼロで分散だけが大きくなる時期にあたる。}
\label{fig:f7-hokkaido-h1}
\end{figure}

\begin{figure}[H]
\centering
\includegraphics[width=1.0\linewidth]{hokkaido/F7-2_wind_roses.pdf}
\caption[風配図（後半）]{同じ格子点の\textbf{月の後半}（16 日--月末）。
前半（図\ref{fig:f7-hokkaido-h1}）と並べると、\textbf{2 枚はよく似ている}。
\SI{30}{\hecto\pascal} の卓越風向は 1・4・7・10 月のいずれも前半と後半で同じである。
これは第 \ref{sec:validation} 章の結論と整合する。
月をさらに細かく割っても、新しい構造は出てこない。}
\label{fig:f7-hokkaido-h2}
\end{figure}

風配図（図 \ref{fig:f7-hokkaido-h1}、\ref{fig:f7-hokkaido-h2}）はこの状況を直接見せる。
平均と分散楕円では 2 つの塊を 1 つの楕円に潰してしまうが、
風配図なら塊が 2 つあればそう見える。

\subsection{冬の 10 hPa が東風であること}

一つ、事前の予想と食い違った点を記しておく。

帯状平均で見れば、中緯度の成層圏は冬に西風（極夜ジェット）である。
ところが本統計では、北海道上空の \SI{10}{\hecto\pascal} は
1 月に東風（\SI{-3.8}{\metre\per\second}）になる。
2016、2017、2018 の各年を個別に見ても、いずれも弱い東風である。
さらに同じ高度を経度方向に見ると、
東へ行くほど東風が強まる勾配がある。

これは誤りではなく、この経度帯の実際の気候である。
真冬の成層圏には北太平洋上に高気圧性の循環（アリューシャン高気圧）が停滞し、
その西縁にあたる日本の東側では帯状平均と逆向きの風になる。
\textbf{帯状平均の知識を、一部の経度だけを切り取った領域にそのまま当てはめてはいけない}という例であり、
このパイプラインの検証項目もその点で修正した。

\subsection{領域の中で答えが変わるか}

\begin{figure}[H]
\centering
\includegraphics[width=1.0\linewidth]{hokkaido/F6_spatial_spread.pdf}
\caption[領域内のばらつき]{領域内での平均風速の幅（最大 $-$ 最小）。これが小さい高度・時期では、領域のどこで上げても風はほぼ同じで、代表 1 地点のプロファイルで語ってよい。大きいところではそうではない。}
\label{fig:f6-hokkaido}
\end{figure}


領域内での平均風速の幅（最大 $-$ 最小）は最大で \SI{31.7}{\metre\per\second} に達する。
これは無視できない大きさで、
\textbf{北海道の領域は「代表 1 地点で語れる」ほど一様ではない}。
南端（\ang{40}N、東北北部）と北端（\ang{46.5}N、宗谷沖）では、
特に冬のジェットの強さが大きく違う。

したがって VOISING の数値を道北や道東にそのまま適用してはいけない。
格子ごとの値は NetCDF と Parquet にあるので、必要な地点を引くこと。

\subsection{地図で見る}

\begin{figure}[H]
\centering
\includegraphics[width=1.0\linewidth]{hokkaido/M1-200_month_grid_200hPa.pdf}
\caption[200 hPa の月別分布図]{\SI{200}{\hecto\pascal}（約 \SI{12}{\kilo\metre}）の平均風を月ごとに並べたもの。塗りが風速、矢印が風向。ジェットが季節とともに動く様子が地図の上で見える。}
\label{fig:m1200-hokkaido}
\end{figure}

\begin{figure}[H]
\centering
\includegraphics[width=1.0\linewidth]{hokkaido/M1-030_month_grid_30hPa.pdf}
\caption[30 hPa の月別分布図]{同じく \SI{30}{\hecto\pascal}（約 \SI{24}{\kilo\metre}）。対流圏とは逆の季節変化をしており、夏に東風が強まる。}
\label{fig:m1030-hokkaido}
\end{figure}

\begin{figure}[H]
\centering
\includegraphics[width=1.0\linewidth]{hokkaido/M3-200_direction_map_200hPa.pdf}
\caption[風向の分布図]{\SI{200}{\hecto\pascal} の平均風向。円環用のカラーマップを使い、定常度 0.3 未満は灰色にしている。灰色の領域では平均風向に意味がない。}
\label{fig:m3200-hokkaido}
\end{figure}

\begin{figure}[H]
\centering
\includegraphics[width=1.0\linewidth]{hokkaido/M2_constancy_map.pdf}
\caption[定常度の分布図]{高度と季節ごとの定常度。濃いところでは風向が安定していて計画が立てやすく、淡いところではビンのなかで風向が入れ替わる。}
\label{fig:m2-hokkaido}
\end{figure}

\begin{figure}[H]
\centering
\includegraphics[width=0.92\linewidth]{hokkaido/M4-01_wind_ellipse_map.pdf}
\caption[風のばらつきを地図に]{1 月前半、\SI{200}{\hecto\pascal} の平均風ベクトルと、その先端に描いた $1\sigma$ の散らばり楕円。丸い楕円はどの向きにもばらつくこと、細長い楕円は主に強さがばらつくことを意味する。}
\label{fig:m401-hokkaido}
\end{figure}


\subsection{落下位置}

\begin{figure}[H]
\centering
\includegraphics[width=1.0\linewidth]{hokkaido/M5_drift_month_grid.pdf}
\caption[月別の平均落下位置]{各格子点から上げた場合の平均落下位置。\textbf{矢印は実寸}で、矢先がそのまま予測落下点である。色が平均の飛行距離。全パネルが同じ枠と同じ色スケールなので、月どうしを直接比べられる。}
\label{fig:m5-hokkaido}
\end{figure}

\begin{figure}[H]
\centering
\includegraphics[width=0.92\linewidth]{hokkaido/M7-01_drift_map.pdf}
\caption[1 月の落下位置]{1 月前半の平均落下位置を拡大したもの。枠は集計対象の領域で、矢先はその外へ出ていることが多い。}
\label{fig:m701-hokkaido}
\end{figure}

\begin{figure}[H]
\centering
\includegraphics[width=0.92\linewidth]{hokkaido/M7-07_drift_map.pdf}
\caption[7 月の落下位置]{同じく 7 月前半。1 月と比べると距離も向きも大きく違う。}
\label{fig:m707-hokkaido}
\end{figure}

\begin{figure}[H]
\centering
\includegraphics[width=1.0\linewidth]{hokkaido/M6_drift_percentile_map.pdf}
\caption[荒れた日の飛行距離]{飛行距離の 90 パーセンタイル。平均 \SI{150}{\kilo\metre} で p90 が \SI{400}{\kilo\metre} という状況と、平均 \SI{150}{\kilo\metre} だけを見て計画することは、運用上まったく違う。}
\label{fig:m6-hokkaido}
\end{figure}

\begin{figure}[H]
\centering
\includegraphics[width=0.92\linewidth]{hokkaido/M8_landing_ellipse_map.pdf}
\caption[代表地点からの落下分布]{1 地点から上げた場合の落下位置を季節別に重ねたもの。楕円は 50 / 90 / 95 \% の範囲で、各解析時刻ごとに求めた着地点の共分散から描いている。散らばりが 2 変量正規に従うという仮定は残る。灰色が陸、白が海である。}
\label{fig:m8-hokkaido}
\end{figure}


飛行モデル（第 \ref{sec:drift} 節、上昇 \SI{5}{\metre\per\second}、
破裂 \SI{30}{\kilo\metre}、降下 \SI{6}{\metre\per\second}）での飛行時間は約 \num{138} 分。
領域平均の飛行距離は年間平均 \SI{143}{\kilo\metre} で、季節差が非常に大きい。

\begin{table}[H]
\centering
\caption{北海道：飛行距離が最短・最長のビン（領域平均）。
方位は機体が\textbf{進む}向きで、風向とは \ang{180} 逆である。}
\begin{tabular}{lrrr}
\toprule
時期 & 平均距離 [\si{\kilo\metre}] & 方位 [度] & p90 [\si{\kilo\metre}] \\
\midrule
7 月 後半 & 62 & 123（東南東） & 102 \\
7 月 前半 & 67 & 102（東南東） & 108 \\
8 月 前半 & 72 & 96（東） & 122 \\
\midrule
12 月 前半 & 215 & 75（東北東） & 292 \\
11 月 後半 & 205 & 78（東北東） & 284 \\
12 月 後半 & 192 & 82（東） & 269 \\
\bottomrule
\end{tabular}
\end{table}

この差は運用の判断に直接つながる。

\begin{itemize}
  \item \textbf{7 月が最も扱いやすい。} 飛行距離が冬の 4 分の 1 になり、
        p90 でも \SI{100}{\kilo\metre} 程度に収まる。
  \item \textbf{11 月から 12 月は避けたほうがよい。} 平均で \SI{200}{\kilo\metre} を超え、
        p90 は \SI{300}{\kilo\metre} に達する。方位は東北東で一定なので、
        十勝から上げた機体はほぼ確実に太平洋に落ちる。
  \item 図 \ref{fig:m8-hokkaido} のとおり、VOISING から上げた場合、
        \textbf{平均落下点はどの季節でも海上}である。
        VOISING は日高の南岸で太平洋に面しており、
        機体は東へ運ばれるので、そのまま沖へ出る。
\end{itemize}

季節で変わるのは「海に落ちるかどうか」ではなく「岸からどれだけ離れるか」である。

\begin{table}[H]
\centering
\caption{VOISING から上げた場合の平均落下点と、そこから最寄りの海岸までの距離。
どの時期も海上だが、夏だけは岸のすぐそばに落ちる。}
\label{tab:voising-offshore}
\sisetup{table-format=3.0}
\begin{tabular}{@{}llSS@{}}
\toprule
月 & ビン & {飛行距離 [\si{\kilo\metre}]} & {岸まで [\si{\kilo\metre}]} \\
\midrule
1 月  & 前半 & 187 & 80 \\
1 月  & 後半 & 184 & 75 \\
4 月  & 前半 & 156 & 67 \\
4 月  & 後半 & 155 & 74 \\
7 月  & 前半 & 71  & 4 \\
7 月  & 後半 & 62  & 2 \\
10 月 & 前半 & 184 & 59 \\
10 月 & 後半 & 185 & 55 \\
\bottomrule
\end{tabular}
\end{table}

\textbf{7 月だけは状況が大きく違う。} 平均落下点が岸から \SIrange{2}{4}{\kilo\metre} しかないので、
散らばり（図 \ref{fig:m8-hokkaido} の楕円）の陸側の裾は陸にかかる。
回収を陸上で行いたいなら、実質的に 7 月しか選択肢がない。
他の季節は \SIrange{55}{80}{\kilo\metre} 沖なので、船を用意するか、回収を諦めるかを決めておく必要がある。

\begin{whybox}{打ち上げ地点を変えるとこの結論は変わる}
これは VOISING という\textbf{特定の地点}の性質であって、北海道の性質ではない。

同じ北海道でも、内陸寄りの地点から上げれば同じ風でも陸上に落ちる。
第 \ref{sec:hokkaido} 章の他の図は領域全体の格子点を扱っているので
地点によらないが、この図と表だけは地点に依存する。

\texttt{out/csv/region\_sites\_hokkaido.csv} に領域内の全 114 地点と
最寄り格子までの距離を出してあるので、別の地点で見たい場合は
\texttt{csv\_out.REFERENCE\_SITE} を変えて再実行すればよい。
\end{whybox}

なお、ここでの散らばりは\textbf{風だけに由来する}。
気球が予定より早く割れたり粘ったりする効果は含まれていないので、
実際の散らばりはこれより大きい。

\subsection{日変化と鉛直シアー}

\begin{figure}[H]
\centering
\includegraphics[width=0.75\linewidth]{hokkaido/F8_diurnal.pdf}
\caption[日変化]{夜（03 JST）と昼（15 JST）の平均風速の差を、季節ごとに高度に対して描いたもの
（\SI{300}{\hecto\pascal} 以下の層のみ、領域平均）。横軸の目盛りに注意してほしい。
差は 3 領域・4 季節のどれでも \SI{1}{\metre\per\second} 未満で、
日々の変動（数十 \si{\metre\per\second}）に比べて小さい。
これが\textbf{時刻を区別せず平均してよい根拠}である。
ビンを時刻で 4 分割すると標本数も 4 分の 1 になるので、
必要のないところで分割しない理由がここにある。}
\label{fig:f8-hokkaido}
\end{figure}

\begin{figure}[H]
\centering
\includegraphics[width=1.0\linewidth]{hokkaido/F9_shear.pdf}
\caption[鉛直シアー]{隣り合う気圧面のあいだの風速差を高度差で割ったもの。
機体そのものはシアーで壊れないが、大きいところでは
高度の推定が少しずれるだけで水平の予測が大きく動く。
最大はジェットの上下と地表付近に出る。}
\label{fig:f9-hokkaido}
\end{figure}


夜（03 JST）と昼（15 JST）の平均風速の差は、最大でも \SI{0.9}{\metre\per\second} である
（図\ref{fig:f8-hokkaido}）。北海道では地表付近より上空（高度 \SI{9}{\kilo\metre} 付近、昼のほうが強い）
のほうが差が大きいが、それでも日々の変動に比べれば小さい。
第 \ref{sec:method} 章で上層の統計を解析時刻で分けなかったのは、
この図が根拠である。

鉛直シアーの最大は \SI{6.7}{(\metre\per\second)\per\kilo\metre}
（5 月前半、高度 \SI{0.1}{\kilo\metre}）。
地表付近を除くと \SI{5.4}{(\metre\per\second)\per\kilo\metre}
（10 月前半、高度 \SI{16.4}{\kilo\metre}）で、対流圏界面のすぐ上にあたる。

\subsection{既存シミュレータの ISA 仮定}

\begin{figure}[H]
\centering
\includegraphics[width=1.0\linewidth]{hokkaido/F10_isa_check.pdf}
\caption[ISA との比較]{気圧から高度への換算。左は JRA-3Q の実際の気圧面高度と、既存シミュレータが使う国際標準大気 ISA。右はその差。ISA は季節も緯度も天気も知らないので、季節によって系統的にずれる。}
\label{fig:f10-hokkaido}
\end{figure}


既存シミュレータ（\texttt{backend/bfps/core/atmosphere.py}）は
気圧から高度への換算に国際標準大気を使っている。
北海道での実際とのずれは最大 \SI{1036}{\metre}（夏）、冬でも \SI{520}{\metre} ある。

これは無視できる量ではない。
たとえば「高度 \SI{30}{\kilo\metre} で破裂」と設定した計算が、
夏には実際の \SI{29}{\kilo\metre} 相当を指してしまう。
高度そのものを扱う場面では、ISA ではなく本統計の
\texttt{alt\_geom\_mean\_m} を使うべきである。

\section{結果：和歌山}
\label{sec:wakayama}

対象は \ang{32.5}--\ang{35.0}N / \ang{134.0}--\ang{137.0}E、
\ang{0.375} 格子で 72 点。紀伊水道から潮岬の南の太平洋までを含む。

北海道と違い、\texttt{sites.json} に和歌山の打ち上げ地点は登録されていない。
そのため本章は\textbf{領域全体の統計}で述べ、
代表地点の表は作っていない。
地点を決めずに 1 点だけ数値を出すと、根拠のない精度を持ち込むことになるためである。
図 \ref{fig:m8-wakayama} だけは、比較のために領域の中央格子点を使っている。

\subsection{年間を通した姿}

\begin{figure}[H]
\centering
\includegraphics[width=0.95\linewidth]{wakayama/F0_sample_counts.pdf}
\caption[標本数]{各ビンに入った解析値の数。前半は必ず 15 日ぶんで 600 個、後半は月の長さで変わり 1 月は 640 個、2 月は閏年が 3 回あるので 532 個になる。この図が合わなければ取得か集計が壊れているので、他のどの図よりも先に確認する。}
\label{fig:f0-wakayama}
\end{figure}

\begin{figure}[H]
\centering
\includegraphics[width=1.0\linewidth]{wakayama/F1_alt_half_speed.pdf}
\caption[時期×高度の風速と風向]{上：平均風速。下：平均風向（吹いてくる向き）。風向には円環用のカラーマップを使っている。順次カラーマップを使うと \ang{0} と \ang{360} の境目に実在しない断層が現れるためである。灰色は定常度が 0.3 未満で塗りを外した領域で、そのビンではそもそも風向が定まっていない。}
\label{fig:f1-wakayama}
\end{figure}

\begin{figure}[H]
\centering
\includegraphics[width=1.0\linewidth]{wakayama/F3_profiles_by_month.pdf}
\caption[月別の高度プロファイル]{月ごとの風速プロファイル。線は各ビンの中央値、灰色の帯はその月の 2 ビンを通した p10--p90 の広がり。帯の幅が平均からのずれの大きさを表しており、線だけを見て計画を立ててはいけない。}
\label{fig:f3-wakayama}
\end{figure}


和歌山の特徴は\textbf{冬のジェットの強さ}である。
\SI{200}{\hecto\pascal}（約 \SI{12.2}{\kilo\metre}）の平均風速は冬に
\SI{75.5}{\metre\per\second}、定常度は 0.98 に達する。
北海道の同じ高度・同じ季節（\SI{40.2}{\metre\per\second}）のほぼ 2 倍である。

これは和歌山が亜熱帯ジェットの真下にあたるためで、
冬に南下した東アジアジェットの核がちょうどこの緯度帯を通る。
定常度 0.98 という値は、
\textbf{冬の 3 か月間、この高度の風はほとんど常に西からの強風だった}ことを意味する。

夏はこれが大きく変わり、\SI{200}{\hecto\pascal} で \SI{22.7}{\metre\per\second}、
定常度 0.57 まで落ちる。ジェットの核が北へ抜けるためである。

\subsection{成層圏}

\begin{figure}[H]
\centering
\includegraphics[width=1.0\linewidth]{wakayama/F2_alt_half_constancy.pdf}
\caption[定常度]{定常度（ベクトル平均風速 ÷ スカラー平均風速）。1 に近いほど風向が安定し、0 に近いほどビンのなかで風向が入れ替わっている。淡い帯が成層圏の風の反転している高度で、季節とともに上下する。}
\label{fig:f2-wakayama}
\end{figure}

\begin{figure}[H]
\centering
\includegraphics[width=1.0\linewidth]{wakayama/F7-1_wind_roses.pdf}
\caption[風配図（前半）]{代表格子点の風配図、\textbf{月の前半}（1--15 日）。
平均と楕円では見せられない「分布そのもの」を示す。
ローブが 2 つに割れている場面が、平均風がほぼゼロで分散だけが大きくなる時期にあたる。}
\label{fig:f7-wakayama-h1}
\end{figure}

\begin{figure}[H]
\centering
\includegraphics[width=1.0\linewidth]{wakayama/F7-2_wind_roses.pdf}
\caption[風配図（後半）]{同じ格子点の\textbf{月の後半}（16 日--月末）。
前半（図\ref{fig:f7-wakayama-h1}）と並べると、ほとんどの月で\textbf{2 枚はよく似ている}。
唯一 10 月だけ \SI{30}{\hecto\pascal} の最頻セクターが前半 \ang{90}、
後半 \ang{248} と入れ替わるが、どちらも頻度は 2 割前後しかない。
つまり向きが変わったというより、\textbf{その時期はどちらの半月も
風向が定まっていない}のであって、最頻セクターを読む意味が薄い。}
\label{fig:f7-wakayama-h2}
\end{figure}

\SI{10}{\hecto\pascal}（約 \SI{31}{\kilo\metre}）は、
夏に \SI{18.2}{\metre\per\second}・定常度 0.99 の\textbf{東風}である。
これは熱帯・亜熱帯の夏の成層圏東風で、
北海道の同じ高度（\SI{13.6}{\metre\per\second}）より強い。

一方、冬と春の定常度はそれぞれ 0.40、0.37 しかない。
定常度が最も低いのは 3 月後半（0.05）、11 月前半（0.11）、
1 月後半（0.18）、3 月前半（0.21）、2 月前半（0.22）である。
北海道と同じく、春と秋の東西逆転期に風向が定まらなくなる。

\subsection{領域の中で答えが変わるか}

\begin{figure}[H]
\centering
\includegraphics[width=1.0\linewidth]{wakayama/F6_spatial_spread.pdf}
\caption[領域内のばらつき]{領域内での平均風速の幅（最大 $-$ 最小）。これが小さい高度・時期では、領域のどこで上げても風はほぼ同じで、代表 1 地点のプロファイルで語ってよい。大きいところではそうではない。}
\label{fig:f6-wakayama}
\end{figure}


領域内での平均風速の幅は最大でも \SI{12.0}{\metre\per\second} である。
北海道の \SI{31.7}{\metre\per\second} と比べるとはるかに小さい。

つまり\textbf{和歌山の領域は、代表 1 地点でおおむね語ってよい}。
領域が狭いこと（北海道の 8 分の 1 の格子点数）に加えて、
緯度の幅が小さくジェットの勾配をまたがないためである。

\subsection{地図で見る}

\begin{figure}[H]
\centering
\includegraphics[width=1.0\linewidth]{wakayama/M1-200_month_grid_200hPa.pdf}
\caption[200 hPa の月別分布図]{\SI{200}{\hecto\pascal}（約 \SI{12}{\kilo\metre}）の平均風を月ごとに並べたもの。塗りが風速、矢印が風向。ジェットが季節とともに動く様子が地図の上で見える。}
\label{fig:m1200-wakayama}
\end{figure}

\begin{figure}[H]
\centering
\includegraphics[width=1.0\linewidth]{wakayama/M1-030_month_grid_30hPa.pdf}
\caption[30 hPa の月別分布図]{同じく \SI{30}{\hecto\pascal}（約 \SI{24}{\kilo\metre}）。対流圏とは逆の季節変化をしており、夏に東風が強まる。}
\label{fig:m1030-wakayama}
\end{figure}

\begin{figure}[H]
\centering
\includegraphics[width=1.0\linewidth]{wakayama/M3-200_direction_map_200hPa.pdf}
\caption[風向の分布図]{\SI{200}{\hecto\pascal} の平均風向。円環用のカラーマップを使い、定常度 0.3 未満は灰色にしている。灰色の領域では平均風向に意味がない。}
\label{fig:m3200-wakayama}
\end{figure}

\begin{figure}[H]
\centering
\includegraphics[width=1.0\linewidth]{wakayama/M2_constancy_map.pdf}
\caption[定常度の分布図]{高度と季節ごとの定常度。濃いところでは風向が安定していて計画が立てやすく、淡いところではビンのなかで風向が入れ替わる。}
\label{fig:m2-wakayama}
\end{figure}

\begin{figure}[H]
\centering
\includegraphics[width=0.92\linewidth]{wakayama/M4-01_wind_ellipse_map.pdf}
\caption[風のばらつきを地図に]{1 月前半、\SI{200}{\hecto\pascal} の平均風ベクトルと、その先端に描いた $1\sigma$ の散らばり楕円。丸い楕円はどの向きにもばらつくこと、細長い楕円は主に強さがばらつくことを意味する。}
\label{fig:m401-wakayama}
\end{figure}


\subsection{落下位置}

\begin{figure}[H]
\centering
\includegraphics[width=1.0\linewidth]{wakayama/M5_drift_month_grid.pdf}
\caption[月別の平均落下位置]{各格子点から上げた場合の平均落下位置。\textbf{矢印は実寸}で、矢先がそのまま予測落下点である。色が平均の飛行距離。全パネルが同じ枠と同じ色スケールなので、月どうしを直接比べられる。}
\label{fig:m5-wakayama}
\end{figure}

\begin{figure}[H]
\centering
\includegraphics[width=0.92\linewidth]{wakayama/M7-01_drift_map.pdf}
\caption[1 月の落下位置]{1 月前半の平均落下位置を拡大したもの。枠は集計対象の領域で、矢先はその外へ出ていることが多い。}
\label{fig:m701-wakayama}
\end{figure}

\begin{figure}[H]
\centering
\includegraphics[width=0.92\linewidth]{wakayama/M7-07_drift_map.pdf}
\caption[7 月の落下位置]{同じく 7 月前半。1 月と比べると距離も向きも大きく違う。}
\label{fig:m707-wakayama}
\end{figure}

\begin{figure}[H]
\centering
\includegraphics[width=1.0\linewidth]{wakayama/M6_drift_percentile_map.pdf}
\caption[荒れた日の飛行距離]{飛行距離の 90 パーセンタイル。平均 \SI{150}{\kilo\metre} で p90 が \SI{400}{\kilo\metre} という状況と、平均 \SI{150}{\kilo\metre} だけを見て計画することは、運用上まったく違う。}
\label{fig:m6-wakayama}
\end{figure}

\begin{figure}[H]
\centering
\includegraphics[width=0.92\linewidth]{wakayama/M8_landing_ellipse_map.pdf}
\caption[代表地点からの落下分布]{1 地点から上げた場合の落下位置を季節別に重ねたもの。楕円は 50 / 90 / 95 \% の範囲で、各解析時刻ごとに求めた着地点の共分散から描いている。散らばりが 2 変量正規に従うという仮定は残る。灰色が陸、白が海である。}
\label{fig:m8-wakayama}
\end{figure}


飛行時間は約 \num{137} 分、領域平均の飛行距離は年間平均 \SI{158}{\kilo\metre}。
北海道（\SI{143}{\kilo\metre}）よりわずかに長い。

\begin{table}[H]
\centering
\caption{和歌山：飛行距離が最短・最長のビン（領域平均）。
方位は機体が\textbf{進む}向き。}
\begin{tabular}{lrrr}
\toprule
時期 & 平均距離 [\si{\kilo\metre}] & 方位 [度] & p90 [\si{\kilo\metre}] \\
\midrule
8 月 後半 & 49 & 285（西北西） & 83 \\
9 月 前半 & 53 & 44（北東） & 87 \\
8 月 前半 & 54 & 259（西南西） & 87 \\
\midrule
2 月 前半 & 276 & 87（東） & 343 \\
1 月 後半 & 270 & 88（東） & 341 \\
12 月 後半 & 269 & 85（東） & 333 \\
\bottomrule
\end{tabular}
\end{table}

北海道と最も大きく違うのは\textbf{夏の飛行方位}である。

\begin{itemize}
  \item 冬（12--2 月）は東へ \SIrange{251}{276}{\kilo\metre}。
        図 \ref{fig:m8-wakayama} のとおり、和歌山中央から上げると
        平均落下点は御前崎の南の海上になる。
  \item 7 月後半から 8 月は\textbf{西寄り}で、方位は \ang{247}--\ang{285}。
        一方で 7 月前半は東（\ang{96}）、9 月前半は北東（\ang{44}）と、
        半月ごとに向きが入れ替わる。
        対流圏上部の弱い西風と成層圏の東風がほぼ釣り合い、
        どちらが少し勝つかで向きが決まるためである。
        北海道の夏（東南東）とは向きが違い、しかも距離が短い。
  \item 8 月後半が年間で最も扱いやすく、平均 \SI{49}{\kilo\metre}、p90 でも \SI{83}{\kilo\metre} に収まる。
\end{itemize}

ただし夏は、東向きの日と西向きの日が打ち消し合うので、
\textbf{平均落下点は打ち上げ地点のすぐ近くになる}。
8 月後半の領域平均は 1 回あたり \SI{49}{\kilo\metre} 飛ぶのに、
平均落下点は \SI{18}{\kilo\metre} しか動かない。
夏の計画は平均の矢印ではなく、散らばり（図 \ref{fig:m8-wakayama} の楕円）で立てる必要がある。
海に落ちやすい北海道とは、回収計画の立て方が大きく変わる。

\subsection{日変化と鉛直シアー}

\begin{figure}[H]
\centering
\includegraphics[width=0.75\linewidth]{wakayama/F8_diurnal.pdf}
\caption[日変化]{夜（03 JST）と昼（15 JST）の平均風速の差を、季節ごとに高度に対して描いたもの
（\SI{300}{\hecto\pascal} 以下の層のみ、領域平均）。横軸の目盛りに注意してほしい。
差は 3 領域・4 季節のどれでも \SI{1}{\metre\per\second} 未満で、
日々の変動（数十 \si{\metre\per\second}）に比べて小さい。
これが\textbf{時刻を区別せず平均してよい根拠}である。
ビンを時刻で 4 分割すると標本数も 4 分の 1 になるので、
必要のないところで分割しない理由がここにある。}
\label{fig:f8-wakayama}
\end{figure}

\begin{figure}[H]
\centering
\includegraphics[width=1.0\linewidth]{wakayama/F9_shear.pdf}
\caption[鉛直シアー]{隣り合う気圧面のあいだの風速差を高度差で割ったもの。
機体そのものはシアーで壊れないが、大きいところでは
高度の推定が少しずれるだけで水平の予測が大きく動く。
最大はジェットの上下と地表付近に出る。}
\label{fig:f9-wakayama}
\end{figure}


夜と昼の平均風速の差は最大 \SI{0.9}{\metre\per\second} で、
北海道と違って高度 \SI{1}{\kilo\metre} 付近（夜のほうが強い）に山がある。
その上では \SI{0.6}{\metre\per\second} 以下に収まる。

鉛直シアーは北海道より大きく、最大 \SI{9.3}{(\metre\per\second)\per\kilo\metre}
（12 月後半、高度 \SI{16.3}{\kilo\metre}）。
地表付近ではなく\textbf{ジェットの直上}が最大になるのが和歌山の特徴で、
冬の亜熱帯ジェットが強く、その上で急に弱まるためである。
この高度帯は破裂高度のすぐ下にあたるので、
高度の推定が数百メートルずれると水平の予測が大きく動く。

\subsection{既存シミュレータの ISA 仮定}

\begin{figure}[H]
\centering
\includegraphics[width=1.0\linewidth]{wakayama/F10_isa_check.pdf}
\caption[ISA との比較]{気圧から高度への換算。左は JRA-3Q の実際の気圧面高度と、既存シミュレータが使う国際標準大気 ISA。右はその差。ISA は季節も緯度も天気も知らないので、季節によって系統的にずれる。}
\label{fig:f10-wakayama}
\end{figure}


ISA からのずれは最大 \SI{787}{\metre}（夏）、冬は \SI{242}{\metre}。
北海道（夏 \SI{1036}{\metre}）より小さいが、やはり無視できる量ではない。

\section{結果：愛媛（愛南町）}
\label{sec:ehime}

3 つ目の領域は、愛媛県南端の愛南町から上げる場合を想定したものである。
\texttt{sites.json} の「愛媛・高知」には愛南町と周辺の 8 地点が登録されており、
代表地点はそのうち\textbf{南レク御荘公園}（\ang{32.9642}N, \ang{132.5521}E）とした。
和歌山と違って打ち上げ地点が決まっているので、
北海道と同じく地点ごとの表（表\ref{tab:ainan-offshore}）も作れる。

対象は \ang{31.5}--\ang{34.5}N / \ang{132.0}--\ang{136.0}E、
\ang{0.375} 格子で 108 点。豊後水道の東半分から四国の南半分、土佐湾、
紀伊水道の西側までを含む。

\begin{whybox}{領域が打ち上げ地点を中心にしていない理由}
愛南町は、取得済みのデータ範囲（マスター箱、\ang{132}--\ang{150}E）の西端から
\ang{0.55}（約 \SI{51}{\kilo\metre}）しか離れていない。
西へ広げるにはデータの取り直しが要るので、
領域は地点を中心にせず、\textbf{東へ伸ばした}。
この地点から上げた機体は 1 年のほとんどで東へ流されるので、
統計が要るのも東側である。

代償が 2 つある。
1 つは、領域平均が地点より東（室戸岬や紀伊水道寄り）の風を多く含むことである。
ただし後で見るとおり、この領域の中では風がほとんど変わらないので（図\ref{fig:f6-ehime}）、
平均が地点の値から大きく離れることはない。
もう 1 つは、夏に西へ流れる機体がデータの外へ出やすいことで、
第 \ref{sec:fidelity} 章の軌道計算ではこれを数えて除外している
（\ref{sec:fidelity-error} 節）。
\end{whybox}

最寄り格子（\ang{32.78}N, \ang{132.38}E）までは \SI{26.5}{\kilo\metre} ある。
格子の対角のほぼ真ん中に地点が入るためで、\ang{0.375} 格子で起こりうる最大に近い。
それでも下で見るように領域内の風の差は小さいので、この距離が結論を変えることはない。

\subsection{年間を通した姿}

\begin{figure}[H]
\centering
\includegraphics[width=0.95\linewidth]{ehime/F0_sample_counts.pdf}
\caption[標本数]{各ビンに入った解析値の数。前半は必ず 15 日ぶんで 600 個、後半は月の長さで変わり 1 月は 640 個、2 月は閏年が 3 回あるので 532 個になる。この図が合わなければ取得か集計が壊れているので、他のどの図よりも先に確認する。}
\label{fig:f0-ehime}
\end{figure}

\begin{figure}[H]
\centering
\includegraphics[width=1.0\linewidth]{ehime/F1_alt_half_speed.pdf}
\caption[時期×高度の風速と風向]{上：平均風速。下：平均風向（吹いてくる向き）。風向には円環用のカラーマップを使っている。順次カラーマップを使うと \ang{0} と \ang{360} の境目に実在しない断層が現れるためである。灰色は定常度が 0.3 未満で塗りを外した領域で、そのビンではそもそも風向が定まっていない。}
\label{fig:f1-ehime}
\end{figure}

\begin{figure}[H]
\centering
\includegraphics[width=1.0\linewidth]{ehime/F3_profiles_by_month.pdf}
\caption[月別の高度プロファイル]{月ごとの風速プロファイル。線は各ビンの中央値、灰色の帯はその月の 2 ビンを通した p10--p90 の広がり。帯の幅が平均からのずれの大きさを表しており、線だけを見て計画を立ててはいけない。}
\label{fig:f3-ehime}
\end{figure}


愛媛の上空は、\textbf{和歌山とほぼ同じ気候}である。
\SI{200}{\hecto\pascal}（約 \SI{12.2}{\kilo\metre}）の平均風速は
冬 \SI{75.4}{\metre\per\second}（定常度 0.98）、春 \SI{49.6}{\metre\per\second}、
夏 \SI{21.9}{\metre\per\second}（定常度 0.55）、秋 \SI{39.4}{\metre\per\second}。
和歌山の冬 \SI{75.5}{\metre\per\second}、夏 \SI{22.7}{\metre\per\second} とほとんど変わらない。
いちばん強いのは 2 月前半の \SI{78.8}{\metre\per\second} である。

2 つの領域は緯度がほぼ同じ（\ang{33}--\ang{34}N 付近）で、
冬の亜熱帯ジェットの核がどちらの真上も通る。
経度が \ang{2} ほど違っても、ジェットの強さはほとんど変わらない。
北海道（冬 \SI{40.2}{\metre\per\second}）との差は、経度ではなく緯度の差である。

\subsection{成層圏}

\begin{figure}[H]
\centering
\includegraphics[width=1.0\linewidth]{ehime/F2_alt_half_constancy.pdf}
\caption[定常度]{定常度（ベクトル平均風速 ÷ スカラー平均風速）。1 に近いほど風向が安定し、0 に近いほどビンのなかで風向が入れ替わっている。淡い帯が成層圏の風の反転している高度で、季節とともに上下する。}
\label{fig:f2-ehime}
\end{figure}

\begin{figure}[H]
\centering
\includegraphics[width=1.0\linewidth]{ehime/F7-1_wind_roses.pdf}
\caption[風配図（前半）]{代表格子点の風配図、\textbf{月の前半}（1--15 日）。
平均と楕円では見せられない「分布そのもの」を示す。
ローブが 2 つに割れている場面が、平均風がほぼゼロで分散だけが大きくなる時期にあたる。}
\label{fig:f7-ehime-h1}
\end{figure}

\begin{figure}[H]
\centering
\includegraphics[width=1.0\linewidth]{ehime/F7-2_wind_roses.pdf}
\caption[風配図（後半）]{同じ格子点の\textbf{月の後半}（16 日--月末）。
前半（図\ref{fig:f7-ehime-h1}）と並べると、ほとんどの月で\textbf{2 枚はよく似ている}。
唯一 10 月だけ \SI{30}{\hecto\pascal} の最頻セクターが前半 \ang{90}、
後半 \ang{248} と入れ替わるが、どちらも頻度は 2 割前後しかない。
つまり向きが変わったというより、\textbf{その時期はどちらの半月も
風向が定まっていない}のであって、最頻セクターを読む意味が薄い。}
\label{fig:f7-ehime-h2}
\end{figure}

\SI{10}{\hecto\pascal}（約 \SI{31}{\kilo\metre}）は、
夏に \SI{18.5}{\metre\per\second}・定常度 0.99 の\textbf{東風}である。
冬と春の定常度は 0.38、0.36 しかなく、
最も低いのは 3 月後半（0.05）、1 月後半（0.13）、11 月前半（0.13）、
2 月前半（0.17）、3 月前半（0.22）である。
順位の入れ替わりはあるが、和歌山と同じ時期が並んでいる。
冬の \SI{10}{\hecto\pascal} が東風（1 月の領域平均 $u = \SI{-4.2}{\metre\per\second}$）
なのも北海道・和歌山と同じで、アリューシャン高気圧の西縁にあたる。

風配図（図\ref{fig:f7-ehime-h1}、\ref{fig:f7-ehime-h2}）の格子点は、
愛南町ではなく領域の中央（\ang{33.15}N, \ang{134.25}E、室戸岬付近）である。
この図は領域の中央で描く決まりにしてあり、3 領域とも同じである。

\subsection{領域の中で答えが変わるか}

\begin{figure}[H]
\centering
\includegraphics[width=1.0\linewidth]{ehime/F6_spatial_spread.pdf}
\caption[領域内のばらつき]{領域内での平均風速の幅（最大 $-$ 最小）。これが小さい高度・時期では、領域のどこで上げても風はほぼ同じで、代表 1 地点のプロファイルで語ってよい。大きいところではそうではない。}
\label{fig:f6-ehime}
\end{figure}


領域内での平均風速の幅は最大 \SI{13.3}{\metre\per\second} で、
和歌山（\SI{12.0}{\metre\per\second}）とほぼ同じ、
北海道（\SI{31.7}{\metre\per\second}）の半分以下である。
和歌山よりわずかに大きいのは、東西に \ang{4} と和歌山より \ang{1} 広いためである。

したがって\textbf{愛媛の領域も、代表 1 地点でおおむね語ってよい}。
領域平均が東へ偏っていること、最寄り格子まで \SI{26}{\kilo\metre} あることは、
この小ささのおかげで問題にならない。
実際、代表地点の最寄り格子で見ても \SI{200}{\hecto\pascal} は
冬 \SI{75.7}{\metre\per\second}、夏 \SI{21.4}{\metre\per\second} で、領域平均と 1 割も違わない。

\subsection{地図で見る}

\begin{figure}[H]
\centering
\includegraphics[width=1.0\linewidth]{ehime/M1-200_month_grid_200hPa.pdf}
\caption[200 hPa の月別分布図]{\SI{200}{\hecto\pascal}（約 \SI{12}{\kilo\metre}）の平均風を月ごとに並べたもの。塗りが風速、矢印が風向。ジェットが季節とともに動く様子が地図の上で見える。}
\label{fig:m1200-ehime}
\end{figure}

\begin{figure}[H]
\centering
\includegraphics[width=1.0\linewidth]{ehime/M1-030_month_grid_30hPa.pdf}
\caption[30 hPa の月別分布図]{同じく \SI{30}{\hecto\pascal}（約 \SI{24}{\kilo\metre}）。対流圏とは逆の季節変化をしており、夏に東風が強まる。}
\label{fig:m1030-ehime}
\end{figure}

\begin{figure}[H]
\centering
\includegraphics[width=1.0\linewidth]{ehime/M3-200_direction_map_200hPa.pdf}
\caption[風向の分布図]{\SI{200}{\hecto\pascal} の平均風向。円環用のカラーマップを使い、定常度 0.3 未満は灰色にしている。灰色の領域では平均風向に意味がない。}
\label{fig:m3200-ehime}
\end{figure}

\begin{figure}[H]
\centering
\includegraphics[width=1.0\linewidth]{ehime/M2_constancy_map.pdf}
\caption[定常度の分布図]{高度と季節ごとの定常度。濃いところでは風向が安定していて計画が立てやすく、淡いところではビンのなかで風向が入れ替わる。}
\label{fig:m2-ehime}
\end{figure}

\begin{figure}[H]
\centering
\includegraphics[width=0.92\linewidth]{ehime/M4-01_wind_ellipse_map.pdf}
\caption[風のばらつきを地図に]{1 月前半、\SI{200}{\hecto\pascal} の平均風ベクトルと、その先端に描いた $1\sigma$ の散らばり楕円。丸い楕円はどの向きにもばらつくこと、細長い楕円は主に強さがばらつくことを意味する。}
\label{fig:m401-ehime}
\end{figure}


\subsection{落下位置}

\begin{figure}[H]
\centering
\includegraphics[width=1.0\linewidth]{ehime/M5_drift_month_grid.pdf}
\caption[月別の平均落下位置]{各格子点から上げた場合の平均落下位置。\textbf{矢印は実寸}で、矢先がそのまま予測落下点である。色が平均の飛行距離。全パネルが同じ枠と同じ色スケールなので、月どうしを直接比べられる。}
\label{fig:m5-ehime}
\end{figure}

\begin{figure}[H]
\centering
\includegraphics[width=0.92\linewidth]{ehime/M7-01_drift_map.pdf}
\caption[1 月の落下位置]{1 月前半の平均落下位置を拡大したもの。枠は集計対象の領域で、矢先はその外へ出ていることが多い。}
\label{fig:m701-ehime}
\end{figure}

\begin{figure}[H]
\centering
\includegraphics[width=0.92\linewidth]{ehime/M7-07_drift_map.pdf}
\caption[7 月の落下位置]{同じく 7 月前半。1 月と比べると距離も向きも大きく違う。}
\label{fig:m707-ehime}
\end{figure}

\begin{figure}[H]
\centering
\includegraphics[width=1.0\linewidth]{ehime/M6_drift_percentile_map.pdf}
\caption[荒れた日の飛行距離]{飛行距離の 90 パーセンタイル。平均 \SI{150}{\kilo\metre} で p90 が \SI{400}{\kilo\metre} という状況と、平均 \SI{150}{\kilo\metre} だけを見て計画することは、運用上まったく違う。}
\label{fig:m6-ehime}
\end{figure}

\begin{figure}[H]
\centering
\includegraphics[width=0.92\linewidth]{ehime/M8_landing_ellipse_map.pdf}
\caption[代表地点からの落下分布]{1 地点から上げた場合の落下位置を季節別に重ねたもの。楕円は 50 / 90 / 95 \% の範囲で、各解析時刻ごとに求めた着地点の共分散から描いている。散らばりが 2 変量正規に従うという仮定は残る。灰色が陸、白が海である。}
\label{fig:m8-ehime}
\end{figure}


飛行時間は約 \num{137} 分、領域平均の飛行距離は年間平均 \SI{156}{\kilo\metre}。
和歌山（\SI{158}{\kilo\metre}）とほぼ同じである。

\begin{table}[H]
\centering
\caption{愛媛：飛行距離が最短・最長のビン（領域平均）。
方位は機体が\textbf{進む}向き。}
\begin{tabular}{lrrr}
\toprule
時期 & 平均距離 [\si{\kilo\metre}] & 方位 [度] & p90 [\si{\kilo\metre}] \\
\midrule
8 月 後半 & 50 & 281（西） & 85 \\
9 月 前半 & 51 & 32（北北東） & 83 \\
7 月 前半 & 54 & 96（東） & 89 \\
\midrule
2 月 前半 & 275 & 87（東） & 340 \\
1 月 後半 & 270 & 88（東） & 339 \\
12 月 後半 & 266 & 85（東） & 327 \\
\bottomrule
\end{tabular}
\end{table}

この表も和歌山とほとんど同じで、夏の飛行方位も同じように振れる。

\begin{itemize}
  \item 冬（12--2 月）は東へ \SIrange{248}{275}{\kilo\metre}。
  \item 7 月後半から 8 月は\textbf{西寄り}（方位 \ang{252}--\ang{281}）、
        9 月前半は北北東、7 月前半は東と、半月ごとに向きが変わる。
        夏は対流圏上部の弱い西風と成層圏の東風がほぼ釣り合うので、
        どちらが少し勝つかで向きが決まる。
  \item 7--8 月の平均距離は \SIrange{50}{58}{\kilo\metre} だが、
        東西に打ち消し合うので\textbf{平均落下点}は打ち上げ地点のすぐ近くになる。
        \SI{50}{\kilo\metre} 飛ぶ機体が、ある日は東へ、別の日は西へ落ちるということである。
\end{itemize}

最後の点は、平均落下点だけを見ると誤解しやすい。
たとえば 7 月前半の領域平均は、1 回あたり \SI{54}{\kilo\metre} 飛ぶのに、
平均落下点は東へ \SI{12}{\kilo\metre} しか動かない。
夏の計画では、平均の矢印ではなく散らばり（図\ref{fig:m8-ehime} の楕円）を見る必要がある。

\subsection{南レク御荘公園から上げた場合}

代表地点から上げたときの平均落下点を、陸か海か、最寄りの海岸まで何キロかで整理した。
海岸線は北海道の表\ref{tab:voising-offshore} と同じもの
（\texttt{frontend/public/geo/land.json}）を使っている。

\begin{table}[H]
\centering
\caption{南レク御荘公園から上げた場合の平均落下点。
飛行距離は 1 回ごとの距離の平均で、平均落下点までの距離とは違う。}
\label{tab:ainan-offshore}
\begin{tabular}{@{}llrrcr@{}}
\toprule
月 & ビン & {飛行距離 [\si{\kilo\metre}]} & {平均落下点の方位 [度]} & 陸／海 & {岸まで [\si{\kilo\metre}]} \\
\midrule
1 月  & 前半 & 264 & 88  & 海 & 59 \\
1 月  & 後半 & 270 & 88  & 海 & 54 \\
4 月  & 前半 & 177 & 88  & 海 & 32 \\
4 月  & 後半 & 155 & 87  & 海 & 24 \\
7 月  & 前半 & 53  & 95  & \textbf{陸} & 5 \\
7 月  & 後半 & 59  & 257 & 海 & 10 \\
10 月 & 前半 & 93  & 81  & 海 & 25 \\
10 月 & 後半 & 151 & 83  & 海 & 15 \\
\bottomrule
\end{tabular}
\end{table}

\textbf{7 月を除くと、平均落下点はどれも四国の南の海上である。}
東へ \SIrange{85}{265}{\kilo\metre} 流されて、土佐湾の沖から紀伊半島の南西沖に落ちる。
1 月は潮岬の西南西、岸から \SIrange{54}{59}{\kilo\metre} の沖で、
4 つの季節で最も岸から遠い。

\textbf{7 月は向きそのものが半月で入れ替わる。}
7 月前半の平均落下点は東へ \SI{9}{\kilo\metre} しか動かず、陸の上になる。
7 月後半は逆に西南西へ \SI{42}{\kilo\metre} 動き、豊後水道の南部の海上になる。
どちらも平均落下点が打ち上げ地点に近いのは、
上で述べたとおり東向きの日と西向きの日が打ち消し合った結果である。
1 回ごとの飛行距離は \SIrange{53}{59}{\kilo\metre} あり、
p90 では \SIrange{88}{96}{\kilo\metre} になる。

\begin{whybox}{北海道の VOISING と比べると}
VOISING は夏に岸から \SIrange{2}{4}{\kilo\metre}、他の季節は \SIrange{55}{80}{\kilo\metre} 沖に落ちる
（表\ref{tab:voising-offshore}）。
愛南町の冬の沖合いの距離（\SIrange{54}{59}{\kilo\metre}）は北海道と同じくらいだが、
春と秋は \SIrange{15}{32}{\kilo\metre} と北海道より岸に近い。
土佐湾の沖や室戸岬の南に落ちるためである。
夏は北海道が「岸のすぐそばの海」なのに対し、愛南町は
\textbf{陸にも海にも落ちうる}。回収の手配を両方に備える必要がある。

なお、この表は代表 1 地点の値である。
\texttt{out/csv/region\_sites\_ehime.csv} に「愛媛・高知」の 8 地点と
最寄り格子までの距離（\SIrange{12}{26}{\kilo\metre}）を出してあるので、
別の地点で見たい場合は \texttt{csv\_out.REFERENCE\_SITE} を変えて再実行すればよい。
\end{whybox}

\subsection{日変化と鉛直シアー}

\begin{figure}[H]
\centering
\includegraphics[width=0.75\linewidth]{ehime/F8_diurnal.pdf}
\caption[日変化]{夜（03 JST）と昼（15 JST）の平均風速の差を、季節ごとに高度に対して描いたもの
（\SI{300}{\hecto\pascal} 以下の層のみ、領域平均）。横軸の目盛りに注意してほしい。
差は 3 領域・4 季節のどれでも \SI{1}{\metre\per\second} 未満で、
日々の変動（数十 \si{\metre\per\second}）に比べて小さい。
これが\textbf{時刻を区別せず平均してよい根拠}である。
ビンを時刻で 4 分割すると標本数も 4 分の 1 になるので、
必要のないところで分割しない理由がここにある。}
\label{fig:f8-ehime}
\end{figure}

\begin{figure}[H]
\centering
\includegraphics[width=1.0\linewidth]{ehime/F9_shear.pdf}
\caption[鉛直シアー]{隣り合う気圧面のあいだの風速差を高度差で割ったもの。
機体そのものはシアーで壊れないが、大きいところでは
高度の推定が少しずれるだけで水平の予測が大きく動く。
最大はジェットの上下と地表付近に出る。}
\label{fig:f9-ehime}
\end{figure}


夜と昼の平均風速の差は最大 \SI{0.9}{\metre\per\second} で、
和歌山と同じく高度 \SI{1}{\kilo\metre} 付近（夜のほうが強い）に山がある。

鉛直シアーの最大は \SI{9.4}{(\metre\per\second)\per\kilo\metre}
（12 月後半、高度 \SI{16.3}{\kilo\metre}）で、
場所も大きさも和歌山（\SI{9.3}{(\metre\per\second)\per\kilo\metre}）とほぼ同じ。
冬の亜熱帯ジェットの真上で風が急に弱まるところである。

\subsection{既存シミュレータの ISA 仮定}

\begin{figure}[H]
\centering
\includegraphics[width=1.0\linewidth]{ehime/F10_isa_check.pdf}
\caption[ISA との比較]{気圧から高度への換算。左は JRA-3Q の実際の気圧面高度と、既存シミュレータが使う国際標準大気 ISA。右はその差。ISA は季節も緯度も天気も知らないので、季節によって系統的にずれる。}
\label{fig:f10-ehime}
\end{figure}


ISA からのずれは最大 \SI{790}{\metre}（夏）、冬は \SI{269}{\metre} である。
和歌山（夏 \SI{787}{\metre}）とほぼ同じで、北海道（夏 \SI{1036}{\metre}）より小さい。

\section{比較と考察}
\label{sec:compare}

\begin{figure}[H]
\centering
\includegraphics[width=0.92\linewidth]{compare/M0-01_overview_map.pdf}
\caption[広域図（1 月）]{1 月後半、\SI{200}{\hecto\pascal} の平均風と 3 つの対象領域。
矢印の長さは 3 領域で共通の縮尺である。
ジェットの核が和歌山・愛媛の緯度帯にあり、北海道側では弱まっているのが見える。
和歌山と愛媛の領域は \ang{134}--\ang{136}E で重なっている。}
\label{fig:m0-jan}
\end{figure}

\begin{figure}[H]
\centering
\includegraphics[width=0.92\linewidth]{compare/M0-07_overview_map.pdf}
\caption[広域図（7 月）]{同じく 7 月後半。冬とは逆に、北海道側のほうが風が強い。
和歌山・愛媛の上空は風が弱く、矢印は南から南東を向く（北寄りの風）。}
\label{fig:m0-jul}
\end{figure}

\begin{figure}[H]
\centering
\includegraphics[width=\linewidth]{compare/F11_region_comparison.pdf}
\caption[北海道と和歌山の比較]{上 2 段は同じカラースケールでの平均風速、
下段はその差（北海道 $-$ 和歌山）。
赤は北海道のほうが強いビン・高度、青は和歌山のほうが強いビン・高度である。}
\label{fig:f11}
\end{figure}

\begin{figure}[H]
\centering
\includegraphics[width=\linewidth]{compare/F11b_region_comparison.pdf}
\caption[和歌山と愛媛の比較]{同じ図を和歌山と愛媛で描いたもの。
上 2 段はほとんど見分けがつかない。下段の差（和歌山 $-$ 愛媛）は
最大でも \SI{1.7}{\metre\per\second} で、カラースケールの幅が
図\ref{fig:f11} の 20 分の 1 以下であることに注意。}
\label{fig:f11b}
\end{figure}

\subsection{ジェットは季節とともに南北に動く}

図 \ref{fig:f11} の差分パネルに、両者の関係がそのまま出ている。

\begin{itemize}
  \item \textbf{冬（11--3 月）}、高度 \SIrange{8}{16}{\kilo\metre} は濃い青、
        すなわち和歌山のほうが最大 \SI{40}{\metre\per\second} 強い（2 月前半、\SI{200}{\hecto\pascal}）。
  \item \textbf{夏から初秋（7--9 月）}、同じ高度帯が赤に変わり、
        北海道のほうが最大 \SI{21}{\metre\per\second} 強い（9 月後半、\SI{200}{\hecto\pascal}）。
\end{itemize}

これは東アジアジェットの季節移動そのものである。
冬にはジェットの核が \ang{30}--\ang{35}N まで南下して和歌山と愛媛の真上を通り、
夏には \ang{40}N 以北へ退いて北海道の上に来る。

ここから読み取れるのは、\textbf{両者が逆位相で動いている}ということである。
一方が荒れているとき、もう一方は穏やかになる。
年間で最も穏やかな時期は、北海道が 7 月後半、和歌山と愛媛が 8 月後半と、
いずれも夏である。
これは対流圏のジェットが弱まる時期が共通しているためで、
差が出るのはその前後の季節である。

一方、\textbf{和歌山と愛媛はほとんど同じ気候である}（図\ref{fig:f11b}）。
経度で \ang{2} ほど離れていても、同じ緯度帯のジェットの下にあるので、
平均風速の差はどのビン・高度でも \SI{1.7}{\metre\per\second} 以内に収まる。
3 領域の違いを決めているのは経度ではなく緯度である。

\begin{table}[H]
\centering
\caption{主な数値の比較（領域平均、2016--2025 年）。
方位は機体が進む向きで、季節の値は正味の変位から求めた。}
\label{tab:compare}
\small
\begin{tabular}{lrrr}
\toprule
 & 北海道 & 和歌山 & 愛媛 \\
\midrule
格子点数 & 570 & 72 & 108 \\
\SI{200}{\hecto\pascal} 冬の平均風速 [\si{\metre\per\second}] & 40.2 & \textbf{75.5} & \textbf{75.4} \\
\SI{200}{\hecto\pascal} 夏の平均風速 [\si{\metre\per\second}] & \textbf{29.9} & 22.7 & 21.9 \\
\SI{200}{\hecto\pascal} 秋の平均風速 [\si{\metre\per\second}] & \textbf{44.5} & 40.3 & 39.4 \\
\SI{10}{\hecto\pascal} 夏の平均風速 [\si{\metre\per\second}] & 13.6 & \textbf{18.2} & \textbf{18.5} \\
領域内の風速差の最大 [\si{\metre\per\second}] & \textbf{31.7} & 12.0 & 13.3 \\
飛行距離 年間平均 [\si{\kilo\metre}] & 143 & 158 & 156 \\
飛行距離 最短のビン [\si{\kilo\metre}] & 62（7 月後半） & 49（8 月後半） & 50（8 月後半） \\
飛行距離 最長のビン [\si{\kilo\metre}] & 215（12 月前半） & 276（2 月前半） & 275（2 月前半） \\
冬（12--2 月）の飛行方位 & 東（\ang{83}） & 東（\ang{87}） & 東（\ang{86}） \\
7 月後半--8 月の飛行方位 & 東〜東南東（\ang{80}--\ang{123}） & \textbf{西寄り（\ang{247}--\ang{285}）} & \textbf{西寄り（\ang{252}--\ang{281}）} \\
ISA からのずれ 最大 [\si{\metre}] & \textbf{1036} & 787 & 790 \\
\bottomrule
\end{tabular}
\end{table}

\subsection{夏の飛行方位が逆を向く}

回収計画を左右するのは、\textbf{夏の飛行方位}である。

北海道は夏でも東から東南東へ飛ぶ（7 月後半--8 月で \ang{80}--\ang{123}）。
弱いながら対流圏上部に西風が残るためである。
和歌山と愛媛は、7 月後半から 8 月に西寄り（\ang{247}--\ang{285}）へ飛ぶ。
対流圏上部の風が弱まり、成層圏の東風（\SI{10}{\hecto\pascal} で
\SI{18}{\metre\per\second}、定常度 0.99）と釣り合うか、わずかに上回るからである。
7 月前半は東、9 月前半は北東と、半月ごとに向きが入れ替わるのもこの釣り合いのためである。

そのため南の 2 領域では、夏の\textbf{平均落下点は打ち上げ地点のすぐ近くに来る}。
1 回ごとには \SIrange{49}{58}{\kilo\metre} 飛ぶのに、東向きの日と西向きの日が打ち消し合う。
愛南町から上げると、7 月前半の平均落下点は陸、7 月後半は豊後水道の海上になる
（表\ref{tab:ainan-offshore}）。
海上への流出が主な懸念である北海道とは、
回収計画も安全評価も設計が変わる。

\subsection{領域を代表点で語れるかどうかも違う}

領域内の風速差は北海道で最大 \SI{31.7}{\metre\per\second}、
和歌山で \SI{12.0}{\metre\per\second}、愛媛で \SI{13.3}{\metre\per\second}。
和歌山と愛媛は緯度幅が \ang{2.5}--\ang{3} と狭くジェットの南北勾配をまたがないので、1 点で近似してよい。
北海道は \ang{6.5} の緯度幅があり、冬のジェット勾配をまたぐので、
\textbf{道南と道北を同じ数字で語ってはいけない}。

これは「解像度を上げたら何が分かったか」の具体例でもある。
\ang{1.25} 格子では北海道の領域に 54 点しか入らず、
この南北勾配を読み取るのは難しかった。

\subsection{上層の定常度はどの領域でも低い}

破裂高度帯（\SI{10}{\hecto\pascal}、約 \SI{31}{\kilo\metre}）の定常度は、
夏こそ 3 領域とも 0.99 だが、冬と春は 0.36--0.44、秋は 0.31--0.52 に落ちる。
ビン単位では 0.2 を切るところがあり、
北海道は 10 月前半（0.15）と 5 月前半（0.17）、
和歌山と愛媛は 3 月後半（どちらも 0.05）が最も低い。

この時期は\textbf{平均風向に意味がない}ので、
落下位置の平均値も参考程度にとどめるべきである。
散らばりのほうを見ること（本統計では p90 と楕円がそれにあたる）。

\subsection{Open-Meteo の上限による誤差はどれくらいか}
\label{sec:ceiling}

第 \ref{sec:intro-ceiling} 節で触れた「\SI{30}{\hecto\pascal} より上が分からない」という問題は、
本報告の統計を使えば定量化できる。
飛行モデルの層滞在時間（式 \eqref{eq:drift} の $T_k$）で重み付けし、
\SI{23.8}{\kilo\metre} より上の層が水平移動にどれだけ寄与するかを求めた。

\begin{table}[H]
\centering
\caption{\SI{23.8}{\kilo\metre}（Open-Meteo の上限）より上の層が、
水平移動に占める割合。1 つのビンの値ではなく、各月 2 ビンの平均。}
\begin{tabular}{lrrrrrr}
\toprule
 & 1 月 & 4 月 & 6 月 & 7 月 & 8 月 & 11 月 \\
\midrule
北海道 & 4.6 \% & 4.8 \% & 18.4 \% & \textbf{53.5 \%} & 20.5 \% & 7.1 \% \\
和歌山 & 1.1 \% & 3.6 \% & 20.0 \% & \textbf{111 \%} & 101 \% & 1.9 \% \\
愛媛   & 1.0 \% & 3.4 \% & 21.7 \% & \textbf{140 \%} & 86 \% & 2.0 \% \\
\bottomrule
\end{tabular}
\end{table}

冬は 1--5 \% にとどまる。
対流圏のジェットが桁違いに強いので、その上の層が加わっても正味の変位はほとんど変わらない。
\textbf{冬に関しては、既存実装の上限は実害がほとんど無い}。

一方、夏は状況が違う。
和歌山の 7 月に 111 \%、愛媛の 7 月に 140 \% という値が出ているのは、
分母が\textbf{正味の}変位で、分子が上層だけの寄与だからである。
上層の東風が下層の流れを打ち消しているため、正味の変位が小さくなり、比が 1 を超える。
言い換えると、\textbf{夏の和歌山と愛媛では、上層を無視すると向きそのものが逆に出る}。

そして第 \ref{sec:hokkaido} 章から第 \ref{sec:ehime} 章で見たとおり、
飛行距離が年間で最も短いのは 3 領域とも夏である。
打ち上げに最も向いた季節に、既存実装の誤差が最も大きくなる。

\subsection{モンテカルロの σ をどこから取るか}

本統計の副産物として、
シミュレータのモンテカルロ計算に与える風の摂動を実測から決められるようになった。

\begin{itemize}
  \item 高度ごとの風のばらつきは \texttt{u\_var}, \texttt{v\_var}, \texttt{uv\_cov}
        に入っている。等方ガウスではなく、実際の異方性を持った楕円が使える。
  \item 落下位置そのもののばらつきは \texttt{drift\_east\_var} などにある。
        こちらは各解析時刻ごとに飛行を積分した結果なので、
        層間の相関を仮定せずに得られている。
\end{itemize}

現在の実装（\texttt{frontend/src/mc/sampling.ts}）は全層に等方の摂動を与えているが、
実際のばらつきは高度によって桁で違い、向きにも偏りがある。
そこを置き換える材料はこの出力に揃っている。

\section{日々の変動と、シミュレーションの当たりやすさ}
\label{sec:fidelity}

ここまでの章は「いつ・どこで・どの高度の風が、ふつうどちらへどれだけ吹くか」を
まとめてきた。本章はそこから一歩進んで、打ち上げを計画する人が実際に抱える
2 つの問いに答える。

\begin{quote}
\begin{enumerate}
\item \textbf{日ごとの風向・風速の変わり方は季節によって違うのか}。
\item \textbf{シミュレーション通りになりやすいのは、どの季節か}。
\end{enumerate}
\end{quote}

3 領域とも同じ手順で測った。
北海道と南の 2 領域（和歌山・愛媛）では季節の効き方がまるで違うので、
片方だけを見て一般論を言うことができない。
一方、和歌山と愛媛はほとんど同じ結果になった。

1 つ目はデータを見れば答えが出る。一方で 2 つ目は、
「シミュレーション通り」が何を指すのかを決めないと答えようがない。
そこで先に、落下点の予測が外れる原因を整理しておく。

\subsection{予測が外れる原因は 4 つあり、測れるのは 3 つ}
\label{sec:fidelity-scope}

\begin{table}[H]
\centering
\caption{落下点の予測を外す 4 つの原因と、再解析から測れるかどうか}
\label{tab:fidelity-sources}
\begin{tabular}{@{}llc@{}}
\toprule
 & 原因 & 再解析から測れるか \\
\midrule
(1) & シミュレータが風を\textbf{打ち上げ時刻}で固定している & 測れる \\
(2) & シミュレータが風を\textbf{打ち上げ地点}で固定している & 測れる \\
(3) & 打ち上げが予定から N 時間ずれる & 測れる \\
(4) & 与えた\textbf{予報自体}が外れている & 測れない \\
\bottomrule
\end{tabular}
\end{table}

(4) を測るには「過去に出された予報」の保管庫が要る。
本研究が持っているのは「過去に実際どうだったか」の再解析であって、
予報の保管庫ではない。したがって (4) は本章の対象外である。

この点は結論の読み方に直結するので、先に言っておく。
実務では (4) がいちばん大きい誤差である可能性が高い。
本章が出す数字は、\textbf{予報が完璧だったとしてもこれだけは外れる}という下限であって、
実際の外れ幅の見積もりではない。

\subsection{シミュレータが置いている 2 つの仮定}

\texttt{backend/bfps/core/wind.py} の \texttt{WindProfile.wind\_at} は、
緯度・経度・時刻を引数で受け取りながら、そのすべてを無視して、
打ち上げ点で 1 回取得したプロファイルを返す。
呼び出し側はそれを飛行全体（約 \SI{137}{\minute}）に使い回す。

その間に機体は水平に \SIrange{60}{300}{\kilo\metre} 移動し、風そのものも変わる。
つまり機体は、もう自分がいない場所の、もう過ぎた時刻の風で飛ばされている。

\begin{whybox}{これは設計ではなく、風の取り出し方だけの問題である}
\texttt{trajectory.py} の \texttt{horizontal\_derivative} は、
時刻 \texttt{tt} を引数で受け取っている。
RK4 の呼び出し側も各段で \texttt{t}、\texttt{t+dt/2}、\texttt{t+dt} をきちんと渡している。
ところが関数の中身は \texttt{tt} を捨てて、打ち上げ時刻を使う。

位置のほうも同じで、RK4 は緯度経度を正しく更新して渡しているのに、
受け取った \texttt{WindProfile} がそれを見ない。

積分の骨組みは 4 次元の風を扱えるようにできていて、
風を返す側だけがそうなっていない。
直すこと自体は難しくない。難しいのは\textbf{直す価値があるかどうかの判断}で、
そのためには誤差の大きさを知る必要がある。本章はその数字を出す。
\end{whybox}

\subsection{日々の変化は季節でどう違うか}
\label{sec:fidelity-variability}

\subsubsection{風はどれだけ動くか}

まず、時間 $L$ が経ったときに風がどれだけ変わるかを見る。
測るのはベクトルの差 $|\mathbf{V}(t+L) - \mathbf{V}(t)|$ である。

\begin{whybox}{なぜ風速の差ではなくベクトルの差なのか}
風速だけを比べると、強さを保ったまま \ang{90} 向きを変えた風が
「まったく変わっていない」と判定されてしまう。
その風の中で上げたスペースバルーンは、まるきり違うところに落ちる。

向きの変化と強さの変化を一度に扱うには、ベクトルの差をとるしかない。
なお、角度そのものを引き算してはいけない理由は
\ref{sec:circular} 節に述べたとおりで、ここでも角度は一度も引いていない。
\end{whybox}

季節ごとの組は、両端が同じ季節に入るものだけを数えている。
遅れ 10 日では 2 月末に始まった組が 3 月に着いてしまうので、
そのまま数えると、季節差がいちばん出てほしい長い遅れのところで
季節どうしが混ざってしまう。

\begin{figure}[H]
\centering
\includegraphics[width=1.0\linewidth]{fidelity/G1-wakayama_change_curve.pdf}
\caption[風の変化量]{遅れ $L$ を置いたときの風の変化量 $|\mathbf{V}(t+L)-\mathbf{V}(t)|$。実線が中央値、帯が p90 まで。横軸は対数で、6 時間から 20 日まで。組は両端が同じ季節に入るものだけを数えている。\SI{10}{\hecto\pascal} と \SI{30}{\hecto\pascal} の曲線が遅れ 24 時間でいったん下がるのは大気潮汐による見かけの効果である。}
\label{fig:g1-wakayama}
\end{figure}

\begin{figure}[H]
\centering
\includegraphics[width=1.0\linewidth]{fidelity/G1-hokkaido_change_curve.pdf}
\caption[風の変化量]{遅れ $L$ を置いたときの風の変化量 $|\mathbf{V}(t+L)-\mathbf{V}(t)|$。実線が中央値、帯が p90 まで。横軸は対数で、6 時間から 20 日まで。組は両端が同じ季節に入るものだけを数えている。\SI{10}{\hecto\pascal} と \SI{30}{\hecto\pascal} の曲線が遅れ 24 時間でいったん下がるのは大気潮汐による見かけの効果である。}
\label{fig:g1-hokkaido}
\end{figure}

\begin{figure}[H]
\centering
\includegraphics[width=1.0\linewidth]{fidelity/G1-ehime_change_curve.pdf}
\caption[風の変化量]{遅れ $L$ を置いたときの風の変化量 $|\mathbf{V}(t+L)-\mathbf{V}(t)|$。実線が中央値、帯が p90 まで。横軸は対数で、6 時間から 20 日まで。組は両端が同じ季節に入るものだけを数えている。\SI{10}{\hecto\pascal} と \SI{30}{\hecto\pascal} の曲線が遅れ 24 時間でいったん下がるのは大気潮汐による見かけの効果である。}
\label{fig:g1-ehime}
\end{figure}


和歌山の \SI{200}{\hecto\pascal} では、遅れ 1 日で春 \SI{16.8}{\metre\per\second}、
夏 \SI{11.7}{\metre\per\second} と開く。
\textbf{春がもっとも荒く、夏がもっとも穏やか}である。
北海道でも季節の順序は同じだが、季節間の開きが小さい。
愛媛の曲線は和歌山とほとんど重なる。

\SI{10}{\hecto\pascal} と \SI{30}{\hecto\pascal} の曲線は、
遅れ 12 時間から 24 時間にかけていったん下がる。
これは大気潮汐で、1 日周期の成分が 24 時間後にはほぼ元に戻るためである。
遅れを 1 日の整数倍にとると変化が小さく見える、という見かけの効果なので、
この 2 つの高度では 6 時間と 12 時間の値を読むときに注意がいる。

\subsubsection{風の偏差が続く時間}

次に、風の偏差がどれだけ持続するかを見る。
各季節・各高度で、偏差の自己相関が $1/e$ を切る遅れを求めた。
偏差はその季節自身の平均からとっている。
そうしないと、強い西風が定常に吹いているだけの季節が、
「持続性が高い」と高く出てしまう。

\begin{figure}[H]
\centering
\includegraphics[width=0.95\linewidth]{fidelity/G2-wakayama_decorrelation.pdf}
\caption[風の偏差が続く時間]{風の偏差の自己相関が $1/e$ を切るまでの日数。偏差はその季節自身の平均からとっているので、定常な風が強いだけの季節が高く出ることはない。棒がなく「20 日超」とある組は、20 日の遅れでもまだ相関が $1/e$ を上回っていた場合で、偏差が続く時間が短いのではなく長いことを意味する。}
\label{fig:g2-wakayama}
\end{figure}

\begin{figure}[H]
\centering
\includegraphics[width=0.95\linewidth]{fidelity/G2-hokkaido_decorrelation.pdf}
\caption[風の偏差が続く時間]{風の偏差の自己相関が $1/e$ を切るまでの日数。偏差はその季節自身の平均からとっているので、定常な風が強いだけの季節が高く出ることはない。棒がなく「20 日超」とある組は、20 日の遅れでもまだ相関が $1/e$ を上回っていた場合で、偏差が続く時間が短いのではなく長いことを意味する。}
\label{fig:g2-hokkaido}
\end{figure}

\begin{figure}[H]
\centering
\includegraphics[width=0.95\linewidth]{fidelity/G2-ehime_decorrelation.pdf}
\caption[風の偏差が続く時間]{風の偏差の自己相関が $1/e$ を切るまでの日数。偏差はその季節自身の平均からとっているので、定常な風が強いだけの季節が高く出ることはない。棒がなく「20 日超」とある組は、20 日の遅れでもまだ相関が $1/e$ を上回っていた場合で、偏差が続く時間が短いのではなく長いことを意味する。}
\label{fig:g2-ehime}
\end{figure}


\textbf{ここで北海道と南の 2 領域が大きく分かれる。}

\begin{table}[H]
\centering
\caption{風の偏差が $1/e$ を切るまでの日数。和歌山と愛媛は季節で 8 倍動くのに、
北海道はほとんど動かない。}
\label{tab:decorrelation-both}
\small
\setlength{\tabcolsep}{3pt}
\sisetup{table-format=2.1}
\begin{tabular}{@{}lSSSSSSSSSSSS@{}}
\toprule
& \multicolumn{4}{c}{北海道} & \multicolumn{4}{c}{和歌山} & \multicolumn{4}{c}{愛媛} \\
\cmidrule(lr){2-5}\cmidrule(lr){6-9}\cmidrule(lr){10-13}
気圧面 & {冬} & {春} & {夏} & {秋} & {冬} & {春} & {夏} & {秋} & {冬} & {春} & {夏} & {秋} \\
\midrule
\SI{10}{\hecto\pascal}  & 4.9 & 7.6 & 14.0 & 8.3  & 4.2 & 6.5 & 17.0 & 8.6  & 4.2 & 6.4 & 17.2 & 8.9 \\
\SI{30}{\hecto\pascal}  & 4.7 & 8.4 & 10.5 & 8.7  & 4.7 & 9.2 & 18.5 & 14.9 & 4.6 & 9.4 & 19.0 & 19.3 \\
\SI{100}{\hecto\pascal} & 3.3 & 2.5 & 3.5  & 4.1  & 2.4 & 3.8 & 15.4 & {20 超} & 2.3 & 4.1 & 15.5 & {20 超} \\
\SI{200}{\hecto\pascal} & 1.8 & 1.2 & 1.8  & 1.7  & 2.0 & 1.7 & 10.7 & 15.3 & 1.9 & 1.9 & 10.8 & 16.3 \\
\SI{500}{\hecto\pascal} & 1.0 & 0.9 & 1.8  & 1.1  & 1.0 & 1.1 & 4.9  & 2.5  & 1.0 & 1.1 & 5.6  & 2.8 \\
\SI{850}{\hecto\pascal} & 0.8 & 0.8 & 1.2  & 0.8  & 0.7 & 0.7 & 1.8  & 0.9  & 0.7 & 0.8 & 2.0  & 0.9 \\
\bottomrule
\end{tabular}
\end{table}

和歌山の \SI{200}{\hecto\pascal} は冬 \num{2.0} 日に対し秋 \num{15.3} 日で、
\SI{100}{\hecto\pascal} の秋にいたっては 20 日を見ても $1/e$ を切らない。
つまり\textbf{夏と秋の風はいったん決まると長く続き、冬と春の風は 1 日で入れ替わる}。
愛媛も冬 \num{1.9} 日、秋 \num{16.3} 日で、和歌山と同じ形である。

ところが北海道の \SI{200}{\hecto\pascal} は
冬 \num{1.8}、春 \num{1.2}、夏 \num{1.8}、秋 \num{1.7} 日と、
\textbf{どの季節でも 2 日を切る}。
対流圏上部については、北海道の風は季節を問わず 1--2 日で入れ替わる。

差が出るのは成層圏だけで、\SI{10}{\hecto\pascal} と \SI{30}{\hecto\pascal} の夏は
北海道でも \numrange{10}{14} 日持続する。
ただし和歌山・愛媛の \numrange{17}{19} 日よりは短い。

\begin{whybox}{なぜ北海道だけ季節差が消えるのか}
和歌山と愛媛は亜熱帯ジェットの南縁にあたり、夏にはジェットが北へ抜けて
穏やかで持続的な流れだけが残る。

北海道は寒帯前線ジェットの通り道で、夏でも低気圧が次々に通過する。
\ref{sec:hokkaido} 章で \SI{200}{\hecto\pascal} の平均風速が
夏でも \SI{29.9}{\metre\per\second} あったのはそのためで、
和歌山の夏（\SI{22.7}{\metre\per\second}）より強い。
流れが止まらないので、偏差も長く続かない。

この違いは第 \ref{sec:validation} 章の結果とも整合する。
北海道で気候値の効きが弱い（誤差削減 17 \%、和歌山は 45 \%、愛媛は 47 \%）のは、
季節が決める分が小さく日々の天気が決める分が大きいからで、
偏差が続く時間の短さも、同じ事情から来ている。
\end{whybox}

\subsubsection{ところが風向は逆の順序になる}

ここまでだけを読むと、夏は予測しやすい季節に見える。
しかし日平均風向の前日差をとると、違う結果になる。

\begin{figure}[H]
\centering
\includegraphics[width=0.95\linewidth]{fidelity/G3-wakayama_direction_change.pdf}
\caption[風向の日々の振れ]{日平均風向の前日差（符号なし）。成分に直してから差をとり、$[-180, 180]$ 度に畳んでいるので、\ang{350} と \ang{10} の差は \ang{20} になる。図\ref{fig:g2-wakayama} と季節の順序が高度によって逆転する点に注意。}
\label{fig:g3-wakayama}
\end{figure}

\begin{figure}[H]
\centering
\includegraphics[width=0.95\linewidth]{fidelity/G3-hokkaido_direction_change.pdf}
\caption[風向の日々の振れ]{日平均風向の前日差（符号なし）。成分に直してから差をとり、$[-180, 180]$ 度に畳んでいるので、\ang{350} と \ang{10} の差は \ang{20} になる。図\ref{fig:g2-hokkaido} と季節の順序が高度によって逆転する点に注意。}
\label{fig:g3-hokkaido}
\end{figure}

\begin{figure}[H]
\centering
\includegraphics[width=0.95\linewidth]{fidelity/G3-ehime_direction_change.pdf}
\caption[風向の日々の振れ]{日平均風向の前日差（符号なし）。成分に直してから差をとり、$[-180, 180]$ 度に畳んでいるので、\ang{350} と \ang{10} の差は \ang{20} になる。図\ref{fig:g2-ehime} と季節の順序が高度によって逆転する点に注意。}
\label{fig:g3-ehime}
\end{figure}


\begin{table}[H]
\centering
\caption{日平均風向の前日差（中央値、度）。高度で季節の順序が入れ替わる。}
\label{tab:direction-change}
\small
\setlength{\tabcolsep}{3pt}
\sisetup{table-format=2.1}
\begin{tabular}{@{}lSSSSSSSSSSSS@{}}
\toprule
& \multicolumn{4}{c}{北海道} & \multicolumn{4}{c}{和歌山} & \multicolumn{4}{c}{愛媛} \\
\cmidrule(lr){2-5}\cmidrule(lr){6-9}\cmidrule(lr){10-13}
気圧面 & {冬} & {春} & {夏} & {秋} & {冬} & {春} & {夏} & {秋} & {冬} & {春} & {夏} & {秋} \\
\midrule
\SI{10}{\hecto\pascal}  & 12.8 & 11.7 & 3.0  & 15.5 & 14.0 & 14.7 & 2.5  & 13.1 & 14.2 & 14.8 & 2.5  & 12.9 \\
\SI{30}{\hecto\pascal}  & 10.2 & 11.4 & 4.7  & 11.7 & 11.2 & 12.0 & 2.8  & 9.7  & 11.1 & 12.0 & 2.7  & 9.2 \\
\SI{100}{\hecto\pascal} & 6.3  & 9.9  & 14.7 & 8.2  & 4.4  & 7.4  & 16.4 & 8.1  & 4.2  & 7.2  & 16.0 & 8.2 \\
\SI{200}{\hecto\pascal} & 9.0  & 16.5 & 18.6 & 12.1 & 5.4  & 10.1 & 21.9 & 10.3 & 5.3  & 9.8  & 22.2 & 10.2 \\
\SI{500}{\hecto\pascal} & 20.2 & 26.1 & 19.4 & 18.0 & 9.0  & 13.9 & 16.3 & 12.6 & 8.8  & 13.3 & 16.1 & 12.3 \\
\SI{850}{\hecto\pascal} & 23.2 & 36.1 & 33.3 & 32.6 & 24.2 & 43.4 & 27.3 & 46.3 & 25.3 & 43.8 & 25.3 & 45.4 \\
\bottomrule
\end{tabular}
\end{table}

\SI{10}{\hecto\pascal} と \SI{30}{\hecto\pascal} では夏が 3 領域とも
\ang{2.5} から \ang{4.7} で、他の季節の 3--5 分の 1 である。
夏の成層圏東風はほとんど向きを変えない。

ところが \SI{100}{\hecto\pascal} から \SI{200}{\hecto\pascal} では順序が反転し、
和歌山では夏が \ang{16.4} から \ang{21.9}、冬が \ang{4.4} から \ang{5.4} になる。
愛媛も同じ値に並ぶ。
北海道でも向きは同じだが、冬が \ang{6.3} から \ang{9.0} と和歌山ほど静かではない。

\SI{500}{\hecto\pascal} では北と南の性格がはっきり分かれる。
和歌山と愛媛の冬は \ang{9} 前後で落ち着いているのに、北海道の冬は \ang{20.2} ある。
中層でも北海道は年じゅう向きが振れている。

\begin{whybox}{持続性が高いのに向きが振れるのはなぜか}
この 2 つは矛盾しない。自己相関が測っているのは\textbf{偏差}の持続性で、
風向の前日差が測っているのは\textbf{風向そのもの}の動きである。

夏の \SI{200}{\hecto\pascal} は平均風が弱い。
弱い風に同じ大きさの偏差が乗ると、向きは大きく回る。
冬は平均風が \SI{70}{\metre\per\second} 前後の強いジェットなので、
同じ偏差では向きがほとんど動かない。

落下点を決めるのは向きと強さの両方なので、
運用上は\textbf{風向の表のほうが直接役に立つ}。
\end{whybox}

\SI{850}{\hecto\pascal} ではまた別の順序になり、南の 2 領域では春（\ang{43}--\ang{44}）と秋（\ang{45}--\ang{46}）が大きい。
地表付近は季節の移り変わりそのものが荒いということで、
これも上空とは別の理由による。

\subsubsection{風の分散はどこにあるか}

最後に、風の分散を 3 つに分けた。

\begin{figure}[H]
\centering
\includegraphics[width=0.85\linewidth]{fidelity/G6-wakayama_variance_split.pdf}
\caption[風の分散の内訳]{風の分散を、季節が説明する分・季節内の日々の変動・1 日の中の変動に分けたもの。時期別の気候値で説明できるのは左の帯だけで、右の 2 つはその日の天気で決まる。}
\label{fig:g6-wakayama}
\end{figure}

\begin{figure}[H]
\centering
\includegraphics[width=0.85\linewidth]{fidelity/G6-hokkaido_variance_split.pdf}
\caption[風の分散の内訳]{風の分散を、季節が説明する分・季節内の日々の変動・1 日の中の変動に分けたもの。時期別の気候値で説明できるのは左の帯だけで、右の 2 つはその日の天気で決まる。}
\label{fig:g6-hokkaido}
\end{figure}

\begin{figure}[H]
\centering
\includegraphics[width=0.85\linewidth]{fidelity/G6-ehime_variance_split.pdf}
\caption[風の分散の内訳]{風の分散を、季節が説明する分・季節内の日々の変動・1 日の中の変動に分けたもの。時期別の気候値で説明できるのは左の帯だけで、右の 2 つはその日の天気で決まる。}
\label{fig:g6-ehime}
\end{figure}


\begin{table}[H]
\centering
\caption{風の分散の内訳（\si{\metre\squared\per\second\squared}、括弧内は全体に占める割合）。}
\label{tab:variance-split}
\begin{tabular}{@{}lrrr@{}}
\toprule
 & 北海道 & 和歌山 & 愛媛 \\
\midrule
全体 & 115.1 & 169.0 & 169.7 \\
季節が説明する分 & 20.9（18 \%） & 76.7（45 \%） & 79.2（47 \%） \\
季節内の日々の天気 & 81.5（71 \%） & 81.7（48 \%） & 80.4（47 \%） \\
1 日の中の変化 & 12.7（11 \%） & 10.6（6 \%） & 10.2（6 \%） \\
\bottomrule
\end{tabular}
\end{table}

和歌山と愛媛では、\textbf{気候値で説明できるのは風の変動のおよそ半分まで}で、
残りの半分はその日の天気が決めている。
北海道では季節の分が 18 \% しかなく、7 割を日々の天気が決めている。
これは \ref{sec:validation} 章の結論と同じことを別の角度から言っている。
半月別の表は計画を立てるための道具であって、予報ではない。

\subsection{打ち上げがずれたら落下点はどれだけ動くか}
\label{sec:fidelity-slip}

準備が遅れて N 時間後に上げ直したとき、落下点はどこへ動くか。
これはシミュレータの誤差ではない。
シミュレータはどちらの打ち上げについても正しく答えたはずである。
測っているのは\textbf{予定がずれることのリスク}で、
打ち上げ計画にどれだけ余裕が要るかを決める数字である。

\begin{figure}[H]
\centering
\includegraphics[width=0.92\linewidth]{fidelity/G4-wakayama_slip_curve.pdf}
\caption[打ち上げが遅れたときの落下点の移動]{打ち上げが遅れたとき、予測落下点がどれだけ動くか。実線が中央値、帯が p90 まで。これはシミュレータの誤差ではなく、予定がずれることのリスクである。}
\label{fig:g4-wakayama}
\end{figure}

\begin{figure}[H]
\centering
\includegraphics[width=0.92\linewidth]{fidelity/G4-hokkaido_slip_curve.pdf}
\caption[打ち上げが遅れたときの落下点の移動]{打ち上げが遅れたとき、予測落下点がどれだけ動くか。実線が中央値、帯が p90 まで。これはシミュレータの誤差ではなく、予定がずれることのリスクである。}
\label{fig:g4-hokkaido}
\end{figure}

\begin{figure}[H]
\centering
\includegraphics[width=0.92\linewidth]{fidelity/G4-ehime_slip_curve.pdf}
\caption[打ち上げが遅れたときの落下点の移動]{打ち上げが遅れたとき、予測落下点がどれだけ動くか。実線が中央値、帯が p90 まで。これはシミュレータの誤差ではなく、予定がずれることのリスクである。}
\label{fig:g4-ehime}
\end{figure}


\begin{table}[H]
\centering
\caption{打ち上げが 1 日遅れたときの落下点の移動（\si{\kilo\metre}）。
北海道はどの季節でも和歌山・愛媛より大きい。}
\label{tab:slip}
\sisetup{table-format=3.1}
\begin{tabular}{@{}lSSSSSS@{}}
\toprule
& \multicolumn{2}{c}{北海道} & \multicolumn{2}{c}{和歌山} & \multicolumn{2}{c}{愛媛} \\
\cmidrule(lr){2-3}\cmidrule(lr){4-5}\cmidrule(lr){6-7}
季節 & {中央値} & {p90} & {中央値} & {p90} & {中央値} & {p90} \\
\midrule
冬 & 55.5 & 108.9 & 47.3 & 89.6 & 45.4 & 86.2 \\
春 & 63.2 & 122.5 & 49.8 & 96.7 & 48.1 & 93.7 \\
夏 & 43.0 & 88.6  & 29.7 & 65.5 & 29.1 & 63.4 \\
秋 & 55.9 & 109.6 & 36.4 & 74.5 & 35.1 & 71.9 \\
\bottomrule
\end{tabular}
\end{table}

1 日の遅れで、和歌山は春 \SI{49.8}{\kilo\metre}、冬 \SI{47.3}{\kilo\metre}、
夏はその 6 割の \SI{29.7}{\kilo\metre}。
\textbf{春がもっとも不安定で、夏がもっとも安定}という、前節と同じ順序である。
愛媛はどの季節も和歌山より \SIrange{1}{2}{\kilo\metre} 小さいだけで、同じ順序に並ぶ。

北海道も順序は同じだが、\textbf{どの季節でも和歌山より大きい}。
いちばん安定した夏でも \SI{43.0}{\kilo\metre} 動き、
これは和歌山のいちばん悪い春（\SI{49.8}{\kilo\metre}）に近い。
北海道では\textbf{いちばん安定した季節を選んでも、和歌山のいちばん不安定な季節と同じくらい落下点が動く}。
偏差が続く時間の短さ（表 \ref{tab:decorrelation-both}）がそのまま表れている。

\subsection{シミュレータ自身の近似誤差}
\label{sec:fidelity-error}

\subsubsection{格子の中を実際に飛ばす}

(1) と (2) の大きさを測るには、固定しない場合の答えが要る。
そこで、機体を格子の中で実際に飛ばす積分器を書いた。
水平は緯度経度で双線形、時間は 6 時間解析の間を線形に内挿し、
移動しながら、そのとき自分がいる場所・時刻の風を読む。

ここから 4 通りを走らせて比べる。

\begin{table}[H]
\centering
\caption{比べた 4 通り。差が測りたいものである。}
\label{tab:variants}
\begin{tabular}{@{}llll@{}}
\toprule
呼び方 & 風を読む時刻 & 風を読む場所 & 対応するもの \\
\midrule
基準 & 飛行中の時刻 & 機体の位置 & 再解析で言える最良の答え \\
時刻のみ固定 & 打ち上げ時刻 & 機体の位置 & (1) だけの影響 \\
地点のみ固定 & 飛行中の時刻 & 打ち上げ点 & (2) だけの影響 \\
両方固定 & 打ち上げ時刻 & 打ち上げ点 & \textbf{現行のシミュレータ} \\
\bottomrule
\end{tabular}
\end{table}

\begin{whybox}{「格子間を内挿しない」という方針に反しないのか}
本パイプラインは格子間の内挿を避けてきた。
2 つの p90 を内挿しても、得られる数は何の p90 でもないからである。

ここで内挿しているのは統計量ではなく\textbf{風の場そのもの}で、
これは格子が標本化している連続な物理場である。
スペースバルーンは格子点の上だけを飛ぶわけではないので、
内挿しないほうがかえって誤差になる。

区別を保つため、この双線形内挿は \texttt{stats/trajectory.py} の中だけに置き、
統計を扱うモジュールには一切持ち込んでいない。
\end{whybox}

鉛直方向は、既存の \texttt{drift.layer\_seconds} と同じ層モデルをそのまま使っている。
気圧面のあいだを内挿すればもう少し忠実にはなるが、
そうすると\texttt{両方固定}と既存の閉形式との差に鉛直の精度差が混ざり、
測ろうとしている「固定したことによる誤差」が分離できなくなる。

高度プロファイルは領域で 1 本に統一した。
4 通りがまったく同じ $z(t)$ を飛ばないと、
差の一部が飛行時間の違いになってしまうからである。
実測すると、領域内で気圧面の高度がばらつく幅は最大 \SI{203}{\metre}、
上昇する \SI{30}{\kilo\metre} の 0.68\,\% であった。

\subsubsection{積分器が正しいことをどう確かめたか}

新しい積分器が出す数字は、積分器が正しいときだけ意味を持つ。
そこで答えが独立にわかる入力で固定した。

もっとも強い確認は、\textbf{両方固定}が既存の閉形式
\texttt{drift.displacements} を再現することである。
両方固定すれば積分は既存の総和 $\sum_k u_k T_k$ に厳密に一致するはずで、
実際、東西方向だけの風（緯度が動かないので地図因子が一定になる）で
機械精度まで一致した。一様な風ではなく、高度・場所・時刻すべてに依存する
風で一致するので、層時間・層の番号付け・固定の切り替えのすべてが検証される。

他に、無風で変位ゼロ、定常一様風で球面上の解析解（航程線の積分）を
相対誤差 $10^{-7}$ で再現、時間だけ変わる風では時刻固定にだけ誤差が出て
地点固定には出ない（空間だけ変わる風ではその逆）、
刻み幅 \SIrange{30}{120}{\second} で着地点が \SI{1}{\kilo\metre} 以内に一致、
といった確認を行った。

\subsubsection{結果}

和歌山は \num{1051992} 回（10 年 $\times$ 6 時間ごと $\times$ 72 格子点）、
北海道は \num{8328270} 回（同じく 570 格子点）、
愛媛は \num{1577988} 回（同じく 108 格子点）の打ち上げを 4 通りずつ飛ばした。

\begin{whybox}{北海道では 7 \% の軌道がデータの外へ出る}
和歌山ではマスター箱の外へ出たものが 0 件だった。
北海道では \num{579271} 件、\SI{6.96}{\percent} が出る。

原因は箱の形である。マスター箱の東端は \ang{150}E で、
北海道領域の東端は \ang{149}E。東へ向かう余裕は 1 度、
\SI{81}{\kilo\metre} しかない（北へは \ang{1.5}、\SI{166}{\kilo\metre} ある）。
冬に \SI{200}{\kilo\metre} 東へ飛ぶ軌道はそこを越えてしまう。

脱出はほぼ東向きである。2020 年 1 月に出た軌道の 82 \%、
7 月では 100 \% が、領域の東端から 1 度以内で上げたものだった。
年間を通すと打ち上げ経度は \ang{138.0} から \ang{148.9} に及ぶので、
まれには西寄りの地点からも出る。それでも主因は東端の余裕不足である。

出た軌道は集計から外してある。
ただし\textbf{これは無作為な脱落ではない}。
速い風に乗って東へ走った軌道だけが抜けるので、
外した集団は飛行距離も誤差も平均より大きい。
10 年ぶんで飛行距離 \SI{206}{\kilo\metre} 対 \SI{132}{\kilo\metre}、
両方固定の誤差 \SI{4.87}{\kilo\metre} 対 \SI{4.11}{\kilo\metre} である。

したがって以下の北海道の誤差は\textbf{わずかに小さめに出ている}。
脱落率と差から見積もると、両方固定の中央値で \SI{1.3}{\percent}、
地点のみ固定で \SI{2.2}{\percent}、時刻のみ固定で \SI{1.1}{\percent}。
どの結論も動かない大きさである。
実装は \texttt{fidelity\_out.exclusion\_bias} が毎回この量を計算して記録する。
箱を広げれば消える制約だが、それには再取得が要る。
\end{whybox}

\begin{whybox}{愛媛では 6 \% の軌道が西端から出る}
愛媛でも \num{93806} 件、\SI{5.94}{\percent} がマスター箱の外へ出た。
こちらは北海道と逆向きで、\textbf{西端}（\ang{132}E）から出ている。
愛媛の領域は箱の西端に接しているので（第 \ref{sec:ehime} 章）、
\ang{132.0}--\ang{133.5}E から上げて夏に西へ流れた軌道が、すぐに外へ出る。

北海道と違い、外れた集団は飛行距離が\textbf{短い}
（\SI{95}{\kilo\metre} 対 \SI{152}{\kilo\metre}）。
夏の弱い風に乗った飛行だからで、固定による誤差も平均より小さい。
したがって以下の愛媛の誤差は\textbf{わずかに大きめに出ている}。
大きさは両方固定の中央値で \SI{1.1}{\percent}、
地点のみ固定で \SI{1.8}{\percent}、時刻のみ固定で \SI{0.8}{\percent} で、
やはり結論は動かない。
\end{whybox}

\begin{figure}[H]
\centering
\includegraphics[width=1.0\linewidth]{fidelity/G5-wakayama_error_breakdown.pdf}
\caption[風を固定したことによる誤差]{基準軌道の着地点からのずれ。棒が中央値、細線が p90 まで。左は \si{\kilo\metre}、右は飛行距離に対する割合。冬は飛行距離が夏の 4 倍あるので、2 つのパネルは季節を逆に並べる。どちらか一方にまとめると、測定ではなく判断になってしまう。}
\label{fig:g5-wakayama}
\end{figure}

\begin{figure}[H]
\centering
\includegraphics[width=1.0\linewidth]{fidelity/G5-hokkaido_error_breakdown.pdf}
\caption[風を固定したことによる誤差]{基準軌道の着地点からのずれ。棒が中央値、細線が p90 まで。左は \si{\kilo\metre}、右は飛行距離に対する割合。冬は飛行距離が夏の 4 倍あるので、2 つのパネルは季節を逆に並べる。どちらか一方にまとめると、測定ではなく判断になってしまう。}
\label{fig:g5-hokkaido}
\end{figure}

\begin{figure}[H]
\centering
\includegraphics[width=1.0\linewidth]{fidelity/G5-ehime_error_breakdown.pdf}
\caption[風を固定したことによる誤差]{基準軌道の着地点からのずれ。棒が中央値、細線が p90 まで。左は \si{\kilo\metre}、右は飛行距離に対する割合。冬は飛行距離が夏の 4 倍あるので、2 つのパネルは季節を逆に並べる。どちらか一方にまとめると、測定ではなく判断になってしまう。}
\label{fig:g5-ehime}
\end{figure}


\begin{table}[H]
\centering
\caption{基準軌道からのずれ（中央値、\si{\kilo\metre}）。
「現行」が実際に動いているシミュレータにあたる。}
\label{tab:fidelity-error}
\sisetup{table-format=1.1}
\begin{tabular}{@{}lSSSS[table-format=3.0]S@{}}
\toprule
季節 & {時刻のみ} & {地点のみ} & {両方（現行）} & {飛行距離 [\si{\kilo\metre}]} & {現行の割合 [\%]} \\
\midrule
\multicolumn{6}{@{}l}{\textbf{北海道}} \\
冬 & 3.6 & 5.8 & 5.1 & 176 & 3.1 \\
春 & 4.1 & 4.7 & 4.2 & 134 & 3.4 \\
夏 & 3.0 & 2.7 & 3.2 & 73  & 4.6 \\
秋 & 3.7 & 5.3 & 4.5 & 164 & 2.9 \\
\addlinespace
\multicolumn{6}{@{}l}{\textbf{和歌山}} \\
冬 & 3.3 & 6.1 & 5.3 & 267 & 2.0 \\
春 & 3.5 & 4.1 & 3.9 & 165 & 2.5 \\
夏 & 2.2 & 2.0 & 2.8 & 63  & 4.6 \\
秋 & 2.7 & 3.2 & 3.5 & 125 & 3.0 \\
\addlinespace
\multicolumn{6}{@{}l}{\textbf{愛媛}} \\
冬 & 3.2 & 5.9 & 5.2 & 267 & 2.0 \\
春 & 3.4 & 4.0 & 3.9 & 165 & 2.5 \\
夏 & 2.2 & 1.9 & 2.8 & 63  & 4.6 \\
秋 & 2.6 & 3.0 & 3.4 & 124 & 3.0 \\
\bottomrule
\end{tabular}
\end{table}

まず絶対値でみると、誤差は 3 領域・4 季節のどれでも
\SIrange{2.8}{5.3}{\kilo\metre} に収まり、飛行距離の \SIrange{2}{5}{\percent} である。
前節の「1 日ずれたら \SIrange{29}{63}{\kilo\metre} 動く」に比べると 1 桁小さい。

\begin{whybox}{風を固定する近似は、直しても精度はほとんど上がらない}
現行シミュレータの近似誤差は、打ち上げが半日ずれたときの移動の
1 割前後にすぎない。
測れなかった予報誤差 (4) はさらに大きいと考えられる。

風を 4 次元で引き直す改修は \texttt{wind.py} だけの話なので難しくはないが、
得られるのは中央値で数キロメートルである。
落下点の精度を上げたいなら、
もっと大きい誤差源から直すほうが効果が大きい。

\textbf{この結論は 3 領域で共通である。} 気候がまるで違う北海道と南の 2 領域で同じ
\SIrange{2}{5}{\percent} に収まったことは、この数字が土地の事情ではなく
飛行時間（約 \SI{137}{\minute}）と風の時間・空間スケールの比で決まっていることを示唆する。
\end{whybox}

\subsubsection{冬には 2 つの固定が打ち消し合う}

表をよく見ると、和歌山の冬だけ\textbf{両方固定（\SI{5.3}{\kilo\metre}）が
地点のみ固定（\SI{6.1}{\kilo\metre}）より小さい}。
誤差を 1 つ足したのに、誤差が減っている。

愛媛の冬も \SI{5.2}{\kilo\metre} 対 \SI{5.9}{\kilo\metre} で同じである。
北海道でも同じことが起きている。冬は \SI{5.1}{\kilo\metre} 対 \SI{5.8}{\kilo\metre}、
秋は \SI{4.5}{\kilo\metre} 対 \SI{5.3}{\kilo\metre} で、
どちらも両方固定のほうが小さい。

これは偶然ではなく、冬の流れの性質から説明がつく。
冬は強い西風で、気圧の谷や尾根も西から東へ動いていく。
機体も東へ運ばれる。

地点を固定すると、機体が着くはずの東側の風ではなく、打ち上げ点の風を使う。
時刻を固定すると、到着時刻の風ではなく、打ち上げ時刻の風を使う。
パターン自体が東へ移動しているので、
「東側の、あとの時刻の風」は「打ち上げ点の、いまの風」に似ている。
2 つのずれが互いを部分的に埋める。

和歌山の夏はこうならない。
両方固定 \SI{2.8}{\kilo\metre} が時刻のみ \SI{2.2}{\kilo\metre} と
地点のみ \SI{2.0}{\kilo\metre} のどちらより大きく、2 つの誤差がそのまま足し合わさる。
夏の和歌山には東進する強い流れがないので、埋め合わせが起きない。
愛媛の夏も \SI{2.8}{\kilo\metre} 対 \SI{2.2}{\kilo\metre}・\SI{1.9}{\kilo\metre} で、同じく足し合わさる。

\textbf{北海道の夏は、そのどちらでもない。}

季節をまとめたことで生じた見かけではないことを確かめるため、
北海道と和歌山について、月ごとに分けて計算し直した。

\begin{table}[H]
\centering
\caption{両方固定の誤差を、大きいほうの片側固定で割った比。
1 を下回れば打ち消し合い、上回れば足し合わせである。}
\label{tab:cancellation}
\sisetup{table-format=1.2}
\begin{tabular}{@{}lSSSS@{}}
\toprule
月 & {時刻のみ [\si{\kilo\metre}]} & {地点のみ [\si{\kilo\metre}]}
   & {両方 [\si{\kilo\metre}]} & {比} \\
\midrule
\multicolumn{5}{@{}l}{\textbf{北海道}} \\
2018年12月 & 3.96 & 7.10 & 5.92 & 0.83 \\
2019年1月  & 3.98 & 6.55 & 5.85 & 0.89 \\
2019年2月  & 4.20 & 6.83 & 6.07 & 0.89 \\
\addlinespace
2019年6月  & 3.23 & 2.63 & 2.95 & 0.91 \\
2019年7月  & 2.58 & 1.69 & 2.53 & 0.98 \\
2019年8月  & 2.66 & 3.12 & 3.09 & 0.99 \\
\addlinespace
\multicolumn{5}{@{}l}{\textbf{和歌山}} \\
2018年12月 & 3.22 & 6.25 & 5.64 & 0.90 \\
2019年1月  & 3.53 & 8.27 & 6.68 & 0.81 \\
2019年2月  & 3.54 & 6.17 & 5.59 & 0.91 \\
\addlinespace
2019年6月  & 3.33 & 3.49 & 4.06 & 1.16 \\
2019年7月  & 2.21 & 1.80 & 2.56 & 1.16 \\
2019年8月  & 2.11 & 2.08 & 2.74 & 1.30 \\
\bottomrule
\end{tabular}
\end{table}

\textbf{冬は 2 領域とも打ち消し合う。} 6 か月すべてが 1 を下回り、
範囲も \numrange{0.81}{0.91} とよく揃っている。
季節をまとめたことによる見かけではない。

\textbf{夏は 2 領域で違う。}
和歌山の夏は 3 か月とも 1 を上回り（\numrange{1.16}{1.30}）、単純な足し合わせである。
北海道の夏は \numrange{0.91}{0.99} で、1 をわずかに下回るところに貼りついている。
足し合わせでも打ち消し合いでもない。

\begin{whybox}{北海道の夏に「中間」が出るのはなぜか}
打ち消し合いは、擾乱が機体と同じ向き・同じくらいの速さで動くときに起きる。

和歌山の夏は亜熱帯ジェットが北へ抜けていて、動く擾乱がほとんどない。
だから 2 つのずれは独立に効き、足し合わさる。

北海道の夏は違う。ジェットは弱まるが消えはせず、
\SI{200}{\hecto\pascal} で平均 \SI{29.9}{\metre\per\second} 残り、
低気圧も通り続ける（\SI{200}{\hecto\pascal} の偏差が \num{1.8} 日しか続かないのはそのためである）。
移動する擾乱が弱いながらも存在するので、打ち消しが部分的にしか働かない。

つまりこの比は、\textbf{その季節・その場所に「東へ動くパターン」が
どれだけあるか}を測る指標として読める。
冬の 0.8 台、和歌山の夏の 1.2 台、北海道の夏の 0.95 前後という並びは、
そのまま擾乱の活発さの順序になっている。
\end{whybox}

\begin{figure}[H]
\centering
\includegraphics[width=1.0\linewidth]{fidelity/G7-01-wakayama_trajectory_map.pdf}
\caption[冬の 1 例]{2020 年 1 月 15 日 00 UTC に和歌山の中心格子から上げた場合の飛行経路。左が全体、右が着地点の拡大。地点のみ固定（橙の破線）と両方固定（黄）が南側にそろい、基準（黒）と時刻のみ固定（青の破線）が北側にそろう。冬は地点の固定による誤差のほうが大きいことが、1 例からも見てとれる。}
\label{fig:g7-01-wakayama}
\end{figure}

\begin{figure}[H]
\centering
\includegraphics[width=1.0\linewidth]{fidelity/G7-07-wakayama_trajectory_map.pdf}
\caption[夏の 1 例]{同じ図の 7 月版。飛行距離が \SI{80}{\kilo\metre} ほどと短く、4 通りの差も \SI{2.5}{\kilo\metre} にとどまる。拡大パネルの縮尺が冬（\SI{14}{\kilo\metre}）と夏（\SI{2}{\kilo\metre}）で違う点に注意。}
\label{fig:g7-07-wakayama}
\end{figure}

\begin{figure}[H]
\centering
\includegraphics[width=1.0\linewidth]{fidelity/G7-01-hokkaido_trajectory_map.pdf}
\caption[冬の 1 例（北海道）]{2020 年 1 月 15 日 00 UTC に VOISING から上げた場合の飛行経路。
左が全体、右が着地点の拡大。
この例では 4 通りの差が \SI{2.9}{\kilo\metre} しかなく、
和歌山の同じ日（図\ref{fig:g7-01-wakayama}、\SI{13.1}{\kilo\metre}）よりずっと小さい。
和歌山では基準（黒）と時刻のみ固定（青）、
地点のみ固定（橙）と両方固定（黄）がそれぞれ対になるのに対し、
ここでは 4 本が順に並ぶだけで、そうした対応は見えない。
\textbf{1 例は 1 例であって統計ではない}。
どちらが典型かは表\ref{tab:fidelity-error} で判断してほしい。}
\label{fig:g7-01-hokkaido}
\end{figure}

\begin{figure}[H]
\centering
\includegraphics[width=1.0\linewidth]{fidelity/G7-07-hokkaido_trajectory_map.pdf}
\caption[夏の 1 例（北海道）]{同じ図の 7 月版。
拡大パネルの縮尺が冬と夏で違う点に注意。}
\label{fig:g7-07-hokkaido}
\end{figure}
\begin{figure}[H]
\centering
\includegraphics[width=1.0\linewidth]{fidelity/G7-01-ehime_trajectory_map.pdf}
\caption[冬の 1 例]{2020 年 1 月 15 日 00 UTC に南レク御荘公園から上げた場合の飛行経路。左が全体、右が着地点の拡大。和歌山の同じ日（図\ref{fig:g7-01-wakayama}）と同じく、地点のみ固定（橙の破線）と両方固定（黄）が南側に、基準（黒）と時刻のみ固定（青の破線）が北側にそろう。現行との差は \SI{11.1}{\kilo\metre}。}
\label{fig:g7-01-ehime}
\end{figure}

\begin{figure}[H]
\centering
\includegraphics[width=1.0\linewidth]{fidelity/G7-07-ehime_trajectory_map.pdf}
\caption[夏の 1 例]{同じ図の 7 月版。いったん東へ出てから西へ戻り、また東へ流れる経路で、4 通りの差は \SI{2.7}{\kilo\metre} にとどまる。拡大パネルの縮尺が冬（\SI{8}{\kilo\metre}）と夏（\SI{4}{\kilo\metre}）で違う点に注意。拡大パネルの左下に見える線は、飛行の途中で着地点の近くを一度通った部分である。}
\label{fig:g7-07-ehime}
\end{figure}


\subsection{結論：どの季節がいちばんシミュレーション通りか}
\label{sec:fidelity-answer}

答えは、何を「通り」と呼ぶかで変わる。そこで両方を出す。
そして\textbf{北海道と南の 2 領域で答えが違う}。和歌山と愛媛の答えは同じである。

\begin{table}[H]
\centering
\caption{総合。絶対値と割合で順位が逆になり、その逆転の仕方が北と南で違う。}
\label{tab:fidelity-ranking}
\sisetup{table-format=2.1}
\begin{tabular}{@{}lS[table-format=3.0]SS[table-format=1.1]SS[table-format=3.0]@{}}
\toprule
季節 & {飛行距離 [\si{\kilo\metre}]} & {近似誤差 [\si{\kilo\metre}]}
     & {割合 [\%]} & {1 日ずれ [\si{\kilo\metre}]} & {同 p90} \\
\midrule
\multicolumn{6}{@{}l}{\textbf{北海道}} \\
冬 & 176 & 5.1 & 3.1 & 55.5 & 109 \\
春 & 134 & 4.2 & 3.4 & 63.2 & 123 \\
夏 & 73  & 3.2 & 4.6 & 43.0 & 89 \\
秋 & 164 & 4.5 & 2.9 & 55.9 & 110 \\
\addlinespace
\multicolumn{6}{@{}l}{\textbf{和歌山}} \\
冬 & 267 & 5.3 & 2.0 & 47.3 & 90 \\
春 & 165 & 3.9 & 2.5 & 49.8 & 97 \\
夏 & 63  & 2.8 & 4.6 & 29.7 & 66 \\
秋 & 125 & 3.5 & 3.0 & 36.4 & 75 \\
\addlinespace
\multicolumn{6}{@{}l}{\textbf{愛媛}} \\
冬 & 267 & 5.2 & 2.0 & 45.4 & 86 \\
春 & 165 & 3.9 & 2.5 & 48.1 & 94 \\
夏 & 63  & 2.8 & 4.6 & 29.1 & 63 \\
秋 & 124 & 3.4 & 3.0 & 35.1 & 72 \\
\bottomrule
\end{tabular}
\end{table}

\paragraph{絶対的な精度がほしいなら、3 領域とも夏。}
近似誤差が最小（和歌山と愛媛 \SI{2.8}{\kilo\metre}、北海道 \SI{3.2}{\kilo\metre}）で、
飛行距離も最短、1 日ずれたときの移動も最小である。
和歌山と愛媛の夏は \SI{200}{\hecto\pascal} の偏差が \numrange{10.7}{10.8} 日続くので、
前日の観測が翌日も生きる。
北海道の夏は偏差が \num{1.8} 日しか続かないが、それでも 4 季節の中では最良である。

\paragraph{比例的に忠実な季節は、領域で分かれる。}
和歌山と愛媛は\textbf{冬}（どちらも \SI{2.0}{\percent}）。
強いジェットで向きの振れも \SI{200}{\hecto\pascal} で \ang{5.4} と小さく、
「どちらへ飛ぶか」がもっとも予測しやすい。

北海道は\textbf{秋}（\SI{2.9}{\percent}）で、冬（\SI{3.1}{\percent}）はわずかに劣る。
差は小さいが、和歌山のように冬が突出して良いということはない。
北海道の冬のジェットは和歌山ほど安定していないためで、
\SI{500}{\hecto\pascal} の風向が冬でも前日から \ang{20.2} 振れる
（和歌山は \ang{9.0}）。

\paragraph{春は 3 領域とも避けたほうがよい。}
1 日ずれたときの移動が和歌山 \SI{49.8}{\kilo\metre}、愛媛 \SI{48.1}{\kilo\metre}、北海道 \SI{63.2}{\kilo\metre} で、
どちらも 4 季節中最大である。
\SI{200}{\hecto\pascal} の偏差が続く時間も最短（和歌山 \num{1.7} 日、愛媛 \num{1.9} 日、北海道 \num{1.2} 日）。
\SI{850}{\hecto\pascal} の風向は前日から \ang{36} から \ang{43} 振れる。

\paragraph{そして領域どうしを比べると、北海道のほうが一貫して扱いにくい。}
1 日ずれたときの移動は\textbf{どの季節でも和歌山・愛媛より大きく}、
北海道のいちばん良い夏（\SI{43.0}{\kilo\metre}）が
和歌山のいちばん悪い春（\SI{49.8}{\kilo\metre}）に近い。
飛行距離は和歌山の冬のほうが長い（\SI{267}{\kilo\metre} 対 \SI{176}{\kilo\metre}）のに、
それでも移動は和歌山のほうが小さい。
距離ではなく\textbf{風の入れ替わりの速さ}が差を生んでいる。

\paragraph{まとめると}、狭い範囲で確実に回収したいなら 3 領域とも夏。
方向を読みたいなら和歌山と愛媛は冬、北海道は秋。春は避ける。
そして北海道では、季節を選んでも和歌山・愛媛ほどの余裕は得られない。

なお繰り返しになるが、これらはすべて
「予報が完璧だったとして」の話である（\ref{sec:fidelity-scope} 節）。
実際の運用では予報誤差 (4) が上乗せされ、本章の数字より外れる。

\subsection{副産物：既存の閉形式にある地図因子のずれ}
\label{sec:fidelity-mapfactor}

積分器と閉形式を突き合わせる過程で、既存の \texttt{drift.displacements} に
系統的なずれがあることがわかった。報告しておく。

閉形式は $\Delta x = \sum_k u_k T_k$ を直接メートルで出していて、
式のどこにも緯度が現れない。
一方で実際の機体は飛行中に南北にも動き、
経度 1 度あたりの距離を決める $\cos(\text{緯度})$ が飛行中に変わる。

このずれは、南北移動が大きい冬の飛行でいちばん大きくなる。
2020 年 1 月で測ると中央値 \SI{0.44}{\kilo\metre}、p90 で \SI{1.12}{\kilo\metre}、
最大 \SI{1.93}{\kilo\metre}。10 年ぶんの最大は \SI{3.45}{\kilo\metre} であった。
7 月は飛行距離が短いので中央値 \SI{0.07}{\kilo\metre} にとどまる。

第 \ref{sec:hokkaido} 章から第 \ref{sec:ehime} 章の落下位置は
すべてこの閉形式で計算しているので、
そこに載せた平均落下位置には冬でおよそ \SI{0.4}{\kilo\metre}、
最悪で数キロメートルの系統誤差が含まれている。
気候値としての比較（どの時期がどちらへ飛ぶか）には影響しない大きさだが、
1 回の飛行の予測に使うには無視できない。
本章の積分器は球面の幾何を正しく扱っているので、この誤差を持たない。

\section{この結果で言えないこと}
\label{sec:limits}

\subsection{解像度の限界（第 \ref{sec:validation} 章より）}

第 \ref{sec:validation} 章の交差検証から、読み取ってはいけない範囲がはっきりしている。

\begin{itemize}
  \item \textbf{日単位の傾向は無い。} 3 領域とも、1 日ビンは最良より
        6 標準誤差以上悪い。本統計に日単位の集計を置いていないのはそのためである。
  \item \textbf{北海道の半月ごとの違いは読み取らない。}
        月単位と統計的に区別できない（$+0.5$ 標準誤差）。
  \item \textbf{冬の半月ごとの違いは読み取らない。}
        冬の中では、時期を細かく知っても落下位置の予測が改善しない（誤差削減 0 \%）。
  \item \textbf{地上付近にこの統計を使わない。}
        \SI{850}{\hecto\pascal} の誤差削減は 4--6 \% しかない。
        打ち上げ当日の地上風は予報で判断すること。
  \item \textbf{気候値は予報の代わりにならない。}
        最良の粒度でも説明できるのは落下位置の分散の 70 \%（和歌山）、
        72 \%（愛媛）、30 \%（北海道）で、残りは当日の天気である。
        風そのものでみても、季節が説明するのは分散の 45--47 \%（和歌山・愛媛）、
        18 \%（北海道）にとどまる
        （第 \ref{sec:fidelity} 章）。
\end{itemize}

なお「半月単位の傾向がある」という言い方自体が正確ではない。
実在するのは 1--2 か月スケールの滑らかな季節変化で、
半月ビンはそれを階段で近似したものである（第 \ref{sec:validation} 章）。

\subsection{予報の誤差は測っていない（第 \ref{sec:fidelity} 章より）}

第 \ref{sec:fidelity} 章はシミュレータが風を固定していることの誤差を測ったが、
測れるのは 4 つある誤差源のうち 3 つだけである。
\textbf{与えた予報そのものが外れている分は測れていない。}
それには「過去に出された予報」の保管庫が要るのに対し、
本研究が持っているのは「過去に実際どうだったか」の再解析だからである。

したがって第 \ref{sec:fidelity} 章の \SIrange{2.8}{5.3}{\kilo\metre} という数字は、
\textbf{予報が完璧だった場合の下限}である。
実際の運用では予報誤差が上乗せされ、これより外れる。
実務ではそちらのほうが大きい可能性が高い。

\subsection{落下位置には地図因子の系統誤差が含まれる（第 \ref{sec:fidelity} 章より）}

第 \ref{sec:hokkaido} 章から第 \ref{sec:ehime} 章の落下位置は
$\Delta x = \sum_k u_k T_k$ という閉形式で計算している。
この式には緯度が現れないので、飛行中に機体が南北へ動いて
経度 1 度あたりの距離が変わる効果を持っていない。

第 \ref{sec:fidelity} 章で球面の幾何を正しく扱う積分器と突き合わせたところ、
和歌山の冬でずれは中央値 \SI{0.44}{\kilo\metre}、
10 年ぶんの最大で \SI{3.45}{\kilo\metre} であった。
ビンどうしの比較には影響しない大きさだが、
1 回の飛行の予測に使う場合はこの分を見込む必要がある。

\subsection{愛媛の領域はデータ範囲の西端に接している}

取得したデータの範囲（\ang{132}--\ang{150}E）は、愛南町の西 \SI{51}{\kilo\metre} で終わる。
そのため愛媛の領域は打ち上げ地点を中心にできず、東へ伸ばしてある
（第 \ref{sec:ehime} 章）。

影響は 2 つある。
第 \ref{sec:fidelity} 章の軌道計算では、愛媛から上げた軌道の \SI{5.9}{\percent} が
西端からデータの外へ出て、集計から外れた。
外れたのは夏に西へ流れた短い飛行が中心で、その偏りは誤差の中央値で 2 \% 未満である。
また、豊後水道より西（九州側）の風は本報告のどの図にも入っていない。
夏の西向きの飛行について、九州寄りの風の違いを確かめるには範囲を西へ広げた取り直しが要る。

\subsection{再解析は観測ではない}

第 \ref{sec:data} 章で述べたとおり、JRA-3Q の値は観測に強く拘束されたモデル推定である。
日本の陸上の対流圏では観測が密なので実態に近いが、
外洋上や高度 \SI{30}{\kilo\metre} 以上では、
その場を直接測ったものは少なく、モデルの寄与が相対的に大きい。
本報告で最もスペースバルーンに関係する高度帯（\SIrange{25}{35}{\kilo\metre}）は、
ちょうどその「観測が薄い」側にある。

\subsection{解像度が地形を解いていない}

\ang{0.375} は約 \SI{40}{\kilo\metre} である。
日高山脈も紀伊山地も四国山地も、この格子では 2--3 点でしか表現されない。
したがって本報告は\textbf{谷筋の風や斜面下降風のような地形性の現象を一切含まない}。
打ち上げ地点そのものの地上風は、この統計ではなく現地の観測で判断すべきである。

上空数キロ以上では地形の影響が薄れるので、
スペースバルーンの飛行の大部分についてはこの制約は緩い。

\subsection{10 年は経年変動をならすには短い}

成層圏の東西風には準 2 年周期振動（QBO）があり、
エルニーニョ・南方振動は対流圏の循環を年ごとに動かす。
10 年（QBO でいえば約 4--5 周期）では、これらの位相が完全にはならされない。
あるビンの値が「その時期のふつう」なのか
「たまたまこの 10 年でそうだった」のかは、本報告では分離できない。

より長い期間を使えばこの点は改善するが、
今度は観測網の変化（衛星観測の導入など）が入ってくる。
どちらを取るかの判断であって、片方が正しいわけではない。

\subsection{気圧面上の統計であって等高度面ではない}

統計は気圧面ごとに取っている。
ある気圧面の実際の高度は季節や天気で数百メートル動くので、
「高度 \SI{20}{\kilo\metre} の統計」と読み替えるときには、
その高度が平均値であることを意識する必要がある。
高度の標準偏差も出力に含めてあるので、必要ならそちらを見てほしい。

\subsection{半月の平均は「その日」を教えてくれない}

これは気候値というものの性質そのものである。
「3 月前半の平均風は西 \SI{30}{\metre\per\second}」と分かっても、
来年の 3 月 10 日がそうだとは限らない。
分位点と分散楕円を併記しているのは、
平均からどれだけずれうるかまで含めて計画に使えるようにするためである。
実際に飛ばす日の判断には、この統計ではなく予報を使うこと。

\clearpage
\appendix

% Number the appendices "付録A" rather than a bare "A", and keep the
% subsections as A.1 so they stay short. Because this goes through
% \thesection, a \ref to an appendix prints "付録E" on its own and the
% citing sentence does not add the word itself.
\renewcommand{\thesection}{付録\Alph{section}}
\renewcommand{\thesubsection}{\Alph{section}.\arabic{subsection}}
\addtocontents{toc}{\protect\appendixtocwidth}

\section{出力ファイル}

\begin{longtable}{lp{0.55\linewidth}}
\toprule
場所 & 内容 \\
\midrule
\endhead
\texttt{out/netcdf/} & 基準となる出力。CF-1.10 準拠の NetCDF4。次元は (month, bin, level, lat, lon)。
  モーメント（$\bar u$, $\bar v$, $\sigma_u^2$, $\sigma_v^2$, $\sigma_{uv}$）と
  分位点・標本数・高度を保持する。 \\
\texttt{out/parquet/} & 同じ内容の long format。region と month で分割。
  モーメントに加えて風向・楕円・シアーなど導出量も入れてある。 \\
\texttt{out/csv/} & 厳選した 6 表（後述）。 \\
\texttt{out/figures/} & 本報告の図。PNG（確認用）と PDF（本報告用）。 \\
\texttt{data/raw/} & GDEX から取得した切り出し済みキャッシュと、
  ファイルごとの検証結果を記録した JSON。 \\
\bottomrule
\end{longtable}

\section{CSV 表の一覧}

\begin{longtable}{lp{0.6\linewidth}}
\toprule
ファイル & 内容 \\
\midrule
\endhead
\texttt{bin\_calendar.csv} & 24 ビンの定義と標本数。区切り方の監査用。 \\
\texttt{region\_grid\_*.csv} & 実際に使った格子点。 \\
\texttt{region\_sites\_*.csv} & 打ち上げ地点と最寄り格子、そのずれ（km）。 \\
\texttt{site\_profile\_*.csv} & 代表地点の (月, ビン, 気圧面) プロファイル。 \\
\texttt{layer\_summary\_*.csv} & 飛行の各段階に対応する 6 層への集約。 \\
\texttt{best\_launch\_halves\_*.csv} & 明示した目的関数によるビンの順位付け。 \\
\texttt{data\_inventory.csv} & 取得元ファイルの一覧と検証結果。 \\
\bottomrule
\end{longtable}

\texttt{best\_launch\_halves} の順位は\textbf{最適ではなく、明示された 1 つの基準}による。
すなわち「0--30 km の領域平均風速の鉛直平均が小さい順、同点なら定常度が高い順」である。
別の基準が欲しければ Parquet から付け直せる。

\section{再現手順}

\codeblock
\begin{lstlisting}
conda env create -f environment.yml
conda activate jra3q
cd research

python -m windclim.cli fetch --years 2016-2025     # 約2.5時間、中断可
for r in hokkaido wakayama ehime; do
  python -m windclim.cli all --region $r
  python -m windclim.cli validate --region $r   # 第 3b 章
  python -m windclim.cli fidelity --region $r   # 第 6b 章
done
python -m windclim.cli compare                  # 3 領域がそろってから

# 付録の 7 日ビン版（本文と別のファイルに書き出される）
python -m windclim.cli all --region wakayama --scheme week
python -m windclim.cli all --region hokkaido --scheme week

cd report && make
\end{lstlisting}

\texttt{fetch} は中断しても安全で、再実行すれば途中から続く。
オフラインのテストは \texttt{pytest}、
GDEX の仕様が変わっていないかの確認は \texttt{pytest -m network} で走る。

\section{和歌山の冬：風速の分布と、そこから出る落下点の散らばり}
\label{sec:appendix-winter}

第 \ref{sec:wakayama} 章で、和歌山の冬は\textbf{風向が安定していて風速がばらつく}
という結果になった。ここではその「ばらつく」を分布として推定し、
その分布から落下点を実際に撒いてみる。

打ち上げ地点は\textbf{和歌山領域の中心格子（「和歌山中央」）}1 点に固定する。
\ang{33.902}N, \ang{135.375}E で、第 \ref{sec:wakayama} 章の各図が扱う
領域と同じ格子である。
用いるのは 12・1・2 月の全解析 \num{3612} 個。

\subsection{風速の分布は一様ではないし、1 つの形でもない}

\begin{figure}[H]
\centering
\includegraphics[width=1.0\linewidth]{wakayama/A1_speed_fits.pdf}
\caption[冬の風速分布]{和歌山の冬、代表的な 6 つの気圧面での風速のヒストグラムと、
当てはめた分布（青）。比較のため一様分布（橙の破線）も重ねてある。
採用した分布の名前・KS 統計量・一様分布との $\Delta$AIC を各パネルに記した。
採用した分布が面によって正規と Weibull に分かれるので、
凡例ではなくパネルごとに名前を書いている。}
\label{fig:a1-wakayama}
\end{figure}


38 の気圧面それぞれについて、正規分布・Weibull 分布・一様分布を当てはめ、
AIC で比べた。

\begin{table}[H]
\centering
\caption{和歌山の冬、代表的な気圧面での当てはめ。
$\Delta$AIC は採用した分布との差で、大きいほどその分布が劣る。
風向 SD は比較のために併記した。}
\label{tab:winter-fit}
\sisetup{table-format=2.2}
\begin{tabular}{@{}llS[table-format=1.2]S[table-format=2.1]S[table-format=2.1]S[table-format=4.0]S[table-format=2.1]@{}}
\toprule
気圧面 & 採用 & {KS} & {平均 [\si{\metre\per\second}]} & {SD}
       & {一様との $\Delta$AIC} & {風向 SD [度]} \\
\midrule
\SI{850}{\hecto\pascal} & 正規    & 0.01 & 10.7 & 4.2  & 3954 & 44.2 \\
\SI{500}{\hecto\pascal} & 正規    & 0.01 & 34.9 & 10.5 & 3981 & 16.8 \\
\SI{200}{\hecto\pascal} & Weibull & 0.03 & 75.9 & 14.3 & 2421 & 11.5 \\
\SI{100}{\hecto\pascal} & 正規    & 0.01 & 48.9 & 8.5  & 3023 & 9.1 \\
\SI{30}{\hecto\pascal}  & Weibull & 0.03 & 9.5  & 6.4  & 3502 & 58.3 \\
\SI{10}{\hecto\pascal}  & Weibull & 0.02 & 12.1 & 8.6  & 4672 & 96.0 \\
\bottomrule
\end{tabular}
\end{table}

\textbf{一様分布は 38 面すべてで負ける。} $\Delta$AIC は最小でも \num{1878}、
最大 \num{4833} である。
「風速がばらばら」という言い方は、中心のない一様な散らばりを思わせるが、
実際には\textbf{はっきりした中心と、行儀のよい裾}を持っている。

形は 1 つではない。38 面のうち 16 面が正規、22 面が Weibull だった。
分かれ方には理由がある。

\begin{whybox}{なぜ高度で形が変わるのか}
\SI{500}{\hecto\pascal} や \SI{850}{\hecto\pascal} では、平均風速が
標準偏差の 3--8 倍ある。風速が 0 に近づくことはまず無いので、
下限 0 という制約が問題にならず、正規分布でよく合う。

\SI{10}{\hecto\pascal} や \SI{30}{\hecto\pascal} では平均が \SIrange{9}{12}{\metre\per\second}、
SD が \SIrange{6}{9}{\metre\per\second} で、変動係数が 0.7 近い。
風速は負になれないので分布は 0 で切られ、右に裾を引く。
これが Weibull の形である。

\SI{200}{\hecto\pascal} だけは別の理由で Weibull が勝つ。
ここは平均 \SI{75.9}{\metre\per\second} と速く 0 からは遠いが、
分布が\textbf{左}に歪んでいる（歪度 $-0.44$）。
ジェットの強さに上限があるためで、Weibull の形状パラメータが大きくなると
左に歪んだ形も表せる。裾の重さを合わせた結果であって、
\SI{10}{\hecto\pascal} と同じ理由ではない。
\end{whybox}

\subsection{層を独立に振ってはいけない}

分布が決まったので、そこから風速のプロファイルを撒いて飛ばす。
ここで最初の注意点がある。

各気圧面を\textbf{独立に}引くのがいちばん単純なやり方だが、それは間違いである。
飛行は全層を足し合わせるので、独立な誤差は互いに打ち消し合い、
落下点の雲が実際よりずっと小さく出る。
実際の風はそうなっていない。
\SI{200}{\hecto\pascal} と \SI{100}{\hecto\pascal} の風速は相関 0.64 で連動する。

そこで観測された順位相関からガウス・コピュラを作り、
各面の周辺分布は当てはめたまま、層どうしの連動を保って撒いた。

コピュラが何をしているかを書いておく。
やりたいことは 2 つあり、そのままでは両立しない。
ひとつは\textbf{各面の分布の形を守ること}である。
前項のとおり、ある面は正規分布、別の面は Weibull 分布が当てはまっており、
面ごとに形が違う。
もうひとつは\textbf{層どうしの連動を入れること}である。

単純に 38 次元の多変量正規分布から引けば連動は入る。
ところがその場合、どの面も正規分布になってしまう。
Weibull が当てはまった面の形が壊れるので、これは使えない。

コピュラは、この 2 つを\textbf{分けて}扱う道具である。
連動だけを正規分布で作り、
分布の形はあとから面ごとに与え直す。
3 段階で書くと次のようになる。

\begin{equation*}
\underbrace{\vec{z} \sim \mathcal{N}(\vec{0},\, R)}_{\text{(1) 連動だけを作る}}
\quad\Longrightarrow\quad
\underbrace{u_j = \Phi(z_j)}_{\text{(2) 一様}[0,1]\text{に均す}}
\quad\Longrightarrow\quad
\underbrace{V_j = F_j^{-1}(u_j)}_{\text{(3) その面の分布に戻す}}
\end{equation*}

$j$ は気圧面の番号（$j = 1, \dots, 38$）、
$\Phi$ は標準正規分布の累積分布関数、
$F_j$ は $j$ 面に当てはめた分布（正規または Weibull）、
$F_j^{-1}$ はその分位点関数である。

段階ごとに見る。
(1) では相関行列 $R$ を持つ正規乱数を引く。
この時点で持っている情報は\textbf{層どうしの連動だけ}で、
分布の形はまだ決めていない。
(2) で $\Phi$ を通すと、$u_j$ は各成分が $[0, 1]$ 上の一様分布になる。
$\Phi$ は単調増加なので、ここで\textbf{連動は壊れない}。
形の情報だけが抜けて、順位の情報だけが残る。
(3) で各面の分位点関数を通すと、$V_j$ はその面に当てはめた分布に\textbf{厳密に}従う。
$F_j^{-1}$ も単調増加なので、ここでも連動は壊れない。

つまり (2) と (3) はどちらも単調変換で、順位を入れ替えない。
連動が (1) から (3) まで素通りし、分布の形だけが (3) で入る。
これが「周辺分布は当てはめたまま、連動を保つ」の中身である。

相関行列 $R$ には、観測の相関をそのまま入れるのではなく、
一段変換したものを入れる。

\begin{equation}
\rho^{(R)}_{jk} = 2 \sin\!\left( \frac{\pi}{6} \rho^{\mathrm{S}}_{jk} \right)
\label{eq:copula-rho}
\end{equation}

$\rho^{\mathrm{S}}_{jk}$ は観測から測った $j$ 面と $k$ 面の\textbf{順位相関}
（Spearman）である。
ここで Pearson ではなく順位相関を使うのには理由がある。
順位相関は単調変換で変わらない量なので、
(2) と (3) を通り抜けても値が保たれる。
一方 Pearson 相関は (3) の $F_j^{-1}$ で変わってしまうので、
入れた値と出てくる値が一致しない。

式 \eqref{eq:copula-rho} は、
「(1) に $\rho^{(R)}$ を入れると (3) の出力の順位相関が
$\rho^{\mathrm{S}}$ になる」という関係である。
ガウス・コピュラではこれが閉じた形で書ける。
$\rho^{\mathrm{S}} = 0.64$ なら $\rho^{(R)} = 0.66$ で、
差は小さいが、1 行で正確にできるものを近似で済ませる理由もない。

最後に 1 つ処理がいる。
$R$ は面のペアごとに別々に測った値を並べたものなので、
相関行列として成立しているとは限らない（正定値でないことがある）。
そのままでは (1) の乱数生成（Cholesky 分解）が失敗するので、
固有値を下から $10^{-8}$ で押さえて最も近い正定値行列に直し、
対角を 1 に直してから使う。

\begin{whybox}{ガウス・コピュラで足りない可能性がある点}
ガウス・コピュラは層間の依存を\textbf{相関行列だけ}で表す。
これで表せないのは\textbf{裾どうしの連動}である。

具体的には、「強風の日は全層いっせいに極端に強い」という性質を考える。
正規分布の連動は、中央付近では強く、裾へ行くほど弱くなる。
したがって実際の大気で裾が中央より強く連動していると、
このモデルは極端な日を\textbf{過小に}出す。

落下点の分布は観測をよく再現しているので大きな誤りは無いはずだが、
裾の連動そのものを確かめたわけではない。
$t$ コピュラなど裾の連動を持つ族と比べるのが筋だが、本報告では行っていない。
\end{whybox}

\subsection{比べる 5 つの条件}

以降の図と表には 5 つの条件が出てくる。
いちばん上が観測そのもので、2 番目がいちばん単純なモデルである。
そこから下へ、\textbf{モデルに要素を 1 つずつ足していく}。
したがって下へ行くほど観測に近づくはずである。
近づかなければ、足した要素に効果がないか、まだ足りない要素が残っている。

\begin{table}[H]
\centering
\caption{5 つの条件の中身。同じ列で値が変わっている行が、
その条件で新しく入った要素である。}
\label{tab:winter-scenarios}
\footnotesize
\begin{tabular}{@{}lllll@{}}
\toprule
条件 & 風速 & 風向 & 層間の連動 & 機体諸元 \\
\midrule
観測（実プロファイル）
  & 実測値
  & 実測値
  & 実測値（仮定なし）
  & 固定 \\
\addlinespace
風速のみ・層独立
  & 当てはめた分布
  & 層ごとの平均に固定
  & \textbf{なし}
  & 固定 \\
風速のみ・鉛直相関あり
  & 当てはめた分布
  & 層ごとの平均に固定
  & \textbf{コピュラ}
  & 固定 \\
＋風向のばらつき
  & 当てはめた分布
  & \textbf{平均のまわりに散らす}
  & コピュラ
  & 固定 \\
＋機体（\texttt{frontend} 既定）
  & 当てはめた分布
  & 平均のまわりに散らす
  & コピュラ
  & \textbf{摂動} \\
\bottomrule
\end{tabular}
\end{table}

1 つずつ中身を書く。

\paragraph{観測（実プロファイル）。}
モデルを一切通さない。
12・1・2 月の解析 \num{3612} 個の $(u, v)$ プロファイルを、
そのまま 1 本ずつ飛ばしたものである。
分布の当てはめもコピュラも使っていないので、
\textbf{これが他の 4 条件の答え合わせになる}。
モデルの雲がこの雲に重なれば、そのモデルは風の散らばりを再現できている。

\paragraph{風速のみ・層独立。}
ここからがモデルである。
各気圧面の風速を、前項で当てはめた分布（正規または Weibull）から引く。
ただし\textbf{面ごとに独立に}引く。
風向は各面の平均風向に\textbf{固定}し、散らさない。
機体諸元も固定する。

つまりこの条件は、
「各高度の風速がどれくらいばらつくか」だけを知っていて、
「その高度どうしがどう連動するか」を知らない場合にあたる。
高度別のヒストグラムだけを見てモンテカルロを組むと、これになる。

\paragraph{風速のみ・鉛直相関あり。}
1 つ前との違いは\textbf{層間の連動だけ}である。
周辺分布も、風向を平均に固定することも、機体を固定することも変えていない。
引き方をコピュラに変えただけで、他は同一の乱数種から引いている。

したがってこの 2 条件の差は、\textbf{鉛直相関を入れたことだけの効果}である。

\paragraph{＋風向のばらつき。}
風速は 1 つ前と同じ（コピュラつき）まま、
風向を平均に固定するのをやめる。
各面の平均風向のまわりに、その面の循環 SD で散らす。
風向どうしも同じコピュラで連動させる。

この条件が答えるのは、
「和歌山の冬は風向が安定している」という第 \ref{sec:wakayama} 章の結果を受けて、
\textbf{では風向を固定してよいのか}という問いである。

\paragraph{＋機体（\texttt{frontend} 既定）。}
風は 1 つ前と同じまま、機体の諸元を摂動する。
上昇速度・破裂高度・降下速度を、
Web アプリのモンテカルロ機能の既定値
（\texttt{frontend/src/mc/sampling.ts}）と同じ分布・同じ大きさで散らす。
上昇速度と降下速度は SD \SI{0.5}{\metre\per\second}、
破裂高度は 5 \% の正規分布である。

これが\textbf{現行の Web アプリが出している散らばりに、
実際の風のばらつきを足したもの}にあたる。
アプリ自身は風を摂動しないので、アプリ単体の雲はこれよりずっと小さい。

\begin{whybox}{5 条件すべてに共通して入っていないもの}
条件ごとの違いは上のとおりだが、\textbf{5 つとも共通に置いている仮定}もある。
図を読むときはこちらも効く。

\begin{itemize}
  \item \textbf{打ち上げ地点は 1 点固定}である。
        「和歌山中央」の \ang{33.902}N, \ang{135.375}E だけを使う。
        打ち上げ場所の違いによる散らばりは入っていない。
  \item \textbf{飛行中に風は変わらない。}
        1 本のプロファイルを飛行全体（約 \num{140} 分）に使い回す。
        これは現行シミュレータと同じ仮定で、
        この仮定自体の誤差は第 \ref{sec:fidelity} 章で別に測った。
  \item \textbf{機体が水平に動いても、参照する風は変わらない。}
        上と同じ理由で、打ち上げ格子の風だけを使う。
  \item \textbf{予報誤差は入っていない。}
        ここで使っているのは再解析、つまり\textbf{後から作った最良推定}である。
        実運用では打ち上げ前に GFS 予報を使うので、
        その予報が外れるぶんの誤差がさらに乗る。
        本報告の手持ちデータでは測れない（第 \ref{sec:limits} 章）。
\end{itemize}

したがって\textbf{観測の雲でさえ、実際の散らばりの下限である}。
\end{whybox}

\subsection{結果}

\begin{figure}[H]
\centering
\includegraphics[width=1.0\linewidth]{wakayama/A2_landing_map.pdf}
\caption[落下点を地図に重ねる]{★の「和歌山中央」（\ang{33.902}N, \ang{135.375}E）から
冬に上げたスペースバルーンの落下点を、
標高と水深に重ねたもの。各条件 \num{4000} 点、観測は \num{3612} 点。
条件の中身は表 \ref{tab:winter-scenarios} にある。
点の描画順は広い雲が先、狭い雲が後なので、
いちばん狭いマゼンタ（層独立）が最前面にある。
地形・水深は NOAA NCEI の DEM を陸と海で別の色階調にし、陰影を重ねたもの。
海岸線は標高 0 の等値線である。}
\label{fig:a2-wakayama}
\end{figure}

\begin{figure}[H]
\centering
\includegraphics[width=1.0\linewidth]{wakayama/A3_landing_clouds.pdf}
\caption[落下点の散らばり（条件別）]{図\ref{fig:a2-wakayama} と同じ点を、
条件ごとに切り離して打ち上げ地点からの変位で示したもの。
灰色は比較用に重ねた観測。
地図の上では重なって見分けにくい 3 つの差が、ここでは分けて読める。}
\label{fig:a3-wakayama}
\end{figure}


\begin{table}[H]
\centering
\caption{落下点の散らばり。★（打ち上げ地点）からの変位で、
東西が飛行方向、南北がそれと直交する方向にあたる。}
\label{tab:winter-spread}
\sisetup{table-format=2.1}
\begin{tabular}{@{}lSSS[table-format=3.1]S[table-format=3.1]@{}}
\toprule
条件 & {東西 SD [\si{\kilo\metre}]} & {南北 SD [\si{\kilo\metre}]}
     & {距離中央値} & {距離 p90} \\
\midrule
観測（実プロファイル）        & 54.5 & 48.6 & 265.8 & 330.9 \\
\addlinespace
風速のみ・層独立            & 15.0 & 3.5  & 270.7 & 289.7 \\
風速のみ・鉛直相関あり        & 49.2 & 5.2  & 270.6 & 332.5 \\
＋風向のばらつき            & 47.7 & 48.8 & 262.2 & 322.1 \\
＋機体（\texttt{frontend} 既定） & 52.2 & 49.5 & 262.5 & 332.2 \\
\bottomrule
\end{tabular}
\end{table}

読み取れることが 3 つある。

図\ref{fig:a2-wakayama} が全体像で、図\ref{fig:a3-wakayama} が条件ごとの切り分けである。
地図のほうは、落下点が\textbf{ほぼすべて海上}で、
しかも南海トラフから伊豆・小笠原の海山域にかかることを示している。
和歌山から冬に上げる場合、回収は洋上前提になる。

\paragraph{1. 層を独立に振ると、散らばりが 3 分の 1 になる。}
東西 SD が \SI{15.0}{\kilo\metre} で、相関を入れた \SI{49.2}{\kilo\metre} の
3 割にしかならない。
分布そのものは同じなのに、層のつなぎ方だけでこれだけ変わる。
モンテカルロで風を摂動するときは、\textbf{層間の相関を入れないと
不確かさを大幅に過小評価する}。

\paragraph{2. 風向を固定すると、横方向の散らばりがほとんど消える。}
相関を入れた風速だけのモデルは、東西 SD を \SI{49.2}{\kilo\metre} と
ほぼ正しく出す（観測 \SI{54.5}{\kilo\metre}）。
ところが南北 SD は \SI{5.2}{\kilo\metre} で、観測の \SI{48.6}{\kilo\metre} に
遠く及ばない。

\begin{whybox}{「風向が安定している」と「風向が落下点を動かさない」は別}
\SI{200}{\hecto\pascal} の定常度は 0.98 で、これは十分に安定と言ってよい値である。
しかしそれでも円環標準偏差は \ang{11.5} ある。

冬の平均飛行距離は \SI{264}{\kilo\metre}。
\ang{11.5} の振れは、横方向に
$264 \times \tan(\ang{11.5}) = \SI{54}{\kilo\metre}$ の散らばりを生む。
観測の南北 SD \SI{48.6}{\kilo\metre} はこれでほぼ説明がつく。

風向を各面の円環標準偏差で振り直すと、南北 SD は \SI{48.8}{\kilo\metre} になり、
観測の \SI{48.6}{\kilo\metre} とほぼ一致する（表\ref{tab:winter-spread} の 4 行目）。

\textbf{定常度が高いことは、風向を固定してよい理由にならない。}
遠くまで飛ぶほど、わずかな角度が大きな距離になる。
\end{whybox}

\paragraph{3. 機体パラメータの寄与は小さい。}
\texttt{frontend} のモンテカルロが振っているのは
上昇速度（$\sigma = \SI{0.5}{\metre\per\second}$）、
破裂高度（$\sigma = 5\ \%$）、
降下速度（$\sigma = \SI{0.5}{\metre\per\second}$）の 3 つで、
\textbf{風は振っていない}（\texttt{frontend/src/mc/sampling.ts}）。

これを風の摂動に上乗せすると、東西 SD が \SI{47.7}{\kilo\metre} から
\SI{52.2}{\kilo\metre} へ、南北 SD が \SI{48.8}{\kilo\metre} から
\SI{49.5}{\kilo\metre} へ動くだけである。

\begin{whybox}{現行のモンテカルロは不確かさの大部分を見ていない}
機体の 3 パラメータだけを振ると、落下点の散らばりは
東西・南北とも \SI{5}{\kilo\metre} 程度にしかならない。
実際の散らばりは \SI{50}{\kilo\metre} 級である。

つまり\texttt{frontend} のモンテカルロが示している雲は、
\textbf{実際に起こりうる散らばりの 1 割程度}である。
残りは風が決めており、そこは現在まったく摂動されていない。

本節のモデル（表\ref{tab:winter-spread} の 5 行目）は観測をよく再現するので、
風の摂動を追加する道筋はある。
必要なのは気圧面ごとの周辺分布と層間の相関行列で、
どちらも \texttt{out/csv/winter\_speed\_fit\_wakayama.csv} と
\texttt{windclim.stats.windsample} にある。
\end{whybox}

\subsection{この付録の限界}

\begin{itemize}
  \item \textbf{1 格子点、1 季節、1 領域}の結果である。
        他の季節や北海道では相関構造も分布の形も違う。
  \item 飛行モデルは\textbf{風を固定する}既存の形（\ref{sec:drift} 節）を使った。
        \texttt{frontend} と揃えるためで、
        第 \ref{sec:fidelity} 章で測った固定による誤差（数キロメートル）は
        ここには入っていない。
        ただし \SI{50}{\kilo\metre} 級の散らばりに対しては小さい。
  \item 地図の標高・水深は NOAA NCEI の DEM mosaic サービスから
        1 度だけ取得してキャッシュしたものである。
        図を描く処理はネットワークに触れない。
        このサービスは著作権表示を返さない。
        米国連邦政府の著作物は通例パブリックドメインであり、
        NOAA の水深プロダクトもふつうその扱いだが、
        サービス自身がそう明示しているわけではない。
        この地図を再配布する場合は各自で条件を確認してほしい。
  \item このサービスは、要求した緯度経度の箱を\textbf{そのまま返すとは限らない}。
        箱の縦横比が要求した画像の縦横比と合わない場合、
        画像に合わせて箱のほうを広げた画像が返ってくる。
        たとえば \ang{134.5}--\ang{141.5}E を $1440 \times 840$ 画素で要求すると、
        返るのは \ang{133.29}--\ang{142.71}E である。
        要求した箱を座標軸に使うと、地図全体がなめらかにずれる。
        一見もっともらしく見えるので気づきにくい。
        そこで取得側は先にメタデータ（\texttt{f=json}）を問い合わせ、
        サービスが申告した範囲をサイドカーに記録し、
        描画はその範囲だけを使う（\texttt{fetch/terrain.py}）。
        あわせて、箱と同じ縦横比の画素数を要求することで、
        要求どおりの箱が返るようにしている。
  \item 再現は \texttt{python -m windclim.cli winter --region wakayama --fetch-dem}
        で行う。初回だけ DEM（\SI{7.1}{\mega\byte}）を取得する。
\end{itemize}

\section{7 日ビンによる旧版の統計}
\label{sec:appendix-week}

本報告は当初、1 か月を 7 日 $\times$ 4 に分けて集計していた
（W1 = 1--7 日、W2 = 8--14 日、W3 = 15--21 日、W4 = 22 日--月末）。
第 \ref{sec:validation} 章の時間解像度の検証を受けて 15 日に改めたが
（\ref{sec:weeks} 節）、旧版もここに残す。

残す理由は 2 つある。第一に、7 日の結果が 15 日と食い違っていないことを
読者が自分で確かめられるようにするため。
第二に、より細かい時期区分が必要な用途のためである。
7 日は「区別できない範囲」の中にあるので、使うこと自体は妨げられない。

\begin{whybox}{図は保存物ではなく、再生成できる}
この付録の図は、集計をやり直して作ったものである。

\begin{lstlisting}
python -m windclim.cli all --region wakayama --scheme week
python -m windclim.cli all --region hokkaido --scheme week
\end{lstlisting}

7 日ビンの出力はファイル名に \texttt{\_week} が付き、
本文の 15 日ビンの出力を上書きしない。
古い図をそのまま保管しておくのではなくコードで再現できるようにしたのは、
本報告の図がすべて \texttt{out/figures/} を直接読む設計だからである。
再生成できない図を 1 枚でも混ぜると、その保証が崩れる。
\end{whybox}

\subsection{標本数}

\begin{figure}[H]
\centering
\includegraphics[width=0.95\linewidth]{hokkaido_week/F0_sample_counts.pdf}
\caption[標本数（7 日ビン、北海道）]{7 日ビンでの標本数。
W1--W3 が各 280 個、W4 だけが 400 個（1 月）や 292 個（2 月）になる。
15 日ビンではこの凹凸が小さくなる（本文 図\ref{fig:f0-hokkaido}）。}
\label{fig:app-f0-hokkaido}
\end{figure}

\subsection{高度 $\times$ 時期の風速}

\begin{figure}[H]
\centering
\includegraphics[width=1.0\linewidth]{hokkaido_week/F1_alt_half_speed.pdf}
\caption[高度 $\times$ 週の平均風速（北海道）]{北海道。横軸は 48 ビン。
本文の図\ref{fig:f1-hokkaido}（24 ビン）と同じ構造が見える。
細かくしても新しい構造が現れるわけではない、というのが第 \ref{sec:validation} 章の
結論であり、この 2 枚を並べるとそれが目で確かめられる。}
\label{fig:app-f1-hokkaido}
\end{figure}

\begin{figure}[H]
\centering
\includegraphics[width=1.0\linewidth]{wakayama_week/F1_alt_half_speed.pdf}
\caption[高度 $\times$ 週の平均風速（和歌山）]{和歌山。同じく 48 ビン。}
\label{fig:app-f1-wakayama}
\end{figure}

\subsection{落下位置}

\begin{figure}[H]
\centering
\includegraphics[width=1.0\linewidth]{wakayama_week/M5_drift_month_grid.pdf}
\caption[月別の平均落下位置（7 日ビン、和歌山）]{7 日ビンでの平均落下位置。
矢印は実寸。本文の図\ref{fig:m5-wakayama} と比べても、
矢印の向きと長さはほとんど変わらない。}
\label{fig:app-m5-wakayama}
\end{figure}

\subsection{旧版の出力ファイル}

\begin{longtable}{@{}p{0.42\linewidth}p{0.52\linewidth}@{}}
\toprule
場所 & 内容 \\
\midrule
\endhead
\texttt{out/netcdf/}\newline\texttt{jra3q\_windclim\_*\_week\_v1.nc} & 7 日ビンの統計。
  次元は本文版と同じ (month, bin, level, lat, lon) で、
  \texttt{bin} の長さが 4 になる。 \\
\texttt{out/figures/<region>\_week/} & 7 日ビンの図一式。 \\
\texttt{out/csv/*\_week.csv} & 7 日ビンの CSV 表（\texttt{bin\_calendar\_week.csv} を含む）。 \\
\bottomrule
\end{longtable}

どちらの NetCDF も、自分がどの区切り方で作られたかを
\texttt{bin\_scheme}、\texttt{bin\_labels}、\texttt{bin\_start\_days} の
属性に持っている。図の軸ラベルも表の列名もこの属性から作られるので、
2 つの製品を取り違えて読むことはできない。

\section{ライセンスと帰属}

本報告および付随データは JRA-3Q の二次的著作物であり、
\textbf{CC BY-NC-SA 4.0}（非商用・継承・表示）を継承する。
リポジトリ本体のライセンスとは別である。

引用：\citet{kosaka2024jra3q}。
データ提供：気象庁、配布：NSF NCAR GDEX（データセット \texttt{d640000}）\citep{gdex_d640000}。

\section{ERA5 の 1 時間値で見た時間帯の違い}
\label{sec:era5-diurnal}

本文の統計は、JRA-3Q の 6 時間ごとの解析値（09・15・21・03 時、日本時間）を
時刻で分けずにまとめたものである。
第 \ref{sec:method} 章ではその根拠として、夜（03 時）と昼（15 時）の差が
\SI{1}{\metre\per\second} 未満であることを示した（図\ref{fig:f8-hokkaido}）。
ただし 6 時間ごとの値では、1 日のうち 4 つの時刻しか見ていない。

そこで、1 時間ごとの値がある ERA5 \citep{hersbach2020era5} を北海道について 1 か月ぶん取り、
朝・昼・夕・晩で風に違いがあるかを調べた。
コードと取得の手順は \texttt{research/alpha/} にある。

\subsection{データと方法}

\begin{table}[H]
\centering
\caption{使ったデータと、時間帯の切り方}
\label{tab:era5-setup}
\begin{tabular}{@{}lp{0.72\linewidth}@{}}
\toprule
データ & ERA5 気圧面解析、NSF NCAR GDEX \citep{gdex_d633000} \\
格子・面 & \ang{0.25} 格子、37 気圧面（\SIrange{1000}{1}{\hecto\pascal}） \\
領域 & 本文の北海道と同じ（\ang{40.0}--\ang{46.5}N / \ang{138.0}--\ang{149.0}E、27 $\times$ 45 点） \\
期間 & 2025 年 7 月。日本時間で 24 時間そろう 7 月 2--31 日の 30 日 \\
時間帯（日本時間） & 朝 05--10 時、昼 11--16 時、夕 17--22 時、晩 23--04 時 \\
\bottomrule
\end{tabular}
\end{table}

月は 7 月にした。夏は日射による地表付近の変化（海陸風や混合層の発達）が最も強く、
時間帯の違いが出るならこの季節に出るはずだからである。
2025 年は本文の統計の期間（2016--2025 年）に含まれる。

日々の天気の違いは時間帯の違いよりずっと大きい。
そのまま時間帯ごとに平均すると、たまたま荒れた日がどの時間帯に多く入ったかで結果が変わってしまう。
そこで\textbf{各日（日本時間の 0--23 時）の平均からの偏差}をとり、その偏差を時間帯ごとに平均した。

信頼区間は、日を 1 つの標本として計算した。
同じ日の 6 時間は互いによく似ているので、1 時間ずつを別の標本に数えると区間が狭く出すぎる。
まず時間帯の 6 時間を日ごとに平均し、その 30 個の値の標準誤差から 95 \% 信頼区間を求めた。

落下点は本文と同じ閉形式（式\eqref{eq:drift}、上昇 \SI{5}{\metre\per\second}、
破裂 \SI{30}{\kilo\metre}、降下 \SI{6}{\metre\per\second}）で、
打ち上げ時刻の風の鉛直分布から計算した。
地点は本文の代表地点（VOISING）に最も近い格子点（\ang{42.25}N, \ang{142.75}E）である。

\subsection{風速の 1 日の変化}

\begin{figure}[H]
\centering
\includegraphics[width=\linewidth]{alpha/E1-202507_hourly_curves.pdf}
\caption[1 時間ごとの風速の変化]{領域平均風速の、日平均からの偏差を 1 時間ごとに並べたもの。
線は 30 日の平均、帯は 95 \% 信頼区間。背景の色が 4 つの時間帯である。
縦軸の目盛りは面ごとに違う。
975 hPa で 06→07 時と 18→19 時に見える段差は、大気の変化ではなく ERA5 の作り方から来ている
（\ref{sec:era5-jumps} 節）。}
\label{fig:e1}
\end{figure}

\begin{figure}[H]
\centering
\includegraphics[width=0.95\linewidth]{alpha/E2-202507_block_profiles.pdf}
\caption[時間帯ごとの偏差の高度分布]{時間帯ごとの平均から日平均を引いた値を、高度に対して描いたもの。
左は風速、右は風ベクトル（向きの変化も含む）の偏差の大きさ。帯は 95 \% 信頼区間。}
\label{fig:e2}
\end{figure}

\begin{table}[H]
\centering
\caption{時間帯ごとの風速の偏差（領域平均、\si{\metre\per\second}、$\pm$ は 95 \% 信頼区間）。
右端は比べるための物差しで、日平均風速そのものの日々の標準偏差である。}
\label{tab:era5-blocks}
\small
\begin{tabular}{@{}lrrrrr@{}}
\toprule
気圧面（高度） & {朝} & {昼} & {夕} & {晩} & {日々の標準偏差} \\
\midrule
\SI{975}{\hecto\pascal}（\SI{0.3}{\kilo\metre}） & $-0.23 \pm 0.19$ & $+0.04 \pm 0.13$ & $-0.07 \pm 0.32$ & $+0.26 \pm 0.24$ & 2.79 \\
\SI{850}{\hecto\pascal}（\SI{1.5}{\kilo\metre}） & $+0.00 \pm 0.20$ & $+0.03 \pm 0.20$ & $-0.04 \pm 0.32$ & $+0.01 \pm 0.31$ & 3.57 \\
\SI{500}{\hecto\pascal}（\SI{5.8}{\kilo\metre}） & $-0.08 \pm 0.34$ & $+0.15 \pm 0.23$ & $+0.09 \pm 0.48$ & $-0.16 \pm 0.37$ & 3.42 \\
\SI{200}{\hecto\pascal}（\SI{12.4}{\kilo\metre}） & $-0.59 \pm 0.54$ & $+0.03 \pm 0.42$ & $+0.64 \pm 0.88$ & $-0.08 \pm 0.70$ & 6.01 \\
\SI{50}{\hecto\pascal}（\SI{21.0}{\kilo\metre}） & $+0.44 \pm 0.23$ & $-0.45 \pm 0.28$ & $-0.02 \pm 0.28$ & $+0.04 \pm 0.17$ & 1.17 \\
\SI{30}{\hecto\pascal}（\SI{24.3}{\kilo\metre}） & $+0.66 \pm 0.26$ & $-0.60 \pm 0.30$ & $-0.10 \pm 0.22$ & $+0.04 \pm 0.20$ & 1.14 \\
\SI{10}{\hecto\pascal}（\SI{31.6}{\kilo\metre}） & $+0.52 \pm 0.39$ & $-0.44 \pm 0.39$ & $-0.10 \pm 0.29$ & $+0.03 \pm 0.40$ & 1.07 \\
\SI{1}{\hecto\pascal}（\SI{48.7}{\kilo\metre}） & $+3.25 \pm 0.63$ & $-1.17 \pm 0.74$ & $-3.71 \pm 0.82$ & $+1.63 \pm 0.85$ & 3.38 \\
\bottomrule
\end{tabular}
\end{table}

結果は高度によって 3 つに分かれる。

\paragraph{対流圏（地上--約 \SI{16}{\kilo\metre}）では、時間帯の違いはほとんど無い。}
高度 \SI{3}{\kilo\metre} より下では、時間帯ごとの偏差は最大でも \SI{0.3}{\metre\per\second} で、
日々の標準偏差（\SIrange{2}{3.6}{\metre\per\second}）の 1 割程度しかない。
信頼区間が 0 をまたがない組もあるが（975 hPa の朝 $-0.23$、晩 $+0.26$ など）、
大きさは本文の統計の精度より小さい。
対流圏上部の \SI{200}{\hecto\pascal} では夕が $+0.64$ とやや大きいが、
信頼区間（$\pm 0.88$）が広く、0 と区別できない。

\paragraph{成層圏（約 \SIrange{20}{32}{\kilo\metre}）には、はっきりした日変化がある。}
\SIrange{50}{10}{\hecto\pascal} では、風速が朝 6 時ごろに最も強く、
昼 14--16 時ごろに最も弱くなる。
1 日の中の振れ幅は \SI{30}{\hecto\pascal} で \SI{1.8}{\metre\per\second}、
\SI{10}{\hecto\pascal} で \SI{1.4}{\metre\per\second} である。
朝と昼の偏差は信頼区間が 0 をまたがず、
しかも日々の標準偏差（約 \SI{1.1}{\metre\per\second}）と同じくらいの大きさがある。
夏の成層圏は風が弱く日々の変動も小さいので、1 日の中の変化が相対的に目立つ。

\paragraph{破裂高度より上（約 \SI{35}{\kilo\metre} 以上）では、日変化はさらに大きい。}
\SI{1}{\hecto\pascal}（約 \SI{49}{\kilo\metre}）では 1 日の振れ幅が \SI{8}{\metre\per\second} を超える。
ただしこの高さは破裂高度 \SI{30}{\kilo\metre} より上なので、落下点には関係しない。

\begin{whybox}{成層圏の日変化は大気潮汐である}
成層圏のオゾン層は太陽の紫外線を吸収して昼に暖まり、夜に冷える。
この加熱が 1 日周期で繰り返されるので、気圧と風も 1 日周期で振動する。
これを大気潮汐と呼ぶ。月の引力による海の潮汐とは別の現象で、原因は日射である。

大気潮汐の振幅は高度とともに大きくなるので、
\SI{30}{\hecto\pascal} より \SI{1}{\hecto\pascal} のほうが振れ幅が大きい。
本文の図\ref{fig:g1-hokkaido} で、\SI{10}{\hecto\pascal} と \SI{30}{\hecto\pascal} の
風の変化量が遅れ 24 時間でいったん下がるのも、同じ潮汐による。
\end{whybox}

\subsection{地表付近の昼と晩}

\begin{figure}[H]
\centering
\includegraphics[width=\linewidth]{alpha/E3-202507_day_night_map.pdf}
\caption[昼と晩の風の差]{昼（11--16 時）の平均風から晩（23--04 時）の平均風を引いたベクトル。
色がその大きさ。左は \SI{975}{\hecto\pascal}（約 \SI{0.3}{\kilo\metre}）、
右は \SI{850}{\hecto\pascal}（約 \SI{1.5}{\kilo\metre}）。
山地では \SI{975}{\hecto\pascal} 面が地面の下にあり、その値は外挿である。}
\label{fig:e3}
\end{figure}

領域平均では小さかった地表付近の差も、格子点ごとに見ると大きいところがある。
\SI{975}{\hecto\pascal} では昼と晩の差が最大 \SI{3.5}{\metre\per\second}、
\SI{850}{\hecto\pascal} でも最大 \SI{2.6}{\metre\per\second} ある。
差が大きいのは海岸や山地の周りで、海陸風や山谷風の影響と考えられる。

それでも領域平均の差が小さくなるのは、差の向きが場所によってばらばらで、
平均すると打ち消し合うからである。
打ち上げ地点ごとの地上風は、本文の第 \ref{sec:limits} 章に書いたとおり、
統計ではなく当日の予報と現地の観測で判断する必要がある。

\subsection{ERA5 の解析窓の境目での跳び}
\label{sec:era5-jumps}

図\ref{fig:e1} の \SI{975}{\hecto\pascal} には、
06→07 時に $-0.65$、18→19 時に $-0.55$ \si{\metre\per\second} の段差がある。
1 時間でこれだけ変わるのは不自然なので、1 時間ごとの風の変化量を時刻別に調べた。

\begin{figure}[H]
\centering
\includegraphics[width=\linewidth]{alpha/E6-202507_hourly_jumps.pdf}
\caption[1 時間ごとの風の変化量]{1 時間ごとの風ベクトルの変化量 $|\mathbf{V}(t+1)-\mathbf{V}(t)|$ を
格子点ごとに求めて領域平均し、変化の始まる時刻（UTC）ごとに並べたもの。
各面の中央値を 1 とした比で描いている。
点線が ERA5 の 12 時間の解析窓の境目（09 UTC と 21 UTC）である。}
\label{fig:e6}
\end{figure}

変化量は、09→10 UTC と 21→22 UTC の 2 か所だけが飛び抜けて大きい。
\SI{850}{\hecto\pascal} ではほかの時刻の 1.7--1.9 倍、
\SI{500}{\hecto\pascal} で 1.3--1.6 倍、\SI{10}{\hecto\pascal} で 1.1--1.2 倍である。
09 UTC と 21 UTC は、日本時間の 18 時と 6 時にあたる。

ERA5 は、09--21 UTC と 21--09 UTC の 12 時間ごとに観測を取り込んで解析を作っている。
1 つの解析窓の中では時間変化が滑らかにつながるが、
窓と窓の境目では、次の窓の観測に合わせて風が不連続に修正される。
2 か所の跳びはその境目と一致しているので、\textbf{大気の変化ではなく解析の作り方による見かけの変化}と考えられる。

したがって、ERA5 の 1 時間値で見た地表付近の「朝の風速低下」の一部は、この跳びによるものである。
1 時間値を使って日変化を調べるときは、この 2 か所をまたぐ比較に注意が要る。
一方、成層圏の日変化（朝に強く昼に弱い）は境目と関係のない滑らかな形をしており、
跳びも小さいので、実在の変化と考えてよい。

\subsection{JRA-3Q との比較}

\begin{figure}[H]
\centering
\includegraphics[width=0.6\linewidth]{alpha/E5-202507_vs_jra3q.pdf}
\caption[ERA5 と JRA-3Q の日変化の比較]{夜（03 時）と昼（15 時）の平均風速の差を、
ERA5（2025 年 7 月）と JRA-3Q（2016--2025 年の 7 月）で比べたもの。
北海道領域の平均。JRA-3Q の時刻別の統計は \SI{300}{\hecto\pascal} より下の面だけにある。}
\label{fig:e5}
\end{figure}

本文の図\ref{fig:f8-hokkaido} と同じ量（03 時 $-$ 15 時）を ERA5 で計算すると、
高度 \SI{4}{\kilo\metre} より下では JRA-3Q とほぼ重なる。
高度 \SIrange{0.5}{1}{\kilo\metre} 付近で夜のほうが \SI{0.2}{\metre\per\second} ほど強い点も共通である。
それより上では ERA5 のほうが昼に強い側へずれるが、差は ERA5 の信頼区間の中に収まっている。
ERA5 は 1 か月、JRA-3Q は 10 年ぶんの平均なので、上層の差は年ごとのばらつきと区別できない。

なお JRA-3Q の 4 つの解析時刻（09・15・21・03 時）は、
朝・昼・夕・晩の 4 つの時間帯に 1 回ずつ入っている。
本文の統計は 4 つの時間帯を同じ重みで含んでいるので、
時間帯による偏りは平均の中で相殺されている。

\subsection{落下点への影響}

\begin{figure}[H]
\centering
\includegraphics[width=\linewidth]{alpha/E4-202507_drift_by_launch_hour.pdf}
\caption[打ち上げ時刻と落下点]{代表地点から上げた場合の落下点を、日平均の落下点からの偏差で示したもの。
左は打ち上げ時刻（日本時間）ごと、右は時間帯ごとの平均。帯と誤差棒は 95 \% 信頼区間。
風は打ち上げ時刻・地点で固定している（本文の閉形式と同じ）。}
\label{fig:e4}
\end{figure}

\begin{table}[H]
\centering
\caption{時間帯ごとの落下点の偏差（代表地点、2025 年 7 月、\si{\kilo\metre}）}
\label{tab:era5-drift}
\begin{tabular}{@{}lrr@{}}
\toprule
時間帯 & {東西（東が正）} & {南北（北が正）} \\
\midrule
朝 & $-1.1 \pm 3.7$ & $+1.2 \pm 2.9$ \\
昼 & $+2.6 \pm 2.4$ & $+2.1 \pm 3.0$ \\
夕 & $-0.4 \pm 5.4$ & $-2.7 \pm 4.5$ \\
晩 & $-1.1 \pm 4.1$ & $-0.5 \pm 4.5$ \\
\bottomrule
\end{tabular}
\end{table}

2025 年 7 月の代表地点の平均飛行距離は \SI{51.5}{\kilo\metre} である。
時間帯ごとの落下点の偏差はどれも \SI{3}{\kilo\metre} 以内で、
信頼区間が 0 をまたがないのは昼の東西（$+2.6 \pm 2.4$）だけである。
昼に上げた場合と晩に上げた場合の差は、東へ \SI{3.7}{\kilo\metre}、北へ \SI{2.6}{\kilo\metre}、
合わせて約 \SI{4.5}{\kilo\metre} になる。

この差を高度帯に分けると、地上--\SI{3}{\kilo\metre} の寄与は東西 \SI{0.1}{\kilo\metre}、南北 $-\SI{0.8}{\kilo\metre}$ と小さく、
大部分は \SIrange{3}{20}{\kilo\metre}（東西 \SI{2.2}{\kilo\metre}、南北 \SI{1.5}{\kilo\metre}）と
\SIrange{20}{30}{\kilo\metre}（東西 \SI{1.3}{\kilo\metre}、南北 \SI{1.8}{\kilo\metre}）から来ている。
地表付近の日変化は上昇・降下とも短時間で通り過ぎるので、落下点にはほとんど影響しない。

比べるための物差しとして、同じ月・同じ地点で
打ち上げが 24 時間ずれたときの落下点の移動は中央値 \SI{35.3}{\kilo\metre}、
日ごとの平均落下点のばらつき（標準偏差）は \SI{47.3}{\kilo\metre} だった。
時間帯による違い（約 \SI{4.5}{\kilo\metre}）はその 1 割程度で、
本文の第 \ref{sec:fidelity} 章で測った「風を固定したことによる誤差」（夏 \SI{3.2}{\kilo\metre}）と同じくらいの大きさである。

\subsection{まとめと限界}

\begin{itemize}
  \item 対流圏の風には、朝・昼・夕・晩ではっきりした違いは無い。
        領域平均の差は \SI{0.8}{\metre\per\second} 以下で、日々の変動の 1 割程度である。
  \item 成層圏（約 \SIrange{20}{32}{\kilo\metre}）には大気潮汐による日変化があり、
        風速は朝 6 時ごろに強く、昼 15 時ごろに弱い（振れ幅 \SIrange{1.4}{1.8}{\metre\per\second}）。
  \item 落下点への影響は、昼と晩の差で約 \SI{4.5}{\kilo\metre} である。
        1 日ずれたときの移動（\SI{35}{\kilo\metre}）よりずっと小さい。
  \item 以上から、本文で統計を時刻で分けなかった判断は、少なくとも夏の北海道については妥当である。
  \item ERA5 の 1 時間値には、09 UTC と 21 UTC の解析窓の境目で不自然な跳びがある。
        1 時間値で日変化を調べるときは、この 2 か所に注意が要る。
\end{itemize}

この付録は 1 年の 1 か月だけを見た試験的な解析であり、次の限界がある。

\begin{itemize}
  \item 30 日の平均なので、信頼区間が広い。とくに対流圏上部の差は年による違いと区別できない。
  \item 夏だけを見ている。冬は成層圏の風が強く日々の変動も大きいので、
        大気潮汐の相対的な大きさは夏より小さくなると予想されるが、確かめていない。
  \item 落下点は、風を打ち上げ時刻・地点で固定した閉形式で計算している。
        飛行中の 2 時間余りの間にも風は時刻とともに変わるが、その効果は入れていない。
\end{itemize}

ERA5 は Copernicus Climate Change Service の製品であり、
本付録の図は Copernicus Climate Change Service の情報を加工したものを含む。

\section{1・2 月に西日本で「雪のない平野」や瀬戸内海に落とせる場所}
\label{sec:winter-sw}

本文の 3 領域は北海道・和歌山・愛媛に限られていた。
本付録では対象を関東より西・南の全域に広げ、1・2 月に次の条件を満たす打ち上げ場所を探す。

\begin{itemize}
  \item 打ち上げ場所に雪がない。
  \item 落下点が陸の平野で、人口集中地区（DID）の外で、雪がない。
\end{itemize}

海上での回収は考えない。したがって海に落ちた回は失敗として数える。
コードは \texttt{research/winter\_sw/} にある。

\subsection{データと方法}

\begin{table}[H]
\centering
\caption{使ったデータ}
\label{tab:wsw-data}
\small
\begin{tabular}{@{}lp{0.68\linewidth}@{}}
\toprule
風 & JRA-3Q（本文と同じ製品）。北緯 24.0--37.5 度、東経 122.0--140.5 度を新たに取得した。
     2016--2025 年の 1・2 月、6 時間ごと \\
地形 & NOAA NCEI の DEM（約 \SI{500}{\metre} 間隔） \\
市街地 & 国土数値情報「人口集中地区データ」2020 年（国土交通省） \\
積雪 & ERA5 の積雪深（水当量）\citep{hersbach2020era5}。毎日 12 時（日本時間）、\ang{0.25} 格子 \\
\bottomrule
\end{tabular}
\end{table}

落下点は本文と同じ閉形式（式\eqref{eq:drift}）で、解析 1 回ごとに求めた。
1・2 月の 10 年ぶんで 2360 回の解析があり、打ち上げ場所ごとにその数だけ落下点が得られる。
落下点の分布に正規分布を仮定する必要はない。

打ち上げ場所の候補は、\ang{0.1} 間隔に並べた陸上の点 2151 か所である。
風はそれぞれの最寄りの JRA-3Q 格子点のものを使った。
あわせて、各地の代表地点 25 か所も同じ方法で評価した。

1 回の打ち上げを「成功」と数える条件は、次のすべてを満たすことである。

\begin{itemize}
  \item 落下点が\textbf{平野}：標高 \SI{100}{\metre} 以下、傾斜 \ang{2} 以下の陸地。
  \item 落下点が \textbf{DID の外}。
  \item その日、落下点と打ち上げ場所の\textbf{どちらにも雪がない}：積雪深（水当量）\SI{1}{\milli\metre} 未満。
\end{itemize}

雪の判定には、打ち上げと同じ日の積雪を使った。
10 年ぶんの平均ではなく、その日ごとに判定している。

\begin{figure}[H]
\centering
\includegraphics[width=\linewidth]{winter_sw/W1_target_and_snow.pdf}
\caption[落下の目標と雪のない日の割合]{左は落下の目標（緑）と DID（赤）。
陸地のうち平野は 20 \% で、平野の 12 \% が DID に入る。
右は 1・2 月に雪のない日の割合。濃いほど雪の日が多い。}
\label{fig:wsw-base}
\end{figure}

\subsection{1・2 月の落下位置の傾向}

\begin{figure}[H]
\centering
\includegraphics[width=\linewidth]{winter_sw/W2_drift_jan_feb.pdf}
\caption[1・2 月の平均落下位置]{標準の飛行条件（上昇 \SI{5}{\metre\per\second}、
破裂 \SI{30}{\kilo\metre}、降下 \SI{6}{\metre\per\second}）での平均落下位置。
点が打ち上げ地点、矢先が平均落下点で、矢印は実寸である。}
\label{fig:wsw-drift}
\end{figure}

1・2 月は、どこから上げても\textbf{ほぼ真東}（方位 \ang{82}--\ang{89}）へ流される。
飛行距離は本土で \SIrange{240}{270}{\kilo\metre}、南ほど短く、
奄美で \SI{233}{\kilo\metre}、那覇で \SI{197}{\kilo\metre}、石垣で \SI{163}{\kilo\metre} である。
冬の亜熱帯ジェットがこの緯度帯の真上にあるためで、
第 \ref{sec:wakayama} 章で見た和歌山の冬と同じ性質が、西日本全体に広がっている。

1 月前半・1 月後半・2 月前半・2 月後半の違いは小さい。
本土では 2 月前半が最も遠く（例：大阪 \SI{271}{\kilo\metre}）、
2 月後半が最も近い（同 \SI{247}{\kilo\metre}）が、差は約 \SI{25}{\kilo\metre} で、向きはほとんど変わらない。
荒れた日（p90）には、本土で \SIrange{315}{345}{\kilo\metre} に達する。

\subsection{標準の飛行条件での結果}

\begin{figure}[H]
\centering
\includegraphics[width=\linewidth]{winter_sw/W3_suitability_standard.pdf}
\caption[標準の飛行条件で目標に落ちる割合]{打ち上げ場所ごとに、
「雪がなく、DID の外の平野」に落ちる割合を半月ごとに示したもの。
打ち上げ場所に雪がある日と、海に落ちた日は失敗に数えている。}
\label{fig:wsw-std}
\end{figure}

\begin{table}[H]
\centering
\caption{代表地点から標準の飛行条件で上げた場合（1・2 月の平均）。
「目標」は雪がなく DID の外の平野に落ちる割合、「平野（DID 含む）」は DID の中も平野として数えた場合の値、「海」は海に落ちる割合で、
「打ち上げ地に雪」は打ち上げ場所に雪がある日の割合である。}
\label{tab:wsw-sites}
\small
\setlength{\tabcolsep}{4pt}
\begin{tabular}{@{}p{0.3\linewidth}rrrrr@{}}
\toprule
地点 & {距離 [\si{\kilo\metre}]} & {目標} & {平野（DID 含む）} & {海} & {打ち上げ地に雪} \\
\midrule
岡山 & 260 & \textbf{13 \%} & 17 \% & 32 \% & 14 \% \\
名古屋 & 261 & \textbf{12 \%} & 21 \% & 52 \% & 13 \% \\
松江 & 253 & 6 \% & 7 \% & 27 \% & 17 \% \\
大阪 & 262 & 5 \% & 7 \% & 57 \% & 18 \% \\
広島 & 259 & 5 \% & 9 \% & 31 \% & 36 \% \\
鳥取 & 253 & 5 \% & 7 \% & 7 \% & 50 \% \\
徳島・和歌山・津 & 263--265 & 4 \% & 5--7 \% & 73--85 \% & 10--11 \% \\
福岡・松山 & 259--262 & 3 \% & 4--5 \% & 35--44 \% & 7--18 \% \\
浜松 & 266 & 3 \% & 3 \% & 92 \% & 6 \% \\
高知・大分・熊本 & 261--265 & 1 \% & 1--2 \% & 77--82 \% & 2--5 \% \\
愛南・室戸・潮岬・宮崎・鹿児島 & 263--269 & 0 \% & 0 \% & 94--100 \% & 0--3 \% \\
種子島・奄美・那覇・石垣 & 163--257 & 0--1 \% & 0--1 \% & 98--100 \% & 0 \% \\
金沢 & 246 & 0 \% & 0 \% & 3 \% & 97 \% \\
\bottomrule
\end{tabular}
\end{table}

結果ははっきりしている。
\textbf{標準の飛行条件では、1・2 月を通した平均で目標に落ちる割合が 25 \% を超える場所は無い}。
平均で最も高いのは岐阜県南部（北緯 35.6 度、東経 137.1 度付近）の 22 \%（半月ごとに 18--24 \%）で、
落下先は関東平野である。
半月ごとに見ると、2 月後半の長野県（北緯 36.0 度、東経 137.9 度付近）の 30 \% が最大である。
代表地点では岡山（13 \%）と名古屋（12 \%）が最も高い。

理由は 2 つある。

\paragraph{東へ 250 km 先に、大きな平野が少ない。}
九州・四国の南岸・紀伊半島・東海から上げると、落下点のほとんどが太平洋に出る
（愛南・潮岬・宮崎などは 94--100 \%）。
中国地方や近畿の内陸から上げると、今度は中部・近畿の山地に落ちる。
250 km 東に平野があるのは、岐阜・名古屋（→ 関東平野）、岡山（→ 濃尾平野・伊勢湾岸）などに限られる。

\paragraph{その平野は大都市圏である。}
250 km 東にある平野の多くは、関東・濃尾・大阪の大都市圏で、DID が広い。
DID を除くと、名古屋では 21 \% が 12 \% に、岐阜県南部では 30--34 \% が 18--24 \% に下がる。

日本海側は、打ち上げ場所の雪で失敗する日が多い。鳥取で約 5 割、金沢で 97 \% である。

\subsection{飛行条件を変えた場合}

飛行距離が短くなれば、落下点は打ち上げ場所の近くの平野に入りやすくなる。
そこで上昇速度と破裂高度を変えた 5 つの条件で、同じ評価を行った。
降下速度は海面相当 \SI{6}{\metre\per\second} で共通にした。

\begin{figure}[H]
\centering
\includegraphics[width=\linewidth]{winter_sw/W4_scenarios.pdf}
\caption[飛行条件ごとの比較]{飛行条件ごとの「目標に落ちる割合」（1・2 月の平均）。
丸は最も割合が高い地点。見出しの飛行距離は、打ち上げ候補全体の平均である。}
\label{fig:wsw-scen}
\end{figure}

\begin{table}[H]
\centering
\caption{飛行条件ごとの飛行距離と、最良地点で目標に落ちる割合（1・2 月の平均）}
\label{tab:wsw-scen}
\begin{tabular}{@{}lrrr@{}}
\toprule
飛行条件 & {飛行距離 [\si{\kilo\metre}]} & {最良地点} & {目標に落ちる割合} \\
\midrule
標準（上昇 5 m/s・破裂 30 km） & 255 & 岐阜県南部 & 22 \% \\
低い破裂（5 m/s・20 km） & 239 & 長野県 & 23 \% \\
速い上昇（8 m/s・30 km） & 188 & 長野県 & 36 \% \\
速い上昇・破裂 25 km & 187 & 長野県 & 37 \% \\
\textbf{速い上昇・破裂 20 km} & \textbf{177} & \textbf{長野県} & \textbf{41 \%} \\
\bottomrule
\end{tabular}
\end{table}

破裂高度を下げるだけでは、飛行距離はあまり縮まない（\SI{255}{\kilo\metre} → \SI{239}{\kilo\metre}）。
第 \ref{sec:wakayama} 章の最後で見たとおり、冬の落下距離の約 6 割は
高度 \SIrange{8}{16}{\kilo\metre} のジェットの層で生じるので、
その層を速く通り抜けるほうが効果が大きい。
上昇速度を \SI{8}{\metre\per\second} にすると、飛行距離は約 7 割（\SI{188}{\kilo\metre}）になる。

\begin{figure}[H]
\centering
\includegraphics[width=\linewidth]{winter_sw/W5_suitability_best.pdf}
\caption[速い上昇・破裂 20 km で目標に落ちる割合]{最も成績のよい飛行条件
（上昇 \SI{8}{\metre\per\second}、破裂 \SI{20}{\kilo\metre}）での、半月ごとの結果。}
\label{fig:wsw-best}
\end{figure}

最も成績のよい条件（上昇 \SI{8}{\metre\per\second}、破裂 \SI{20}{\kilo\metre}）では、
長野県の佐久〜八ヶ岳付近（北緯 36.0--36.1 度、東経 138.2--138.4 度）から上げた場合に、
37--45 \% が目標に落ちる。
落下先は北関東の平野（埼玉・群馬・茨城の DID の外）である。

\begin{figure}[H]
\centering
\includegraphics[width=\linewidth]{winter_sw/W6_landing_clouds.pdf}
\caption[落下点の分布]{最良地点から上げた場合の、解析 1 回ごとの落下点。
上は標準の飛行条件、下は上昇 \SI{8}{\metre\per\second}・破裂 \SI{20}{\kilo\metre}。
青が目標に落ちた回、灰色がそれ以外（DID・山地・海、または雪の日）。}
\label{fig:wsw-clouds}
\end{figure}

\begin{table}[H]
\centering
\caption{地方ごとの最良地点（1・2 月の平均）。地方は、海岸線で本州・九州・四国と島を分けたうえで、
本州の中を経度・緯度で分けた。「最良の条件」は 5 つの飛行条件のうち最もよいもの。}
\label{tab:wsw-areas}
\small
\begin{tabular}{@{}lp{0.42\linewidth}rr@{}}
\toprule
地方 & 最良地点（標準条件 / 最良の条件） & {標準条件} & {最良の条件} \\
\midrule
中部（長野・岐阜・山梨・北陸） & 岐阜県南部 / 長野県 佐久〜八ヶ岳 & 22 \% & \textbf{41 \%} \\
東海（静岡・愛知・三重） & 愛知県北部 / 山梨・静岡県境（富士川沿い） & 17 \% & 20 \% \\
近畿 & 兵庫県南西部 / 兵庫県中部 & 12 \% & 19 \% \\
中国 & 岡山県 & 14 \% & 15 \% \\
四国 & 香川県 & 9 \% & 9 \% \\
九州 & 福岡県北九州 / 長崎県北部 & 7 \% & 7 \% \\
離島 & 小豆島 / 五島列島 & 12 \% & 15 \% \\
南西諸島 & 先島・沖縄本島 & 3 \% & 4 \% \\
\bottomrule
\end{tabular}
\end{table}

中部以外の地方では、飛行条件を変えても 20 \% 前後かそれ以下にとどまる。
四国と九州の本土はどの条件でも 10 \% に届かない。
東へ 180--260 km の範囲が、ほとんど海（太平洋や瀬戸内海）か山地になるためである。
離島では、上昇を速くすると五島列島から佐賀・熊本の平野に落とせる日が増え、15 \% になる。
南西諸島は、どの条件でも落下点のほとんどが海になる。

\subsection{瀬戸内海への着水を認める場合}
\label{sec:wsw-seto}

ここまでは海に落ちた回をすべて失敗に数えた。
波の穏やかな瀬戸内海であれば船で回収できるので、瀬戸内海への着水も成功に数えて評価し直した。
瀬戸内海の範囲は、関門海峡・豊予海峡（関崎〜佐田岬）・鳴門海峡・紀淡海峡を境とする内側の海とし、
外洋に近い紀伊水道と豊後水道は含めない（別府湾は含める）。
着水した回は落下点の積雪を問わないが、打ち上げ場所に雪がある日は引き続き失敗に数える。
コードでは \texttt{python winter\_sw/analyze.py --allow-seto} で切り替わる。

\begin{figure}[H]
\centering
\includegraphics[width=\linewidth]{winter_sw_seto/W3_suitability_standard.pdf}
\caption[瀬戸内海を含めた場合の割合（標準の飛行条件）]{瀬戸内海への着水も成功に数えた場合の、
標準の飛行条件での「目標に落ちる割合」。色の目盛りは図\ref{fig:wsw-std} と違い、70 \% までである。}
\label{fig:wsw-seto-std}
\end{figure}

結果は大きく変わる。
1・2 月は真東へ約 \SI{250}{\kilo\metre} 流されるので、
九州北西部や山口県西部から上げると、ちょうど周防灘から伊予灘・広島湾に落ちる。

\begin{table}[H]
\centering
\caption{瀬戸内海への着水を認めた場合の、地方ごとの最良地点（1・2 月の平均）}
\label{tab:wsw-seto-areas}
\small
\begin{tabular}{@{}lp{0.42\linewidth}rr@{}}
\toprule
地方 & 最良地点（標準条件 / 最良の条件） & {標準条件} & {最良の条件} \\
\midrule
離島 & 壱岐 & \textbf{47 \%} & \textbf{60 \%} \\
九州 & 佐賀県北部（唐津・呼子付近） & 37 \% & \textbf{62 \%} \\
中国 & 山口県西部（長門付近） & 39 \% & 47 \% \\
中部（長野・岐阜・山梨・北陸） & 岐阜県南部 / 長野県 佐久〜八ヶ岳 & 22 \% & 41 \% \\
東海（静岡・愛知・三重） & 愛知県北部 / 山梨・静岡県境 & 17 \% & 20 \% \\
四国 & 愛媛県北部（松山の北） & 10 \% & 19 \% \\
近畿 & 兵庫県南西部 / 兵庫県中部 & 12 \% & 19 \% \\
南西諸島 & 先島 & 3 \% & 4 \% \\
\bottomrule
\end{tabular}
\end{table}

\paragraph{壱岐は標準の飛行条件のままで 47 \% に達する。}
半月ごとに 40--60 \% で、2 月後半が最も高い。
成功した回のほとんど（全体の 36--55 \%）が瀬戸内海への着水で、陸の平野に落ちるのは 5 \% 前後である。
1・2 月に打ち上げ場所に雪がある日は、10 年を通して無かった。
ただし壱岐のような小さな島では、ERA5 の格子のほとんどが海になるので、
積雪が少なめに出ている可能性がある（\ref{sec:wsw-notes} 節）。

\paragraph{上昇を速くすると、九州北部の本土から 6 割に届く。}
上昇 \SI{8}{\metre\per\second} にすると飛行距離が約 \SI{190}{\kilo\metre} に縮み、
佐賀県北部（唐津・呼子付近）から上げた場合に 55--66 \% が目標に入る。
打ち上げ場所に雪がある日は 1--5 \% である。

\begin{figure}[H]
\centering
\includegraphics[width=\linewidth]{winter_sw_seto/W6_landing_clouds.pdf}
\caption[瀬戸内海を含めた場合の落下点]{上は壱岐から標準の飛行条件で、
下は唐津付近から上昇 \SI{8}{\metre\per\second}・破裂 \SI{20}{\kilo\metre} で上げた場合の、解析 1 回ごとの落下点。
水色が瀬戸内海、青い点が目標に入った回である。
外れた回は、手前の北九州・山口の山地か、瀬戸内海を越えた四国の山地、または豊後水道側の海に落ちている。}
\label{fig:wsw-seto-clouds}
\end{figure}

瀬戸内海を目標に含めても、打ち上げ場所の選択肢は西日本の西の端に限られる。
東へ \SI{250}{\kilo\metre} 先が瀬戸内海になるのは、九州北西部・壱岐・山口県西部だけである。
それより東から上げると、瀬戸内海を通り越して近畿・中部の山地や太平洋に落ちる。

\subsection{読むときの注意}
\label{sec:wsw-notes}

\begin{itemize}
  \item \textbf{瀬戸内海は船の往来が多い。}
        着水を認める場合は、回収船の手配に加えて、航路との位置関係を事前に確認する必要がある。
        また瀬戸内海の境界は本付録で決めたもので、法令上の「瀬戸内海」より狭い。
  \item \textbf{「DID の外の平野」は、回収しやすい場所と同じではない。}
        北関東の平野は DID の外でも人家や道路が多い。
        また首都圏の上空は、成田・羽田の空港に近い空域にかかる。
        打ち上げの届出や飛行の可否は、この統計とは別に確認する必要がある。
  \item \textbf{ERA5 の積雪は \ang{0.25}（約 \SI{25}{\kilo\metre}）格子である。}
        沿岸の都市でも、同じ格子に入る山地の雪を拾う。
        広島（36 \%）や大阪（18 \%）の「打ち上げ地に雪」は、実際より多めに出ている可能性が高い。
        逆に壱岐のような小さな島では、格子のほとんどが海になるので少なめに出る。
  \item \textbf{平野の判定は標高と傾斜だけで行っている。}
        湖（琵琶湖など、湖面の標高が \SI{1}{\metre} を超えるもの）、河川、森林、農地を区別していない。
  \item \textbf{最良地点の長野県は標高 \SIrange{700}{1300}{\metre} の高地である。}
        計算では地上高度を 0 としているので、実際の飛行時間はわずかに短くなる。
        また 1・2 月に打ち上げ場所に雪がある日が 12--37 \% ある。
  \item \textbf{風は打ち上げ時刻・地点で固定している。}
        本文の第 \ref{sec:fidelity} 章のとおり、この近似による誤差は冬で \SIrange{5}{6}{\kilo\metre} で、
        結論は変わらない。ただし当日の予報誤差は含まれていない。
\end{itemize}

\subsection{まとめ}

\begin{itemize}
  \item 1・2 月は西日本のどこから上げても、真東へ \SIrange{240}{270}{\kilo\metre}（南西諸島は \SIrange{160}{230}{\kilo\metre}）流される。
        前半と後半の違いは小さい。
  \item 陸の平野（DID の外、雪なし）だけを目標にすると、標準の飛行条件では 1・2 月の平均でどこでも 25 \% 未満である。
        上昇を速くして飛行距離を 7 割に縮めても、長野県の佐久〜八ヶ岳付近から北関東の平野へ落とす形で約 4 割が上限になる。
  \item 瀬戸内海への着水を認めると、九州北西部と壱岐が有力になる。
        壱岐は標準の飛行条件のままで 47 \%、佐賀県北部は上昇を速くすると 62 \% が目標に入る。
        どちらも 1・2 月に打ち上げ場所の雪がほとんど無い。
\end{itemize}

地形は NOAA NCEI の DEM、市街地は国土数値情報（人口集中地区データ、国土交通省）、
積雪は ERA5（Copernicus Climate Change Service の情報を加工したものを含む）による。


\printbibliography[title={参考文献}]

\end{document}
